\documentclass[11pt]{article}
\usepackage[T1]{fontenc}
\usepackage{lmodern,amsmath,amssymb,amsthm,mathtools}
\usepackage[margin=28mm]{geometry}
\usepackage[hidelinks]{hyperref}
\usepackage{microtype}
\setlength{\emergencystretch}{3em}
\newtheorem{theorem}{Theorem}[section]
\newtheorem{lemma}[theorem]{Lemma}
\newtheorem{proposition}[theorem]{Proposition}
\newtheorem{remark}[theorem]{Remark}
\newcommand{\R}{\mathbb R}
\newcommand{\C}{\mathbb C}
\newcommand{\T}{\mathbb T}
\newcommand{\rank}{\operatorname{rank}}
\newcommand{\E}{\mathcal E}
\newcommand{\Tr}{\operatorname{Tr}}
\newcommand{\nct}{\mathsf n}
\newcommand{\ip}[2]{\langle #1,#2\rangle}
\title{Coarse information and finite rank obstructions\newline for flat lattice Hodge forms}
\author{Lluis Eriksson}
\date{7 November 2026 --- integrated revision candidate}
\usepackage{booktabs,array}
\newtheorem{matrixTheorem}[theorem]{Theorem}
\newtheorem{matrixProposition}[theorem]{Proposition}
\newtheorem{matrixExample}[theorem]{Example}
\providecommand{\tr}{\operatorname{tr}}
\providecommand{\im}{\operatorname{im}}
\providecommand{\id}{\mathrm I}
\providecommand{\Nlow}{N_{<c}}
\providecommand{\diag}{\operatorname{diag}}
\hypersetup{pdftitle={Coarse information and finite rank obstructions for flat lattice Hodge forms},pdfauthor={Lluis Eriksson}}
\begin{document}
\maketitle
\begin{abstract}
On a finite periodic lattice with unit spacing, a flat Hodge form cannot
have a positive coercivity constant uniform under refinement when supplemented
only by linear observations in a fixed finite-dimensional real space and
linear fine-space repairs of uniformly bounded rank. We construct a real
physical Fourier band, prove a necessary information budget and a separate
bounded-rank sector obstruction, and explain why a classical singular-Schur
condition prevents reversing that budget. Constructed derivative repairs
and differentiable observations with locally upper-quadratic penalties
admit tangent obstructions, including an exact regular nonlinear-sector
argument. For nonnegative weights, a subquadratic directional spatial mean
forces output-rank divergence. For signed weights, the spatial mean of the
pointwise positive part gives a sufficient condition; ordinary signed means
do not, as an actual fixed-output counterexample with every mean zero shows.
An attributed classical square cancellation delimits the sufficient criterion;
a finite-sum extension covering every eligible pair is displayed analytically.
A full volume coexact packet also gives a quantitative rank budget under
spatial means of the pointwise positive total coefficient on its two occupied
field components. A derived integer cutoff yields positive output rank density
on every sufficiently large lattice, while exact trace averaging improves a
sufficient low frequency constant. Formal endpoints are separately pinned, and additional analytic arguments
are identified. These are finite flat statements, with no interacting or
continuum Yang--Mills mass-gap claim or new general spectral-principle claim.
The integrated continuation also proves a critical inverse-moment sufficient
condition for global spectral observations, together with limits on spatial
localization. Rectangular means have a general upper bound with exponent
$d/4$ and the target exponent $d/2$ for an explicitly restricted scaled-cube
class. Two constructive counterexamples exclude a single-global-threshold
algorithm and a universal exact cube decomposition, while preserving successful
rectangular partitions for those same profiles. The general critical
rectangular question remains open. Self-contained local and supplied-cube proofs
are included, along with an actual physical fixed-map obstruction and its
separately scoped exact finite penalty threshold.
The present integrated revision adds positive matrix-gradient estimates,
finite-rank signed spectral repair and an exact critical-penalty obstruction
for prescribed rectangular means. Its leading penalty coefficient persists
on actual flat coexact fields; the complete ambient scalar optimum is
proved and accompanied by paired exact rational replay. These supporting
results do not settle the general critical rectangular rank-order question.
\end{abstract}

\section{Purpose and relation to the existing record}
This text consolidates a continuation of the exact-kernel obstruction in
\cite{eriksson}. That paper exhibits fields in the kernel of a particular
line-integral block map. Changing the map can remove a particular witness.
Here the obstruction depends instead on the amount of coarse information:
a whole low-energy subspace is too large to be seen by a fixed-dimensional
measurement together with a bounded-rank repair. The map need not factor
through the original block variable, be local, or obey a norm bound; it is
real-linear in the unrestricted-penalty results. The penalty applied to its output may be nonlinear. Later sections treat differentiable observations with a locally upper-quadratic penalty and exact nonlinear regular sectors.
The restricted-sector result is proved separately and does not follow merely
from failure of an estimate on the full space.

Fourier diagonalization, rank-nullity and finite-dimensional Schur complements
are classical. For the singular Schur complement and pseudoinverse background,
see Albert and Carlson--Haynsworth--Markham \cite{albert,carlson}. Finite-difference Hodge theory has a long history; Dodziuk's
work \cite{dodziuk} provides historical context, rather than this paper's
specific lattice normalization. We make no claim to a new general spectral
principle or an exhaustive priority search. The contribution is a concrete,
reproducibly checked constraint on proposed repair routes in this programme.
This consolidation is intended for the existing paper dossier, rather than
an additional deposit for each lemma.

Gauge-fixing and block averaging in renormalization-group constructions
require their own domain and scaling analysis. Dimock \cite{dimock} develops
covariant axial gauge for an abelian model and discusses its relation to
Balaban's programme. His introduction distinguishes a finite unit lattice
from a torus with small spacing, equivalently a large unit-spacing volume.
Our unrestricted, unit-spacing flat form is not identified with every
constrained fluctuation form in those constructions. This negative diagnostic
does not refute their gauge-fixed estimates or establish an interacting bound.

\section{Spaces and the actual flat form}
Fix integers $d\geq1$, $n\geq1$ and $N_c\geq2$. Set
$X_n=(\mathbb Z/n\mathbb Z)^d$ and $m=N_c^2-1$. The real one-cochain space is
\[
 E_n=\{A:X_n\times\{1,\ldots,d\}\longrightarrow\R^m\},\qquad
 \|A\|^2=\sum_{x\in X_n}\sum_{i=1}^d\|A_i(x)\|^2.
\]
The internal coordinates are those of the repository's Euclidean Lie-coordinate
model. No identification with a matrix Lie algebra is required here.
For positive unordered plaquettes $i<k$, define
\begin{align*}
 (\mathrm dA)_{ik}(x)&=A_k(x+e_i)-A_k(x)-A_i(x+e_k)+A_i(x),\\
 (\delta A)(x)&=\sum_i\bigl(A_i(x)-A_i(x-e_i)\bigr).
\end{align*}
The flat operator $K_{0,n}$ is the actual repository Hodge operator, with
\begin{equation}\label{eq:hodge}
 \ip{A}{K_{0,n}A}=\|\mathrm dA\|^2+\|\delta A\|^2.
\end{equation}
The plaquette norm sums each pair once. All differences have unit spacing;
there is no hidden $n^2$ or block-side squared factor.

Choose a direction $j$ and two distinct internal unit coordinates $u,v$.
They exist since $m\geq3$. For $f:X_n\to\C$, put
\[
 (E_jf)_i(x)=\begin{cases}(\Re f(x))u+(\Im f(x))v,&i=j,\\0,&i\ne j.\end{cases}
\]
This is a real-linear injection and $\|E_jf\|^2=\sum_x|f(x)|^2$.
Writing $D_i f(x)=f(x+e_i)-f(x)$, the actual curl and divergence give
\begin{equation}\label{eq:stencil}
 \ip{E_jf}{K_{0,n}E_jf}
 =\sum_{i\ne j}\sum_x|D_i f(x)|^2+
   \sum_x|f(x)-f(x-e_j)|^2
 =\sum_i\sum_x|D_i f(x)|^2.
\end{equation}
The last equality is finite translation invariance. The exclusive two
plaquette incidences and the divergence term account for every direction
exactly once, with no extra factor two.

\section{The original block kernel and its repair budget}\label{sec:old}
We retain the principal results of the historical paper \cite{eriksson}
before explaining the new information-budget obstruction. The original
Appendix~A (an even-block square-wave witness with quotient $8/\ell$) remains
in the preserved v1; it is not reproduced here. Use positive integers
$\ell,N'$, set $n=\ell N'$, and for a coarse site $y$ put
$B_y=\{\ell y+r:0\leq r_i<\ell\}$. The original line-integral map is
\[
 (Q_\ell A)_i(y)=\ell^{-d}\sum_{x\in B_y}\sum_{k=0}^{\ell-1}A_i(x+ke_i).
\]
For a constant field $A_i(x)=v_i$ and $S=\sum_i\|v_i\|^2>0$,
\[
 \|A\|^2=\ell^d(N')^dS,\qquad
 \|\ell^\sigma Q_\ell A\|^2=\ell^{2\sigma+2}(N')^dS.
\]
A uniform Poincare estimate must therefore have
$\sigma\geq(d-2)/2$ if it is to survive this sector. Equality alone makes
the two powers identical. Larger exponents pass this necessary test too;
neither equality nor the inequality is a sufficient full-space condition.
This follows by cancelling $(N')^dS$ to obtain
$\ell^{d-2-2\sigma}\leq C_P$. The real-exponent comparison is analytic here.

For $d\geq2$ choose distinct directions $p,q$ and the real isometry
$\iota(z)=(\Re z)u+(\Im z)v$ from Section~2. For $\ell\geq2$ define
\[
 A^{(\ell)}_i(x)=\mathbf1_{i=p}\,
 \iota\bigl(e^{2\pi\mathrm i x_q/\ell}\bigr),\qquad
 \lambda_\ell=|e^{2\pi\mathrm i/\ell}-1|^2
             =4\sin^2(\pi/\ell).
\]
The character is well defined on $X_{\ell N'}$. The sum over the $q$-offset
in every block is a nontrivial root-of-unity sum and vanishes. Shifts in the
line direction $p$ do not change the profile. Hence
\begin{equation}\label{eq:old-witness}
 Q_\ell A^{(\ell)}=0,\quad \delta A^{(\ell)}=0,\quad
 \|A^{(\ell)}\|^2=(\ell N')^d,\quad
 \ip{A^{(\ell)}}{K_{0,n}A^{(\ell)}}=(\ell N')^d\lambda_\ell.
\end{equation}
Only the plaquette pair containing $p,q$ contributes to the curl energy;
its difference has constant squared norm $\lambda_\ell$ at every site.
This also explains the exact energy in~\eqref{eq:old-witness}, including
periodic boundary jumps.

Consequently any scale-wise estimate
\[
 \|A\|^2\leq C_P(\ell)\bigl(\ip{A}{K_{0,n}A}+\Psi_\ell(Q_\ell A)\bigr),
 \qquad \Psi_\ell(0)=0,
\]
forces $C_P(\ell)\geq\lambda_\ell^{-1}\geq\ell^2/(4\pi^2)$.
No fixed $C_P>0$ before all block sides can work. This includes
$\Psi_\ell(y)=\|T_\ell y\|^2$ for arbitrary linear postconditioners, and
all critical or faster scalar rescalings. There is no regularity or sign
assumption on the scalar function $\Psi_\ell$.

The original spectral strengthening is retained analytically. For
$1\leq r\leq R\leq\lfloor(\ell-1)/2\rfloor$, replace the profile by
$\iota(c e^{2\pi\mathrm i r x_q/\ell})$, $c\in\C$. These are mutually
orthogonal real two-dimensional spaces in $\ker Q_\ell$. Derivative
orthogonality must be checked too: every forward difference multiplies a
character by its frequency multiplier, so the same root-of-unity Gram sum
annihilates cross terms between distinct frequencies. Thus for
$H_\ell=K_{0,n}+(T_\ell Q_\ell)^*(T_\ell Q_\ell)$ the sum of these spaces
has dimension $2R$ and maximum Rayleigh quotient at most
$4\sin^2(\pi R/\ell)\leq4\pi^2R^2/\ell^2$. Finite-dimensional min--max
gives $\mu_{2R}(H_\ell)\leq4\pi^2R^2/\ell^2$. Taking
$R=\lfloor\ell^\alpha\rfloor$, $0<\alpha<1$, for sufficiently large
$\ell$, yields a growing cluster below $4\pi^2\ell^{-2(1-\alpha)}$.
Eigenvalue language is restricted to this linear self-adjoint operator,
not applied to a general nonlinear penalty.

Finally, for any linear changed measurement $\widetilde Q_\ell$, linear
postconditioner $T_\ell$ and positive self-adjoint fine term $V_\ell$, put
\[
 \varepsilon_\ell=\frac{\|T_\ell\widetilde Q_\ell A^{(\ell)}\|}
                         {\|A^{(\ell)}\|},\qquad
 \nu_\ell=\frac{\ip{A^{(\ell)}}{V_\ell A^{(\ell)}}}
                    {\|A^{(\ell)}\|^2}.
\]
Set $\widetilde H_\ell=K_{0,n}+
(T_\ell\widetilde Q_\ell)^*(T_\ell\widetilde Q_\ell)+V_\ell$.
Expansion gives the exact repaired quotient
\[
 \frac{\ip{A^{(\ell)}}{\widetilde H_\ell A^{(\ell)}}}
      {\|A^{(\ell)}\|^2}
 =\lambda_\ell+\varepsilon_\ell^2+\nu_\ell.
\]
A full-space constant $c>0$ therefore requires
$\varepsilon_\ell^2+\nu_\ell\geq c-\lambda_\ell$ at every scale, hence
$\liminf_{\ell\to\infty}(\varepsilon_\ell^2+\nu_\ell)\geq c$.
On the single witness line, uniform coercivity is equivalent to
$\inf_{\ell\geq2}(\lambda_\ell+\varepsilon_\ell^2+\nu_\ell)>0$;
this is not full-space sufficiency. In particular if
$\widetilde Q_\ell=Q_\ell+R_\ell$, $T_\ell=\ell^{(d-2)/2}I$ and
$\eta_\ell=\|R_\ell A^{(\ell)}\|/\|A^{(\ell)}\|$, then
$\varepsilon_\ell^2=\ell^{d-2}\eta_\ell^2$. In dimension four and with
$V_\ell=0$, leakage $o(\ell^{-1})$ cannot even repair this line.
The liminf, real-exponent, spectral multiplicity and witness-line equivalence
are analytic statements. The older exact kernel, norm-form energy, finite
budget and dimension-four factorized no-go have their own preserved formal
source and reproductions, separately anchored below.
\section{A constructed low band}
For a nonnegative integer $R$ with $2R<n$, let $\Lambda_R=\{-R,\ldots,R\}^d$ and
$B=|\Lambda_R|=(2R+1)^d$. Define
\[
 \chi_k(x)=\exp(2\pi\mathrm i\,k\cdot x/n),\qquad
 f_a=\sum_{k\in\Lambda_R}a_k\chi_k,\qquad
 S_R=\{E_j f_a:a\in\C^{\Lambda_R}\}\subset E_n.
\]
\begin{lemma}\label{lem:band}
The real subspace $S_R$ has dimension $2B$, and for every $A\in S_R$,
\begin{equation}\label{eq:band}
 \ip{A}{K_{0,n}A}\leq b_R(n)\|A\|^2,
 \qquad b_R(n)=\frac{4\pi^2dR^2}{n^2}.
\end{equation}
\end{lemma}
\begin{proof}
For $k,\ell\in\Lambda_R$, the product root-of-unity sum is
$\sum_x\overline{\chi_k(x)}\chi_\ell(x)=n^d\mathbf1_{k=\ell}$.
The no-alias condition ensures that coordinate differences, of absolute value
at most $2R<n$, cannot be nonzero multiples of $n$. Thus synthesis is injective
and its complex image has dimension $B$. As a real space that image has
dimension $2B$; applying the real-linear injection $E_j$ preserves dimension.
The factor two represents independent real and imaginary coefficients stored
in two real internal coordinates, not two copies of a real scalar band.

For completeness, energy orthogonality is needed as well as mass orthogonality.
For each $i$,
\[
 D_i f_a=\sum_k a_k\bigl(e^{2\pi\mathrm i k_i/n}-1\bigr)\chi_k.
\]
Applying the same Gram sum to these derivative amplitudes gives
\begin{align*}
 \|E_jf_a\|^2&=n^d\sum_k|a_k|^2,\\
 \ip{E_jf_a}{K_{0,n}E_jf_a}
 &=n^d\sum_k |a_k|^2\sum_i|e^{2\pi\mathrm i k_i/n}-1|^2.
\end{align*}
Since $|e^{\mathrm it}-1|\leq|t|$ and $|k_i|\leq R$, each sum of multipliers
is at most $4\pi^2dR^2/n^2$. This proves the bound for every member of the
constructed image, not only its individual generators.
\end{proof}

\section{The finite information budget}
Let $F$ be a finite-dimensional real vector space, $C:E_n\to F$ a linear
measurement and $V:E_n\to E_n$ a linear operator. No symmetry or positivity
of $V$ is needed. Let $\Phi:F\to\R$ be any function with $\Phi(0)=0$.
\begin{theorem}\label{thm:budget}
If $c>b_R(n)$ and
\begin{equation}\label{eq:estimate}
 c\|A\|^2\leq\ip{A}{K_{0,n}A}+\Phi(CA)+\ip{A}{VA}
 \quad\text{for all }A\in E_n,
\end{equation}
then $2(2R+1)^d\leq\dim F+\rank V$.
\end{theorem}
\begin{proof}
If the budget is smaller than $\dim S_R$, restrict the map $(C,V)$ to $S_R$.
Its image has dimension at most $\dim F+\rank V$. Rank-nullity yields a
nonzero $A\in S_R$ with $CA=VA=0$. Equation~\eqref{eq:estimate} and
Lemma~\ref{lem:band} imply $c\|A\|^2\leq b_R(n)\|A\|^2$, a contradiction.
\end{proof}
The codomain dimension is a conservative measurement budget. We do not label
this formal endpoint as a sharper theorem with $\rank C$ in its place.
No regularity, sign, polynomial form or growth restriction on $\Phi$ enters.

\begin{proposition}\label{prop:sector}
If $P:E_n\to E_n$ is linear, $c>b_R(n)$ and
$\dim F+\rank P<2(2R+1)^d$, the inequality
\[
 c\|A\|^2\leq\ip{A}{K_{0,n}A}+\Phi(CA)
 \qquad(A\in\ker P)
\]
is false.
\end{proposition}
\begin{proof}
Apply the same rank-nullity argument to $(C,P)|_{S_R}$. Its nonzero kernel
vector belongs to the stipulated sector as well as the unobserved band.
\end{proof}
For example, $d=4,n=16,R=1$ gives dimension $162$ and
$b_R(n)=\pi^2/16<1$. Twelve measurement coordinates and a repair of rank at
most $149$ cannot yield~\eqref{eq:estimate} with $c=1$. This illustrates a
necessary budget, not the programme's chosen target constant.

\section{Uniform refinement and sector exclusion}
Fix a positive integer $N'$ and put $n_L=(L+1)N'$ for nonnegative integers $L$.
The finite-dimensional space $F$ is fixed. Allow arbitrary families
of linear maps $C_L:E_{n_L}\to F$, arbitrary functions $\Phi_L:F\to\R$
with $\Phi_L(0)=0$, and linear maps $V_L:E_{n_L}\to E_{n_L}$ with
$\rank V_L\leq r$ for one fixed nonnegative integer $r$.
\begin{theorem}\label{thm:uniform}
There is no $c>0$ for which~\eqref{eq:estimate} holds for every $L$ and every
$A\in E_{n_L}$ with these families. Likewise, for linear $P_L$ of rank at
most $r$, there is no positive uniform $c$ for the sector estimates of
Proposition~\ref{prop:sector} on all $\ker P_L$.
\end{theorem}
\begin{proof}
Set $M=\dim F$ and $R=M+r+1$. Since $d\geq1$,
$2(2R+1)^d>M+r$. Choose $L$ so large that $2R<n_L$ and $b_R(n_L)<c$.
Here is a fully finite choice argument, without an assumed limit theorem:
with $a=4\pi^2dR^2$, take an integer
$L>\max\{2R,a/c,1\}$. Since $N'\geq1$, $n_L\geq L+1>L>1$, and
$a<cL\leq cn_L^2$. Theorem~\ref{thm:budget} now contradicts the rank bound.
Use Proposition~\ref{prop:sector} for the separate sector conclusion.
\end{proof}
The constants $c,r$ precede the quantifier over $L$. Arbitrary rescalings of
$C_L$ preserve its fixed codomain and do not escape the theorem. Removing
harmonic or constant directions is covered only when a concrete linear
constraint describing the proposed sector has uniformly bounded rank.
A quotient or its norm is not silently identified with such a sector.
Linearity is essential: if nonlinear observations were allowed, the single
real observation $C_L(A)=\|A\|^2$, with $\Phi_L(y)=y$ and $V_L=0$, would give
the uniform lower bound $c=1$ since $K_{0,n_L}\succeq0$. This counterexample
explains why finite-dimensional output alone is not the hypothesis.

\section{Constructed derivative repairs and tangent observations}\label{sec:tangent}
The bounded-rank obstruction does not exhaust the available repairs. For the
actual componentwise forward gradient $\mathcal G A=(D_k A_i)_{x,i,k}$, finite
periodic summation and cancellation of the mixed differences give
\[
 \|\mathcal G A\|^2=\ip{A}{K_{0,n}A}.
\]
This is an identity on the whole cochain space, with unordered plaquettes
counted once; it is not an assumed comparison on a supplied band. Let $W$ be
the diagonal multiplication by real weights $w_{x,i,k}$, and construct
$V_w=\mathcal G^*W\mathcal G$. If $w_{x,i,k}\le M$ and $M\ge0$, then
\begin{equation}\label{eq:weighted-actual}
 \ip{A}{V_w A}=\sum_{x,i,k}w_{x,i,k}\|(D_k A_i)(x)\|^2
 \le M\ip{A}{K_{0,n}A}.
\end{equation}
A lower bound on the weights is unnecessary for this upper estimate.
Consequently the constructed Fourier band satisfies
\[
 q_w(A):=\ip{A}{(K_{0,n}+V_w)A}
 \le (1+M)b_R(n)\|A\|^2 \qquad(A\in S_R).
\]
These operator identities and inequalities are formally proved in the
separately pinned source reported below.

\begin{theorem}[Finite-dimensional tangent obstruction]\label{thm:tangent}
Let $E$ be a finite-dimensional real inner-product space, $F$ a finite-dimensional
real normed space, and $S\subseteq E$ a linear subspace with $\dim S>\dim F$.
Let $q:E\to\R$ satisfy $q(x)\le b\|x\|^2$ on $S$. Suppose $C:E\to F$ is
Fr\'echet differentiable at zero, $C(0)=0$, and for some $\delta>0$ and finite
$K\in\R$,
\begin{equation}\label{eq:phi-local}
 \Phi(y)\le K\|y\|^2\qquad(\|y\|<\delta).
\end{equation}
If $c>b$, the full-space estimate
$c\|x\|^2\le q(x)+\Phi(C(x))$ is impossible.
\end{theorem}
\begin{proof}
Put $L=DC(0)$. Rank-nullity applied to the actual restriction $L|_S$ produces
$v\in S\cap\ker L$, $v\ne0$. Differentiability gives
$\|C(tv)\|/\|tv\|\to0$ as nonzero real $t\to0$. For small such $t$,
\eqref{eq:phi-local} applies and the proposed bound implies
\[
 c\le b+K\frac{\|C(tv)\|^2}{\|tv\|^2}.
\]
Taking the limit gives $c\le b$, a contradiction. The kernel vector and the
remainder limit are derived, not introduced as hypotheses of the theorem.
\end{proof}
The three focal Lean endpoints include this theorem, its remainder lemma and
the actual physical specialization with threshold $(1+M)b_R(n)$. The proof
does not require an explicit assumption $\Phi(0)=0$. A local-ball version is
also valid by the same shrinking-ray argument, while the original formal endpoint above is stated on the full space. The later local kernel and average-weight endpoint in Section~\ref{sec:average} has its own separately reproduced proof.

This result trades hypotheses. The global linearity of $C$ is replaced by
differentiability at zero, while an arbitrary zero-vanishing penalty is
replaced by genuine quadratic upper control. Merely asking $C$ and $\Phi$
to be $C^1$ with zero derivatives at zero is insufficient: for $1<\alpha<2$,
\[
 C(x)=\|x\|^\alpha,\qquad \Phi(y)=|y|^{2/\alpha}
 \quad\Longrightarrow\quad \Phi(C(x))=\|x\|^2.
\]
Both derivatives at zero vanish. Off zero,
$\|DC(x)\|=\alpha\|x\|^{\alpha-1}$ and
$|\Phi'(y)|=(2/\alpha)|y|^{2/\alpha-1}$ tend to zero, so both maps are indeed
$C^1$ there and at zero. But
$\Phi(y)/|y|^2=|y|^{2/\alpha-2}$ is unbounded above near zero. On the actual
lattice, setting $w=0$ gives a uniform lower bound with $c=1$ because
$K_{0,n}\succeq0$. This is an exact counterexample, not a numerical test.

\section{Exact nonlinear regular sectors}\label{sec:nonlinear-sector}
The next theorem is an analytic consequence of rank-nullity and the classical
implicit function theorem. The actual curve construction and local-energy obstruction now have separately pinned formal proofs, identified in the evidence section.
The pinned formal library supplies an implicit parametrization with derivative
the inclusion of the derivative kernel \cite{mathlib-implicit}; we construct
and use that parametrization explicitly below rather than assuming the curve
needed for the conclusion.

\begin{theorem}[Regular-sector tangent obstruction]\label{thm:regular-sector}
Let $E$ be a finite-dimensional real inner-product space and $F,G$ finite-dimensional
real normed spaces. Let $T:E\to E$ be linear and set $q(x)=\ip{x}{Tx}$.
Suppose a linear subspace $S$ satisfies $q(v)\le b\|v\|^2$ on $S$ and
$\dim S>\dim F+\dim G$. Let $C$ be differentiable at zero with $C(0)=0$ and
let $\Phi$ satisfy~\eqref{eq:phi-local}. Let $P:E\to G$ be $C^1$ on a
neighborhood of zero with $P(0)=0$ and $A=DP(0)$ surjective. For every
$c>b$ and every $r>0$, there is no estimate
\begin{equation}\label{eq:nonlinear-sector-bound}
 c\|x\|^2\le q(x)+\Phi(C(x))
 \qquad\bigl(P(x)=0,\ \|x\|<r\bigr).
\end{equation}
\end{theorem}
\begin{proof}
Put $L=DC(0)$. Apply rank-nullity to $(L,A)|_S:S\to F\times G$ to obtain
$v\in S$, $v\ne0$, with $Lv=Av=0$. Choose a linear complement $W$ of
$H=\ker A$. Surjectivity implies $A|_W:W\to G$ is bijective. Apply the
implicit function theorem to $(u,w)\mapsto P(u+w)$ at $(0,0)\in H\times W$.
It produces a $C^1$ map $\psi$ near zero in $H$ with
\[
 \psi(0)=0,\quad P(u+\psi(u))=0,\quad
 D\psi(0)=-(A|_W)^{-1}A|_H=0.
\]
Thus $\gamma(t)=tv+\psi(tv)$ lies in the exact sector for small real $t$,
$\gamma(0)=0$, and $\gamma(t)=tv+e(t)$ with $\|e(t)\|=o(|t|)$. In particular
$\|\gamma(t)\|/|t|\to\|v\|>0$, so the nonzero points of this curve eventually
lie in the stipulated ball.

The curve need not remain in $S$. This is why the quadratic form assumption
is used: direct expansion gives
\[
 q(\gamma(t))=t^2q(v)+t\ip{v}{Te(t)}
       +t\ip{e(t)}{Tv}+\ip{e(t)}{Te(t)}.
\]
Since $T$ is bounded in finite dimension, division by $t^2$ gives
$q(\gamma(t))/t^2\to q(v)$. Differentiability of $C$ and $Lv=0$ give
$C(\gamma(t))=Le(t)+o(|t|)=o(|t|)$. Apply the assumed bound along this exact
curve and use~\eqref{eq:phi-local}. Division by $t^2$ and the limit yield
$c\|v\|^2\le q(v)\le b\|v\|^2$, contrary to $c>b$.
\end{proof}
Unlike the preceding full-space result, an arbitrary energy with an upper
bound only on $S$ would not justify this proof: the constructed exact-sector
curve generally leaves $S$. The displayed expansion supplies the missing
comparison for the actual quadratic energy.

A concrete curved example also shows why an exact vector in $S$ is not the
right object. Write $E=\R^4$ with coordinates $(u_1,u_2,y,z)$, put
$P(u,y,z)=z-\|u\|^2$, $C(u,y,z)=y$, $\Phi(y)=y^2$, $S=\{z=0\}$ and
$q(u,y,z)=z^2$. Then $DP(0)$ is surjective and $\dim S=3>1+1$.
But $S\cap\ker C\cap P^{-1}(0)=\{0\}$: the linear band contains no
nonzero exact-sector unseen vector. Nevertheless
$\gamma(t)=(t,0,0,t^2)$ lies in the exact sector and has repaired Rayleigh
quotient $t^2/(1+t^2)\to0$. The implicit-curve argument is genuinely needed.

Nor can quadratic $q$ in Theorem~\ref{thm:regular-sector} be replaced merely
by a $C^1$ energy stationary at zero with the same upper bound on $S$.
With the same $E,C,\Phi,S$, instead put
\[
 P(u,y,z)=z-\|u\|^{4/3},\qquad e(u,y,z)=|z|^{3/2}.
\]
Both maps are $C^1$, $DP(0)(u,y,z)=z$ is surjective, $De(0)=0$, and $e=0$
on $S$, so $b=0$. On the exact sector,
\[
 e(x)+\Phi(C(x))=\|u\|^2+y^2,\qquad
 \|x\|^2=\|u\|^2+y^2+\|u\|^{8/3}.
\]
For $\|x\|<1$, $\|u\|^{8/3}\le\|u\|^2$, so the local sector bound with
$c=1/2>b$ holds. Differentiability of the norm powers at zero follows from
their exponents being greater than one; their gradients off zero tend to
zero. This exact counterexample refutes a proposed $C^1$-energy weakening,
while the actual quadratic Hodge/repair energy retains the required expansion.

The regularity condition cannot simply be dropped. Take $E=\R^3$,
$F=G=\R$, $C=0$, $\Phi=0$, $T=0$, $S=E$, and $P(x)=\|x\|^2$.
Then $P$ is smooth, $DP(0)=0$, and the strict dimension budget holds, but
$P^{-1}(0)=\{0\}$ and the estimate with $c=1$ is vacuous on the sector.
This refutes the proposal to replace submersion by bounded derivative rank.
It does not refute Theorem~\ref{thm:regular-sector}, because $DP(0)$ is not
surjective. Critical constraint maps require a separate analysis.

\begin{proposition}[Uniform physical consequence]\label{prop:regular-uniform}
Let $n_L\to\infty$ with fixed $d,N_c$, and let $F,G$ have fixed finite dimensions.
For each $L$, allow its own $C_L,P_L,\Phi_L$ satisfying the hypotheses of
Theorem~\ref{thm:regular-sector}. Let $q_L$ be the actual flat Hodge form plus
the constructed weighted derivative repair with weights $w_L\le M_L$,
$M_L\ge0$, and $M_L/n_L^2\to0$. Even on the exact sectors $P_L=0$ and
arbitrarily small balls of positive radius $r_L$, no positive uniform
coercivity constant exists.
\end{proposition}
\begin{proof}
Choose a fixed integer $R$ with $2(2R+1)^d>\dim F+\dim G$. For all large $L$,
$2R<n_L$ and the constructed band obeys
$q_L\le (1+M_L)b_R(n_L)\|\cdot\|^2$. Its coefficient tends to zero.
For any proposed $c>0$, choose one such $L$ with this coefficient below $c$
and apply Theorem~\ref{thm:regular-sector} at that fixed lattice.
The curve is shrunk only after choosing $L$: no uniform implicit-function
radius, inverse-derivative norm, penalty constant or $r_L$ is required.
\end{proof}
These are finite unit-spacing obstructions for regular nonlinear sectors,
not identifications of a physical gauge-fixed domain or a quotient norm.


\section{Directional mean weights and local nonlinear observations}\label{sec:average}
Pointwise bounds on a derivative repair are stronger than the following
actual spatial-mean condition. Fix $d\geq2$, $N_c\geq2$, and distinct
directions $i,j$. For nonnegative componentwise weights $w_n(k,(x,l))$, put
\[
 Q_n(A)=\ip{A}{K_{0,n}A}+\sum_{x,k,l}w_n(k,(x,l))
                  \|A_l(x+e_k)-A_l(x)\|^2,
 \qquad b_{ij}(n)=n^{-d}\sum_x w_n(i,(x,j)).
\]
The repair is the actual $G^*W_nG$ of the componentwise forward gradient.
In particular $Q_n\geq0$; signed weights are not included in this section.
Let $D_n=\ker\delta$ be the actual flat divergence-free space.

\begin{theorem}[Local nonlinear average-weight obstruction]\label{thm:average}
For each $n\geq1$, let $F_n$ be a finite-dimensional real normed space,
$C_n:E_n\to F_n$ be Fr\'echet differentiable at zero with $C_n(0)=0$,
and let $\Phi_n:F_n\to\R$ satisfy
\[
 \Phi_n(y)\leq K_n\|y\|^2\quad(\|y\|<\delta_n),\qquad \delta_n>0.
\]
Suppose one fixed $c>0$ and positive radii $r_n$ satisfy
\begin{equation}\label{eq:local-average}
 c\|A\|^2\leq Q_n(A)+\Phi_n(C_n(A))
       \quad(A\in D_n,\ \|A\|<r_n).
\end{equation}
For every integer $0\leq R<n$,
\begin{equation}\label{eq:mean-finite}
 c(2R-\dim F_n)\leq 2(1+b_{ij}(n))
       \sum_{s=1}^{R}|e^{2\pi\mathrm{i}s/n}-1|^2
 \leq \frac{8\pi^2(1+b_{ij}(n))R^3}{n^2}.
\end{equation}
Consequently, if $(1+b_{ij}(n))/n^2\to0$, then $\dim F_n\to\infty$.
No uniform bound on $K_n,\delta_n^{-1}$ or $r_n^{-1}$ is required.
\end{theorem}
\begin{proof}
Write $L_n=DC_n(0)$. If $v\in D_n\cap\ker L_n$, differentiability gives
$C_n(tv)=o(|t|)$. For all sufficiently small nonzero $t$, both the local
coercivity and penalty bound apply. Divide by $t^2$ and let $t\to0$ to get
\begin{equation}\label{eq:kernel-average}
 c\|v\|^2\leq Q_n(v).
\end{equation}
This argument is performed separately at each $n$; it does not impose
uniform neighborhoods or derivative remainder rates.

Use the two internal coordinates $u,v$ of Section 2 to regard a complex
number as a real internal vector. For $1\leq s\leq R$ and $a\in\{1,\mathrm{i}\}$,
define $\psi_{s,a}$ by taking only component $j$ nonzero and setting it to
$n^{-d/2}$ times the internal vector corresponding to
$a e^{2\pi\mathrm{i}s x_i/n}$. The finite-character orthogonality identity
and the real orthogonality of the two quadratures give an orthonormal
family of $2R$ fields. This also covers even $n$; the two quadratures use
two internal coordinates, not two allegedly independent scalar sine and
cosine functions. Since $i\ne j$, these fields are independent of $x_j$,
so $\delta\psi_{s,a}=0$. Their span $S_{R,n}$ is thus an actual
$2R$-dimensional subspace of $D_n$.

Every nonzero forward difference of $\psi_{s,a}$ is in direction $i$ and
has the same squared norm at every site, namely
$n^{-d}|e^{2\pi\mathrm{i}s/n}-1|^2$. The exact Hodge and repair energies are
therefore
\[
 Q_n(\psi_{s,a})=(1+b_{ij}(n))|e^{2\pi\mathrm{i}s/n}-1|^2.
\]
In particular the trace of the compression of $Q_n$ to $S_{R,n}$ is the
middle expression in \eqref{eq:mean-finite}; no invariance of $S_{R,n}$
under the weighted operator is asserted or needed.

For completeness, the positive trace/rank argument is as follows. Set
$U=S_{R,n}\cap\ker L_n$. Rank-nullity gives
$\dim U\geq2R-\dim F_n$. Extend an orthonormal basis of $U$ to one of
$S_{R,n}$. Equation \eqref{eq:kernel-average} bounds the terms on $U$
below by $c$; positivity of $Q_n$ makes every remaining diagonal term
nonnegative. Hence $c(2R-\dim F_n)\leq\operatorname{tr}_{S_{R,n}}Q_n$.
Finally $|e^{2\pi\mathrm{i}s/n}-1|\leq2\pi s/n\leq2\pi R/n$, which proves
the last inequality in \eqref{eq:mean-finite}.

Fix $R$. The right side tends to zero under the stated mean condition.
Eventually it is less than $c$. If $\dim F_n<2R$, integrality would make
the left side at least $c$. Thus $\dim F_n\geq2R$ eventually. Arbitrarily
large fixed $R$ prove the claimed divergence.
\end{proof}

The condition concerns one directional spatial average, not the maximum
of all weights. For example, a weight of size $n^{d+1}$ at one site in
component $(i,j)$ and zero elsewhere has $b_{ij}(n)=n$, while its maximum
is not subquadratic. Theorem~\ref{thm:average} still excludes bounded
output dimension under its other hypotheses. This example is a direct
calculation, not an additional formal endpoint. The nonnegative-weight
condition is essential to the trace argument as used here: the complement
of the derivative kernel must contribute nonnegative energy. This is not
a universal extension of older statements allowing signed weights with a
pointwise absolute cap.

\subsection{A constructed quadratic-weight boundary example}
The exponent $2$ cannot be replaced by an unrestricted allowance of
quadratic-size mean weights. Let
$\mu_j(A)=n^{-d}\sum_x A_j(x)$, with output in the same fixed space
$F=(\R^{N_c^2-1})^d$ for every $n$. A finite path comparison on the actual
periodic lattice gives
\[
 \|A-\operatorname{lift}\mu(A)\|^2\leq d n^2\|GA\|^2,
 \qquad \|\operatorname{lift}\mu(A)\|^2=n^d\|\mu(A)\|^2.
\]
One way to obtain the first inequality is to write each difference of
site values as a path of at most $dn$ forward edges, apply finite
Cauchy--Schwarz, and average the pairs of sites; the displayed constant
is deliberately nonoptimal. The exact modular-torus reindexing and its
transport to the physical cochain norm are also proved in the companion
formal construction.
Using $\|x+y\|^2\leq2\|x\|^2+2\|y\|^2$, choose actual weights
$w_n(k,(x,l))=2dn^2$ and the quadratic penalty
$\Phi_n(y)=2n^d\|y\|^2$. Then
\[
 \|A\|^2\leq\ip{A}{G^*W_nGA}+\Phi_n(\mu(A))
                  \leq Q_n(A)+\Phi_n(\mu(A))
\]
on the entire space, hence also on $D_n$. The mean observation is an
actual continuous linear map with finite-dimensional output type
independent of $n$, not a postulated observation certificate. Here
$b_{ij}(n)=2dn^2$ and the subquadratic hypothesis fails. The example
allows $K_n=2n^d$ to vary, exactly as the obstruction does. It proves
neither an optimal boundary constant nor positive coercivity for an
interacting gauge operator. The path/mean mechanism is classical;
see the primary exposition in \cite{levin-peres}.

\subsection{Why the penalty assumption remains}
The $C^1$ counterexample after Theorem~\ref{thm:tangent} also applies to
the actual nonnegative $Q_n$ and to $D_n$. A second useful counterexample
retains a polynomial observation and a globally Lipschitz penalty:
\[
 C(A)=\|A\|^2\in\R,\qquad \Phi(y)=c|y|,
 \qquad DC(0)=0,\qquad \Phi(C(A))=c\|A\|^2.
\]
This example already gives the bound for the actual admissible choice $w=0$,
with $Q_n=K_{0,n}\succeq0$. The zero-form interpretation is only an abstract
illustration, not an instance of this nonnegative-weight lattice family.
No finite local
quadratic upper constant for $\Phi$ exists, since $c|y|/y^2=c/|y|$.
Likewise, $C(A)=\|A\|^{3/2}$ and $\Phi(y)=c|y|^{4/3}$ are both $C^1$
with zero derivatives at zero, yet their composition is exactly quadratic.
These are complete analytic counterexamples, not new Lean claims.

There is a distinct analytic regularity tradeoff, not a hypothesis removal
for merely differentiable observations. If, on each fixed derivative-kernel
direction, $C(tv)=t^2H_v+o(t^2)$, and $\Phi(0)=0$, $D\Phi(0)=0$, then
$\Phi(C(tv))=o(t^2)$. Indeed $\|C(tv)\|/t^2$ is bounded and
$|\Phi(y)|\leq\epsilon\|y\|$ near zero, including at zero; multiply the
two estimates and let $\epsilon\downarrow0$. A $C^2$ observation and a
nonnegative differentiable penalty vanishing at zero suffice for this
step, because a differentiable local minimum has zero derivative.
No actual observation's $C^2$ regularity is proved here, and no formal
composition of this alternate route is claimed.

The strengthened mean-weight result is an output-rank divergence theorem.
Its fixed-$R$ argument does not assert a positive density of observations
per lattice volume. Fixed $c>0$, nonnegative weights and the stated
penalty condition remain explicit. In particular no interacting or
continuum Yang--Mills hypothesis, ultraviolet activity estimate, or
physical spectral gap is supplied by this section.

\section{Signed weights and the limits of mean criteria}\label{sec:signed}
Write $R_w=G^*W G$ and $Q_w(A)=\ip{A}{K_{0,n}A}+\ip{A}{R_wA}$.
For signed weights this form need not be nonnegative. Its spatial means
alone no longer supply the positive trace comparison used above.
The following sufficient condition is nevertheless available without a
lower bound on any original weight. Define the \emph{pointwise} positive part
and its selected directional mean by
\[
 w^+(k,(x,l))=\max\{w(k,(x,l)),0\},\qquad
 b^+_{ij}(n)=n^{-d}\sum_x w^+(i,(x,j)).
\]
In general $b^+_{ij}$ is not the positive part of the signed mean.

\begin{theorem}[A sufficient signed-weight condition]\label{thm:signed-positive}
Keep the observation, penalty, local divergence-free domain and fixed
$c>0$ of Theorem~\ref{thm:average}, but use arbitrary signed actual weights
in the local coercivity inequality. For every $0\leq R<n$,
\[
 c(2R-\dim F_n)\leq
       \frac{8\pi^2(1+b^+_{ij}(n))R^3}{n^2}.
\]
If $(1+b^+_{ij}(n))/n^2\to0$, then $\dim F_n\to\infty$.
The original $C_n,L_n,\Phi_n,c,K_n,\delta_n,r_n$ are retained.
\end{theorem}
\begin{proof}
The exact weighted-gradient energy gives, for every actual cochain,
\[
 \ip{A}{R_w A}=\sum_{x,k,l}w(k,(x,l))\|D_k A_l(x)\|^2
       \leq\sum_{x,k,l}w^+(k,(x,l))\|D_k A_l(x)\|^2
       =\ip{A}{R_{w^+}A}.
\]
The inequality uses only nonnegativity of the squared differences.
Thus the original local coercivity inequality transfers upward to the
constructed nonnegative repair $R_{w^+}$, and Theorem~\ref{thm:average}
applies. Nonnegativity of $w^+$ and the comparison are proved from their
actual definitions, not supplied as fields of a comparison structure.
\end{proof}
This changes the explicit mean hypothesis for signed weights. It neither
proves the original weights nonnegative nor removes the local regularity,
quadratic penalty or fixed positive coercivity requirements.

\subsection{Zero signed means with fixed output}
There is no counterpart using only ordinary signed directional means,
even when \emph{every} such mean is zero. Fix a site $o\in X_n$, set
$V=n^d$, $a=2dn^2$, and construct actual weights
\[
 w_n(k,(x,l))=a\bigl(1-V\mathbf1_{x=o}\bigr).
\]
Since $|X_n|=V$, their ordinary spatial means are all exactly zero.
Let $\mu(A)$ be the directional mean above, and let
$J_n(A)=(D_k A_l(o))_{k,l}$ collect all actual forward differences at $o$.
Use the same finite output type for every side:
\[
 F=(\R^m)^d\times(\R^m)^{d\times d},\quad
 C_n(A)=(\mu(A),J_n(A)),\quad
 \|(y,z)\|_F=\max\{\|y\|_2,\|z\|_2\}.
\]
Its dimension is $(d+d^2)m$, independent of $n$. The companion code proves
the fixed finite output type; this displayed numerical dimension is an
analytic calculation, not an additional named oracle endpoint.

\begin{proposition}\label{prop:signed-zero}
The actual continuous linear observation $C_n$ and the nonnegative
quadratic penalty
\[
 \Phi_n(y,z)=2V\|y\|_2^2+aV\|z\|_2^2
\]
satisfy $C_n(0)=0$, $DC_n(0)=C_n$,
$\Phi_n(y,z)\leq V(2+a)\|(y,z)\|_F^2$, and
\[
 \|A\|^2\leq Q_{w_n}(A)+\Phi_n(C_n(A))\qquad(A\in E_n).
\]
Consequently all the original local observation and penalty hypotheses
can hold with fixed output and $c=1$, although all signed means vanish.
\end{proposition}
\begin{proof}
Linearity follows directly from the finite sums and differences. The
penalty bound follows from each projection norm being at most the
specified product max norm. The exact energy expansion is
\[
 \ip{A}{R_{w_n}A}=a\|GA\|^2-aV\|J_n(A)\|_2^2.
\]
Adding the penalty cancels the negative term and leaves
\[
 Q_{w_n}(A)+\Phi_n(C_n(A))
   =\ip{A}{K_{0,n}A}+a\|GA\|^2+2V\|\mu(A)\|_2^2
   \geq\|A\|^2.
\]
The last step is the actual mean/path estimate of the preceding section,
with $a=2dn^2$, together with $K_{0,n}\succeq0$.
\end{proof}
The coefficient $K_n=V(2+a)$ may grow with the lattice side, exactly as
allowed by the proposed ordinary-mean extension. This is not a transfer
of the desired coercivity into an assumed observation certificate.
The weights, observation, penalty and cancellation are constructed.
For $n>1$ their positive-part mean is $a(1-V^{-1})$, of quadratic rather
than subquadratic order, so they do not contradict
Theorem~\ref{thm:signed-positive}.

As an analytic diagnostic, take $n\geq3$, $i\ne j$, a unit internal vector
$u$, and $A_j(x)=u\mathbf1_{x_i=o_i}$, with other components zero. Then
$\delta A=0$, $\|GA\|^2=2n^{d-1}$ and $\|J_n(A)\|^2=1$, giving
\[
 Q_{w_n}(A)=n^{d-1}\bigl(2-a(n-2)\bigr)<0.
\]
This explains the failure of the original positive trace argument.
The diagnostic field calculation is analytic; the formal fixed-output
family theorem does not claim to check this additional field explicitly.

\subsection{An attributed square cancellation}
Fix distinct directions $i,k$, a component $j\ne i$ (for instance $j=k$),
and $o\in X_n$. For an amplitude $A_0\geq0$, define the repair
weights additively by
\begin{align*}
 w(p)&=A_0\mathbf1_{p=(i,(o,j))}
       -3A_0\mathbf1_{p=(k,(o,j))}\\
     &\quad-3A_0\mathbf1_{p=(i,(o+e_k,j))}
       -3A_0\mathbf1_{p=(k,(o+e_i,j))}.
\end{align*}
Writing $D_s A_j(x)=A_j(x+e_s)-A_j(x)$, commuting shifts give
\[
 D_i A_j(o)=D_k A_j(o)+D_i A_j(o+e_k)-D_k A_j(o+e_i).
\]
The squared norm of this sum is at most three times the sum of the
three squared norms. The actual finite weighted-energy sum therefore
proves $\ip{A}{R_wA}\leq0$ for every cochain.
For $n\geq2$ the positive edge is separate from all negative contributions,
so $w^+=A_0\mathbf1_{p=(i,(o,j))}$ and $b^+_{ij}=A_0/n^d$.
Choosing $A_0=n^d n^2$ gives $b^+_{ij}=n^2$.

If the original local coercivity hypotheses hold for this actual repair,
they imply the same hypotheses with zero repair, since $Q_w(A)\leq
\ip{A}{K_{0,n}A}$. Applying Theorem~\ref{thm:average} with zero weights
again forces $\dim F_n\to\infty$. The selected normalized positive-part mean is
$1+n^{-2}$ and does not tend to zero. This last limit interpretation is
an elementary analytic consequence of the formal exact mean identity,
not a separate negated-limit oracle endpoint. This single sparse cycle
does not make every directional criterion fail: other eligible means
are zero. Its formal endpoints establish the selected mean and derived
order, not an all-directions counterexample.

There is, however, a finite sum that covers every eligible pair. This
extension is analytic here. For every $j$ and every $i\ne j$, use the
preceding cycle with $k=j$, the same site $o$ and amplitude $A_0$, and add
all its weight functions to obtain $\overline w$. Linearity of the exact
weighted-energy sum gives $\ip{A}{R_{\overline w}A}\leq0$.
For every off-diagonal pair $i\ne j$, direct inspection of the four
indicators gives
\[
 \overline w(i,(x,j))=
 A_0\mathbf1_{x=o}-3A_0\mathbf1_{x=o+e_j}.
\]
Every diagonal component has only nonpositive contributions. For $n\geq2$
the two displayed sites are distinct, so each eligible pair has
$\overline b^+_{ij}=A_0/n^d$. With $A_0=n^d n^2$, \emph{all} off-diagonal
positive-part means equal $n^2$, while the derived domination by zero
still gives the original rank obstruction. Thus no eligible mean condition
in Theorem~\ref{thm:signed-positive} applies to this sum. The conclusion
follows from its actual energy order, not from an assumed comparison
field. The finite sum and its all-pairs mean calculation are not silently
included in the single-cycle formal oracle reports.

The cancellation mechanism is classical, rather than a new general
spectral principle. Restricted to the four active vertices, the negative
of this scalar repair has three positive path conductances $3A_0$ and a
direct negative conductance $-A_0$. For $A_0>0$ the path resistance is
$1/A_0$, so this is the critical case of Zelazo and B\"urger's
Theorem III.3~\cite{zelazo}. Isolated vertices contribute zero blocks;
the internal vector coordinates sum scalar copies. This comparison
attributes the mechanism; it does not establish originality of the
whole consolidated dossier.

The conditional obstruction is nonvacuous. An analytic check takes
$F_n=E_n$, $C_n=I$ and
$\Phi_n(y)=\|y\|^2-\ip{y}{R_w y}\geq0$ (or replace $w$ by $\overline w$).
The operator norm gives the upper quadratic bound
$(1+\|R_w\|)\|y\|^2$, while $Q_w(A)+\Phi_n(A)
=\ip{A}{K_{0,n}A}+\|A\|^2\geq\|A\|^2$.
This example has growing output, as the theorem requires; it is not an
additional formal endpoint or a compressed observation construction.
None of these signed-weight statements proves a volume-density bound,
identifies an interacting Hessian, or supplies a physical mass gap.


% Source-only mathematical section for the next consolidated candidate.
% Assembled source; exact-PDF review remains separate and version-specific.
\section{Full volume rank from signed total coefficient means}\label{sec:volume-mean}
The directional strip in Section~\ref{sec:average} detects divergence of the
observation dimension. A full dimensional coexact packet yields a positive
volume density under a stronger, explicitly stated mean hypothesis. Original
weights may have either sign. The relevant quantity is the spatial mean of
$\max\{1+w,0\}$, not an ordinary signed mean or a positive part taken after
averaging. The extra $1$ is the actual flat Hodge coefficient, already absorbed
into the envelope; it must not be added a second time.

In this section write $D\geq2$ for the number of lattice directions,
$V_n=n^D$, and $E_n$ for the actual real cochain space on
$(\mathbb Z/n\mathbb Z)^D$, with fixed $N_c\geq2$. Use field components
$0,1$ and the internal real plane representing a complex scalar from
Section~2. Put $D_i A_l(x)=A_l(x+e_i)-A_l(x)$ and retain the original form
\[
 Q_{w_n}(A)=\langle A,K_{0,n}A\rangle+
       \sum_{x,i,l}w_n(i,(x,l))\|D_i A_l(x)\|^2.
\]
The actual all-cochain flat Hodge identity gives
$Q_{w_n}(A)=\sum_{x,i,l}(1+w_n(i,(x,l)))\|D_i A_l(x)\|^2$.

\begin{theorem}[Finite full volume mean budget]\label{thm:volume-mean-finite}
Let $1\leq R$ and $2R<n$. Assume $m\geq0$ and
\begin{equation}\label{eq:total-positive-means}
 \sum_x\max\{1+w_n(i,(x,l)),0\}\leq V_n m
       \qquad(0\leq i<D,\ l=0,1).
\end{equation}
No mean bound on the other field components, original weight lower bound or
prescribed uniform pointwise cap is assumed. Let $C_n:E_n\to F_n$ be
Fr\'echet differentiable at zero, $C_n(0)=0$, with $F_n$ finite dimensional.
Suppose $\Phi_n(y)\leq K_n\|y\|^2$ for $\|y\|<\delta_n$, $\delta_n>0$,
and the original small divergence-free coercivity inequality holds with
$c\geq0$ and some radius $r_n>0$. Set
\[
 H=2R(2R+1)^{D-1},\qquad
 \Lambda(k)=4\sum_{i=0}^{D-1}\sin^2(\pi k_i/n),
 \qquad\mathcal K=\{1,\ldots,R\}\times\{-R,\ldots,R\}^{D-1}.
\]
Then
\begin{equation}\label{eq:volume-mean-finite}
 c(H-\dim_{\R}F_n)\leq m\sum_{(k,q)\in\mathcal K\times\{0,1\}}\Lambda(k)
 \leq mH\,\frac{4\pi^2 D R^2}{n^2}.
\end{equation}
\end{theorem}
\begin{proof}
For $k\in\mathcal K$ use the actual scalar character $\chi_k$ and the
coexact field with components given by its backward differences in
directions $1,0$, with opposite signs. Commuting differences give actual
divergence zero. Since $k_0\neq0\pmod n$, the normalization denominator is
positive. Counting Fourier orthogonality and both internal real quadratures
produce an orthonormal family of $H$ actual real fields, spanning $S_n$.
The condition $2R<n$ prevents aliasing of their scalar frequencies; no
scalar real-valued conjugacy constraint is imposed on the internal plane.
Every vector of $S_n$ is supported on field components $0,1$.

Define actual coefficients $v_n(i,(x,l))=\max\{1+w_n(i,(x,l)),0\}$ for
$l=0,1$, and zero otherwise, and the positive form $P_n=G^*V_n^{\rm wt}G$
with those coefficients. The notation $V_n^{\rm wt}$ here denotes the
weight multiplication operator, not lattice volume. Then $P_n\succeq0$
on all cochains and $Q_{w_n}(A)\leq\langle A,P_nA\rangle$ on $S_n$.
This comparison is needed only on the actual packet image, and is not
asserted on every cochain or every divergence-free field.

Write $L_n=DC_n(0)$. The same shrinking-ray proof as in
Section~\ref{sec:average}, now followed by the proved envelope comparison,
gives $c\|v\|^2\leq\langle v,P_nv\rangle$ on $S_n\cap\ker L_n$.
Its dimension is at least $H-\dim F_n$. Complete an orthonormal basis of
that kernel within $S_n$. Positivity of $P_n$ controls the complementary
trace terms. Thus $c(H-\dim F_n)\leq\operatorname{tr}(P_n|_{S_n})$, where
the restriction means the genuine orthogonal compression.

For an individual normalized packet $U_{k,q}$, its squared component
amplitudes are constant in $x$ and are $V_n^{-1}a_{k,l}$ for $l=0,1$,
where $a_{k,l}\geq0$ and $a_{k,0}+a_{k,1}=1$. Consequently
\[
 \langle U_{k,q},P_nU_{k,q}\rangle
 =V_n^{-1}\sum_{i=0}^{D-1}\sum_{l=0}^{1}
   |\chi_k(e_i)-1|^2 a_{k,l}\sum_x v_n(i,(x,l))
 \leq m\sum_i|\chi_k(e_i)-1|^2=m\Lambda(k).
\]
Sum over all $H$ basis vectors. Finally every signed integer frequency
satisfies $|k_i|\leq R$, and $|\sin t|\leq|t|$ gives
$\Lambda(k)\leq4\pi^2DR^2/n^2$. All $D$ derivative directions remain in
this sum even though only two field components are occupied.
\end{proof}

\begin{theorem}[One positive density on the full integer tail]\label{thm:volume-full-tail}
Fix $D,N_c$, $c>0$ and $m\geq0$. For each integer $n$ under consideration,
allow its own $F_n,C_n,L_n,\Phi_n,K_n,\delta_n,r_n,w_n$ satisfying the
original local hypotheses and~\eqref{eq:total-positive-means}, with these
same $c,m$. There are $k\geq3$ and $\varepsilon>0$, depending only on
$D,m,c$, such that
\[
 \dim_{\R}F_n\geq\varepsilon n^D\qquad(n\geq k)
\]
whenever those finite-side hypotheses hold. If the original family hypotheses hold only for $n\geq n_0$, the conclusion applies for $n\geq\max\{n_0,k\}$. One may take
$\varepsilon=1/(2k^D)$ for any integer $k\geq3$ satisfying
$2m\,4\pi^2D/k^2\leq c$. No common quadratic penalty constant,
differentiability neighborhood or coercivity radius is required.
\end{theorem}
\begin{proof}
Choose such a $k$ by the Archimedean property. For each $n\geq k$ put
$R=\lfloor n/k\rfloor$. Integer division yields
\[
 R\geq1,\qquad kR\leq n<k(R+1)\leq2kR,\qquad2R<kR\leq n.
\]
Hence the actual packet is admissible and $R/n\leq1/k$. The finite theorem
gives $c(H-\dim F_n)\leq(c/2)H$, so $H\leq2\dim F_n$. Also
$H\geq(2R)^D\geq(n/k)^D$. Combining these gives the displayed density.
Each shrinking-ray argument was applied at a fixed $n$ before using these
integer inequalities. No limit of derivatives or local neighborhoods is taken.
\end{proof}

\begin{proposition}[Exact trace average and a sharper analytic constant]\label{prop:volume-average}
For $R\geq1$, put $A_R(x)=\sum_{r=1}^R\sin^2(rx)$. The actual multiplier
average in~\eqref{eq:volume-mean-finite} is
\begin{align*}
 B_{D,R,n}&=\frac1H\sum_{k,q}\Lambda(k)
   =4A_R(\pi/n)\frac{2DR+1}{R(2R+1)},\\
 A_R(x)&=\frac R2-\frac{\sin(Rx)\cos((R+1)x)}{2\sin x},\\
 B_{D,R,n}&\leq\frac{2\pi^2}{3n^2}(R+1)(2DR+1).
\end{align*}
For fixed $0<\alpha<1/2$ and $R=\lfloor\alpha n\rfloor$, the original
family hypotheses therefore imply the analytic bound
\[
 \liminf_{n\to\infty}\frac{\dim F_n}{n^D}
 \geq(2\alpha)^D\left[1-\frac mc\,2D
       \left(1-\frac{\sin(2\pi\alpha)}{2\pi\alpha}\right)\right].
\]
Its elementary quadratic lower estimate replaces the bracket by
$1-4\pi^2Dm\alpha^2/(3c)$. This improves the low-frequency quadratic
coefficient of the maximum-coordinate estimate by a factor three.
\end{proposition}
\begin{proof}
The positive first coordinate contributes $4A_R/R$ to the average; each
of the $D-1$ centered coordinates contributes $8A_R/(2R+1)$. Both real
quadratures occur in numerator and denominator. The cosine sum follows by
summing $2\sin x\cos(2rx)=\sin((2r+1)x)-\sin((2r-1)x)$ and telescoping;
$\sin^2(rx)=(1-\cos(2rx))/2$ gives the formula for $A_R$.
Bounding $\sin^2(rx)\leq r^2x^2$ and summing $r^2$ gives the finite upper
bound, including its lower-order terms. Floor inequalities give
$H/n^D\to(2\alpha)^D$, and $\sin x/x\to1$ gives
$B_{D,R,n}\to2D[1-\sin(2\pi\alpha)/(2\pi\alpha)]$. Divide the finite
rank bound by $n^D$ and take the lower limit. Negative right sides are merely
weak bounds. For $m=0$, taking $\alpha\uparrow1/2$ gives lower density one.
\end{proof}
These are classical Fourier and finite trace arguments, with the standard
discrete-torus spectrum already present in the literature~\cite{lin-wan-zhang}. The exact average
and its improved constant are analytic calculations here, not extra formal
endpoints or sharp optimal observation-rank claims. They do not identify an
interacting background or a continuum Yang--Mills spectral gap.


\section{Why dimension is not a sufficiency criterion}
The following analytic statement is included to prevent the reverse inference
from a necessary rank budget. It is not a new Lean-verified endpoint here. The proof below is the classical
singular-Schur range condition \cite{albert,carlson}, expressed as boundedness
below on each measurement fiber; we claim no new general matrix criterion.
Let $E$ be a finite-dimensional real inner-product space, $H=H^*$ and
$B=H-cI$. Fix linear $C:E\to F$, and decompose $E=N\oplus W$ with
$N=\ker C$, $W=N^\perp$. In this decomposition write
\[
 B=\begin{pmatrix}A&D\\D^*&G\end{pmatrix}.
\]
\begin{proposition}\label{prop:schur}
A finite real-valued function $\Phi$ with $\Phi(0)=0$ can make
$\ip{x}{Bx}+\Phi(Cx)\geq0$ for every $x$ if and only if
$A\succeq0$ and $\ker A\subseteq\ker D^*$. In that event a quadratic
penalty suffices.
\end{proposition}
\begin{proof}
Taking $x\in N$ proves $A\succeq0$. For $z\in\ker A$ and fixed $w\in W$,
the value at $tz+w$ contains $2t\ip{z}{Dw}$, while $C(tz+w)=Cw$ is fixed.
Boundedness below for all real $t$ forces this coefficient to vanish.

Conversely, $\ker A\subseteq\ker D^*$ implies $D(W)\subseteq\operatorname{ran}A$.
Let $A^\dagger$ be the inverse on $\operatorname{ran}A$ and zero on $\ker A$.
Completing the square gives
\[
 \ip{n+w}{B(n+w)}=
 \ip{n+A^\dagger Dw}{A(n+A^\dagger Dw)}
 +\ip{w}{(G-D^*A^\dagger D)w}.
\]
The restriction $C|_W$ is injective; extend its inverse on $C(W)$ to a
linear map $T:F\to W$. The quadratic function
$\Phi(y)=\ip{Ty}{(D^*A^\dagger D-G)Ty}$ cancels the second term and vanishes
at zero. If $F$ has a chosen Euclidean norm, a sufficiently large nonnegative
multiple of $\|y\|^2$ also suffices, by finite-dimensional norm bounds.
\end{proof}
For a concrete counterexample let $H=\operatorname{diag}(0,2)$, $c=1$ and
$C(x,y)=x+y$. On $\ker C$, $\ip{x}{(H-I)x}=0$, and $C$ detects the negative
eigendirection of $H-I$. Nevertheless, on $(t,1-t)$ the repaired form is
$1-2t+\Phi(1)$, unbounded below for every finite $\Phi(1)$.
Kernel nonnegativity and low-mode detection alone do not suffice; cross-coupling
must be checked.

\section{Operator scaling and the actual consumer}
In the repository, the flat fixed-volume covariance endpoint uses
$K_0+aQ^*Q$ with the physical counting $\ell^2$ norm and the same unscaled
flat Hodge operator. More precisely, the declared precision shell includes
$-\Sigma$. Specializing every perturbation and its summable majorant to zero
and taking $a>0$ gives the available certified coercivity lower bound
$\min(1,a)/C_P$, with $C_P>0$ from the full-periodic block-Poincare estimate.
This is not a claim that this value is optimal. At fixed $d,N_c,N'$ the block
codomain is fixed and finite. Inspection therefore identifies the unscaled
flat base route with the setting of Theorem~\ref{thm:uniform}, with
$\Phi_L(y)=a_L\|y\|^2$ and no fine repair. Even allowing arbitrary scalar
$a_L$ or rescaling $Q_L$ does not yield a positive uniform full-space constant.
This is a source-level specialization in this text; no new Lean adapter to
that consumer is claimed. At the new source anchor below, the module
\texttt{PhysicalGaugeFixedPrecision} defines
\texttt{flatGaugeFixedPrecisionCLM}; its explicit-Poincare coercivity theorem
is recorded by its full name and specialization in the linked declaration map.

The inspected endpoint is explicitly fixed-volume. Its positive rate may
depend on $L,N'$ and $C_P$, so there is no contradiction. Nor does the
obstruction exclude all possible uniform Combes--Thomas arguments. Its
locality radius and matrix-entry bounds also vary with $L$; coercivity alone
would not settle the rate. An interacting Hessian, a growing measurement
budget, a different sector norm or a physically rescaled operator must each
be audited separately.

In particular, writing $\ell=L+1$ so that $n_L=\ell N'$, replacing $K_0$
by $\ell^2K_0$ changes the band threshold exactly to
$4\pi^2dR^2/(N')^2$. The proof above
no longer drives it below every fixed $c>0$. This observation is an escape
from this proof, not an established coercivity theorem for the programme's
physical operator. No ultraviolet activity estimate, continuum limit,
Wilson-Hessian identification or Yang--Mills mass gap is obtained.

\section{Formal evidence and publication context}
The source-tree anchor for the new physical-band and uniform-rank results is
\begin{center}\small\texttt{e31b428917aa767b4a1aad097db5cf92fb9cd23e}.\end{center}
The original factorized-kernel and repair development in
Section~\ref{sec:old} instead has the distinct tested source anchor
\begin{center}\small\texttt{7b25bfb39357406cf61880e3ac0428d7d1ba3449}.\end{center}
Its exact modules are not present in the new main tree merely because a
reproduction supplement is present there. The two full preserved focal runs
at that older SHA agree in compiled object, oracle log and source manifest;
their six named reports use only the three standard axioms. The reproduced
factorized endpoint is dimension four. The general-dimensional witness
calculation, trigonometric closed form, spectral multiplicity, liminf and
witness-line statements are analytic here; they are not all silently included
in that formalization claim. The external declaration map keeps both anchors. The public source repository is
\href{https://github.com/lluiseriksson/THE-ERIKSSON-PROGRAMME}{THE-ERIKSSON-PROGRAMME on GitHub}.
The preserved parent v5 map is
\href{https://github.com/lluiseriksson/THE-ERIKSSON-PROGRAMME/blob/d9e2a805796e6901dc9b5565f67f975cd5e41760/papers/flat-coarse-rank/v5-candidate/declaration-map.json}{\texttt{declaration-map.json}}
and the preserved evidence is
\href{https://github.com/lluiseriksson/THE-ERIKSSON-PROGRAMME/tree/e31b428917aa767b4a1aad097db5cf92fb9cd23e/reproducibility}{\texttt{reproducibility/}}. Exact source SHAs, not moving branch heads,
identify the formal claims. That parent map does not certify the new full volume extension. The map packaged with this manuscript records each added anchor and the separately reproduced full-tail extension. The accompanying exact-artifact report binds the final source and PDF hashes.

The geometric and Fourier modules are ordinary proof dependencies, not
assumptions about a newly supplied band. Each previously verified stage identified below has complete
same-source fresh-clone Colab evidence; the new extension boundary is stated separately. For the uniform endpoints the tested
source is
\begin{center}\small\texttt{17ca395d4e9de992b1c575454a112da9e6f97e86}.\end{center}
Its two runs
agree in compiled object, raw oracle log and source-manifest hashes. The
standalone budget lemma uses \texttt{propext, Quot.sound}; the other four
uniform-module reports also use \texttt{Classical.choice}. No extra axiom is
reported. Actual-log acceptance and three real-field rejection tests pass
under normal and optimized Python. These tests concern evidence handling,
not mathematical independence.

The additional whole-space gradient/Hodge identity and constructed weighted
repair bounds use tested source
\begin{center}\small\texttt{dfd49d0b0730a4d36f58a4cbb7a17d6a41fb60bc}.\end{center}
The full-space tangent endpoints in Section~\ref{sec:tangent} use tested source
\begin{center}\small\texttt{81cf07b2e0d30543682c4073a4a2d82950d638b7}.\end{center}
Each has two complete same-source fresh Colab Linux reconstructions with
matching compiled object, raw oracle and LF-source-manifest hashes. For the
tangent source, all three named endpoints use only \texttt{propext},
\texttt{Classical.choice} and \texttt{Quot.sound}. Its oracles and pinned
\href{https://github.com/lluiseriksson/THE-ERIKSSON-PROGRAMME/tree/e3bef02c87f4f6ea3392b3ed8fe8b54acb13a264/reproducibility/nonlinear-tangent-20261030}{Colab reproduction drivers}
are separate from the manuscript render. Failed source candidates and their
full logs are retained in the programme evidence archive. Later, separately reproduced sources formalize the actual implicit-function
curve construction at
\begin{center}\small\texttt{ef81f028ce6e7e8aaca7e9229eb00aa683752a41}.\end{center}
The local quadratic-energy regular-sector obstruction uses
\begin{center}\small\texttt{9bb90cfd77ec5cea5c4447667fe8baf8c21429a1}.\end{center}
The displayed sector proofs remain readable analytic arguments; their
uniform consequence is not silently included in those four-endpoint reports.
No theorem about a nonlinear sector is inferred from a full-space no-go alone.

Short labels and an external declaration map replace long wrapped Lean names
in this manuscript. The companion source map records exact names, source SHAs,
and evidence directories. Reproduction commands remain in the existing
copyable notebooks rather than being split into misleading line fragments.
The fixed-map Schur criterion, its counterexample calculation and consumer
specialization above are analytic or source-inspection arguments; the whole
manuscript is not described as machine checked.

The actual mean-weight finite/asymptotic theorem of Section~\ref{sec:average}
uses tested source
\begin{center}\small\texttt{ac60c0df885d8866329a839e2b4e73c90595d806}.\end{center}
Its two own fresh Colab Linux clones 076/076b pass six exact standard-axiom
reports and all nine stages. Compiled module SHA-256 is
\begin{center}\small\texttt{28fcc004795a1d8ab74ca87b5e67cdf90f557551bc02b1bb912da38324c26e50}.\end{center}
Raw oracle and LF-source-manifest hashes also match; full records and the
strict mutation-tested validator are preserved in the
\href{https://github.com/lluiseriksson/THE-ERIKSSON-PROGRAMME/tree/c3f2b1d31e50b37e44d895e0b08ee834bd6af488/reproducibility/average-weight-rank-divergence-20261031}{mean-weight reproduction package}.
The proof composes the actual character strip, its exact trace, positivity,
local derivative-kernel bound and multiplier estimate; these are derived
lemmas, not desired conclusions supplied as hypotheses. Nonnegative
weights, fixed positive c and the quadratic penalty upper bound remain.
The constructed quadratic-weight mean observation instead uses tested source
\begin{center}\small\texttt{7f37d6437a8f69a1d3fddf3dc0bcfad74c6bf9b6}.\end{center}
It has two same-source fresh-clone reproductions 067/067b with six exact
standard-axiom reports and matching compiled, raw-oracle and LF-source hashes.
The formal output is an actual fixed finite-dimensional type; the displayed
numeric dimension and nonoptimal path explanation are not additional oracle
endpoints. The penalty examples and regularity tradeoff are analytic only.
Classical quadratic-form, Fourier and trace mechanisms, as discussed in
\cite{spielman}, establish no new general principle or priority here.

The signed-weight results in Section~\ref{sec:signed} have distinct exact
formal anchors. The positive-part construction, actual energy order and
finite/asymptotic sufficient condition use
\begin{center}\small\texttt{810785b17b2f1bb14dd6caa5435865eb51a86ef3}.\end{center}
Its fresh own Colab Linux clones 077/077b pass six exact standard-axiom
reports. The actual zero-signed-mean fixed-output family instead uses
\begin{center}\small\texttt{497e3b4ef4046cee25b1069b8aa0fd8c0cfaa178}.\end{center}
Its clones 079/079b pass ten exact reports, including actual Hodge coercivity
and the constructed family on one fixed finite output type. The individual
square-cycle order, selected positive-part mean and rank consequence use
\begin{center}\small\texttt{4c5728d147721523f5a521011351ac4bf2d3cf86}.\end{center}
Its clones 081/081b pass ten exact reports. In every pair the compiled
module, raw oracle and LF-source-manifest hashes match, with real mutation
controls under normal and optimized Python. The declaration map records
exact names and hashes, and links the complete reproduction packages.
Failed probes 078 and 080, their sources, full diagnostics and corrections
remain preserved; no passing result is inherited across source changes.

The individual cycle has zero positive-part means in other pairs. Its
formal selected-pair identity is not an all-directions counterexample.
The finite sum covering all off-diagonal pairs, its displayed mean calculation,
the negative hyperplane diagnostic, numerical output dimension, normalized
critical-mean limit interpretation and full-state identity-observation
non-vacuity check are analytic here. They are not additional oracle endpoints.
No density conclusion from ordinary signed spatial means alone has been established.
The code indexes positive side length by $k+1$, equivalent to the positive
integer indexing used in the text.

\subsection{Evidence boundary for the full volume extension}
The actual normalized image and genuine compression trace are separately
reproduced at source \texttt{9a8946b7}, the whole-image support and constructed
positive total envelope at \texttt{88445ac6}, and the original local nonlinear
finite rank composition at \texttt{3a2d5305}. The actual numerical bound and
constructed uniform dilation density have two matching own Colab reproductions
at \texttt{031ca274}, with five exact standard-axiom reports and full manifests.
These are distinct source artifacts, not grades inherited by this manuscript.

The full integer-tail source \texttt{067fed75} now has two fresh own Colab
reproductions, 094/094b, with all nine stages exiting zero and five exact
standard-axiom reports. Their compiled module, raw oracle and LF-source
manifest hashes match. Real mutation controls pass under both normal and
optimized Python. The earlier browser loss and absent saved output are
preserved separately; no outcome is reconstructed for that earlier event.
The live session manager reported no active sessions before and after these
fresh runs. This is execution reproducibility, not independent scientific
review. The exact averaged multiplier and sinc limit remain analytic
calculations. Preparation and rendering alone do not constitute a completed
exact-artifact review.
\bigskip
\subsection*{Historical model assessments and review context}
This candidate was prepared by the Codex agent in the author's continuing
programme chat, with access to earlier objections, source audits and supplied
reproduction records. The served backend model identifier and reasoning
setting are not exposed, and no technical isolation or scientific independence
is claimed for this preparation.

The public v2 of this record has canonical PDF SHA-256
\begin{center}\small\texttt{b37f2fe86b7a8a526ad02f8b28491c300a9ff72d5653536bc5532e8c9468ed8f}.\end{center}
A subagent explicitly authorized by the author and configured as
\texttt{gpt-6-astra} assessed that exact eight-page artifact after first
preserving the earlier candidate's objections. Actual backend identity and
reasoning setting were not exposed. The workflow used the same account,
provider and programme context; it was not blind or organizationally
independent. The complete earlier and corrected-artifact reports remain in
the v2 package. The corrected-artifact report regarded O1--O5 as resolved
and retained O6: no independent complete Lean rebuild or dependency-closure
audit. No numerical score or human peer review was supplied.

The immediately preceding public artifact is AIRR
\href{https://airr.science/papers/ARR-2026-7EPM3V3DP796BBRR/versions/v5/}{v5},
a 33-page Working paper released on 4 November 2026 under CC BY 4.0,
with CC0 metadata. Its canonical PDF SHA-256 is
\begin{center}\small\texttt{69d802bfea8f8f23bfe6aba03cdc4835feac777e024c939040c7f099a11462c8}.\end{center}
It is public and citable, but has not completed AIRR screening and is not
editorially accepted or human peer reviewed. Earlier versions and their complete
reports remain preserved. The author is also the AIRR founder/operator; no
independent editorial decision is inferred from publication.

The historical revision136 corrected verified weight-symbol and cross-reference
defects in that artifact; those changes alone were editorial. This consolidated
successor additionally supplies the classical local and geometric proofs,
corrects a false one-cluster-per-leaf inference, and incorporates the separately
scoped actual physical embedding and exact finite penalty calculation. A newly rendered PDF will be a different
artifact: its hash, visual review, assessment and any necessary deposit permission
must be recorded separately. No public version number, private submission number,
license grant, historical review score or acceptance is assigned by this source.

\section{From necessary rank budgets to sufficient observations}
The preceding sections establish necessary information budgets and their
formal and analytic boundaries. The continuation below asks when actual
coefficient distributions suffice to construct observations of the same
order. It is integrated into this revision, not a second paper or deposit.
For these sections set $D=d\geq3$ and $g=N_c^2-1$, $h=Dg$.
The symbol $\mu\geq1$ denotes the scalar coefficient mean budget here;
it does not replace the earlier symbol $m=N_c^2-1$ for internal dimension.
The field is denoted $U$ instead of $A$. All norms and original forward-edge
occurrences use the same counting, unit-spacing flat convention. Restrictions
on scalar coefficients and the distinction between spectral coordinates and
rectangular means are stated explicitly below.

Historical main-text and continuation artifacts were assessed on different
PDF bytes. The immediately preceding integrated v5 now has a public exact-artifact
AI report, described in Section~\ref{sec:model-context}. That report is historical
evidence for v5, not an assessment of this changed candidate. A first exact-PDF
review must follow rendering and inspection of this candidate.
\section{Actual objects and the original observation problem}\label{ap:sec:objects}
Fix $D\geq3$, $N_c\geq2$, $c>0$, and the unit-spacing torus
$X_n=(\mathbb Z/n\mathbb Z)^D$. The actual real adjoint coefficient space has dimension
$g=N_c^2-1$; there are $D$ field components, so write $h=Dg$.
For the actual forward difference $d_i u(x)=u(x+e_i)-u(x)$, take one scalar total
coefficient $a(x)=1+w_{i,x,j}>0$ shared by all derivative directions and field components.
Use counting norms and the actual componentwise flat energy
\begin{equation}\label{ap:eq:energy}
 Q_w(U)=\E_a(U)=\sum_{x,i,j}a(x)\|d_iU_j(x)\|^2.
\end{equation}
The existing flat Hodge identity identifies these forms; its formal source is listed
in Section~\ref{ap:sec:evidence}. All continuation adapters remain analytic.
There is no substitution of a covariant or interacting derivative in \eqref{ap:eq:energy}.

The original local data consist of a differentiable map $C$ at zero, with $C(0)=0$,
into an actual finite-dimensional real normed space $F$, and a penalty satisfying
$\Phi(y)\leq K\|y\|^2$ near zero, such that
\begin{equation}\label{ap:eq:local}
 c\|U\|^2\leq Q_w(U)+\Phi(C(U))
\end{equation}
for sufficiently small actual divergence-zero fields. ``Local'' here means a
configuration-space neighbourhood; spatial support of $C$ is an additional condition.
The constant $c$ is fixed, while observation maps and neighbourhoods may depend on $n$.

A useful necessary condition follows directly. If an actual $J$-dimensional
divergence-zero subspace has $Q_w(U)\leq\lambda\|U\|^2$ with $\lambda<c$, then
$\dim F\geq J$. Otherwise a nonzero $U$ in that subspace lies in $\ker DC(0)$.
Along $tU$, $C(tU)=o(t)$ and the penalty is at most $o(t^2)$. Dividing
\eqref{ap:eq:local} by $t^2$ and taking $t\to0$ contradicts $\lambda<c$.
No rank or band certificate is supplied as an extra hypothesis in this argument.

\section{The critical global moment theorem}
Put $p=D/2>1$. For $\mu\geq1$ suppose that
\begin{equation}\label{ap:eq:moment}
 n^{-D}\sum_{x\in X_n}(\mu/a(x))^p\leq M<\infty.
\end{equation}
\begin{theorem}\label{ap:thm:critical}
Under \eqref{ap:eq:moment}, there is an actual real linear spectral observation $C$ such that
\[
 c\|U\|^2\leq Q_w(U)+c\|C(U)\|^2
 \quad\hbox{for every actual cochain }U,
\]
and
\begin{equation}\label{ap:eq:rank}
 \dim F\leq C_*\max\{1,(n/\sqrt{1+\mu})^D\},
\end{equation}
where, with $I_p=(p-1)^2/[p(p+1)]$ and
$\Gamma_p=\frac p{p-1}I_p^{-(p-1)}$, one safe choice is
\[
 C_*=h\left[1+\Gamma_p(3/4)^D D^{p+1}c^pM2^p\right].
\]
Thus $\Phi(y)=c\|y\|^2$, $K=c$ and any positive neighbourhood radii supply the original
local data. In particular the assertion holds on the actual divergence-zero subspace.
\end{theorem}
The original subclass $1\leq a\leq2\mu$ with exact arithmetic mean $\mu$ is included,
but those extra normalizations are not needed for sufficiency. The distribution
condition \eqref{ap:eq:moment} remains additional; it is not derived from an arithmetic mean.

\subsection{A finite normalized Fourier lemma}
For a finite abelian group of size $V$, let $\mathcal F$ be its counting-unitary Fourier
matrix, whose entries have modulus $V^{-1/2}$. For a frequency multiplier $g$ put
\[
 L=V^{-1}\sup_{t>0}t^p\#\{k:|g(k)|^2>t\}.
\]
For a nonnegative matrix $B$, write $\nct(s,B)$ for its number of eigenvalues strictly
greater than $s$, counted with ordinary multiplicity.
\begin{lemma}\label{ap:lem:finite}
For $K=M_f\mathcal F^*M_g\mathcal F$ and $s>0$,
\[
 \nct(s,K^*K)\leq\Gamma_p Ls^{-p}\sum_x|f(x)|^{2p}.
\]
The trace and eigenvalue counts are not volume-normalized.
\end{lemma}
\begin{proof}
Remove phases by unitaries and set $A=\mathcal F^*M_{|g|^2}\mathcal F$.
If $L=0$ the assertion is immediate. Otherwise let $B_0=pL/(p-1)$ and let $A^+$
be the inverse on $\operatorname{ran}A$, zero on its kernel. For $0\leq\gamma\leq A$
write $\rho(x)=\langle\delta_x,\gamma\delta_x\rangle$. Positivity makes $\gamma$
annihilate $\ker A$. On the range let $P_E=\mathbf1_{\{A^+\leq E\}}$ and
$R_E=\mathbf1_{\{A^+>E\}}$. The weak-frequency tail gives
\[
 \langle\delta_x,P_EAP_E\delta_x\rangle
 =V^{-1}\sum_{|g(k)|^2\geq1/E}|g(k)|^2\leq B_0E^{p-1}.
\]
Indeed the tail above $t$ is $t$ times its count plus the integral of counts above $t$;
the resulting bound is $VLp\,t^{1-p}/(p-1)$. Non-strict thresholds follow by limits.
The triangle inequality for $\gamma^{1/2}\delta_x$ split by $P_E+R_E$ yields
\[
 \langle\delta_x,R_E\gamma R_E\delta_x\rangle
 \geq\bigl(\sqrt{\rho(x)}-\sqrt{B_0}E^{(p-1)/2}\bigr)_+^2.
\]
Exact spectral layer integration consequently gives
\begin{equation}\label{ap:eq:density}
 \Tr(A^+\gamma)\geq I_pB_0^{-1/(p-1)}\sum_x\rho(x)^{p/(p-1)}.
\end{equation}
The scalar integral reduces to
$\int_0^1(1-t^{(p-1)/2})^2dt=I_p$.
Now take $P$ to be the spectral projection of $A^{1/2}M_{|f|^2}A^{1/2}$ above $s$,
with rank $\nu>0$, and set $\gamma=A^{1/2}PA^{1/2}$. Then
$0\leq\gamma\leq A$, $\Tr(A^+\gamma)\leq\nu$, and
\[
 s\nu<\sum_x|f(x)|^2\rho(x)
 \leq\Bigl(\sum_x|f(x)|^{2p}\Bigr)^{1/p}
       \Bigl(\sum_x\rho(x)^{p/(p-1)}\Bigr)^{(p-1)/p}.
\]
Insert \eqref{ap:eq:density} and raise to $p$. The constant becomes
$B_0I_p^{-(p-1)}=\Gamma_pL$. This also explicitly handles a zero-frequency kernel.
\end{proof}
The density method is classical and follows the Rumin--Frank approach
\cite[Section 2]{ap:frank}. The proof above supplies its finite normalization and kernel
adapter; it is not a claim of a new Cwikel mechanism or independent discovery.

\subsection{Actual lattice frequencies and the weighted inverse}
For $n\geq2$ let $L_0=d^*d$ and work first on the actual counting-mean-zero scalar
space $H_0$. Its Fourier symbols are
\[
 z_i(k)=e^{2\pi\mathrm i k_i/n}-1,\qquad
 \lambda(k)=4\sum_i\sin^2(\pi k_i/n).
\]
Choose centred integer representatives $r_i$ with $|r_i|\leq n/2$, using only one
representative at an even endpoint. Concavity of sine implies
\[
 \lambda(k)\geq16|r|^2/n^2,\qquad
 |g_i(k)|:=|z_i(k)/\lambda(k)|\leq n/(4|r|),\quad k\ne0.
\]
Set $g_i(0)=0$. If $|g_i|>\tau$, then $0<|r|<R=n/(4\tau)$.
For $R\leq1$ there are none, and for $R>1$ there are at most $(3R)^D$ representatives.
Thus
\begin{equation}\label{ap:eq:weak}
 n^{-D}\sup_{t>0}t^{D/2}\#\{|g_i|^2>t\}\leq(3/4)^D.
\end{equation}
This is uniform in $n$, including $n=2$, not an imported infinite-volume bound.

Let $L_a=d^*M_ad$. On $H_0$ it is positive definite. For $f\in H_0$, the flow
$j_*=M_adL_a^{-1}f$ minimizes $\langle j,M_{a^{-1}}j\rangle$ among $d^*j=f$:
write $j=j_*+k$, $d^*k=0$, and the cross term is
$\langle dL_a^{-1}f,k\rangle=0$. Inserting the explicit feasible flow
$j=dL_0^{-1}f$ gives the actual quadratic-form comparison
\begin{equation}\label{ap:eq:inverse}
 L_a^{-1}\leq\sum_{i=1}^DT_i^*T_i,\qquad
 T_i=M_{a^{-1/2}}d_iL_0^{-1}.
\end{equation}
Extend $L_0^{-1}$ by zero on constants. Lemma~\ref{ap:lem:finite} and \eqref{ap:eq:weak}
bound each $\nct(s,T_i^*T_i)$ by

\[
 \Gamma_p(3/4)^Ds^{-p}\sum_xa(x)^{-p}.
\]

For nonnegative matrices, the count of their sum above $t$ is no larger than the
sum of their counts above $t/D$: intersect the complements of their high spectral
subspaces and apply min--max. No commutation is required. With $t=1/c$ in
\eqref{ap:eq:inverse}, inversion on $H_0$ and the single scalar constant mode yield
\begin{equation}\label{ap:eq:count}
 \#\{\lambda(L_a)<c\}\leq
 1+\Gamma_p(3/4)^DD^{p+1}c^pM n^D/\mu^p.
\end{equation}
No spectral-band dimension is assumed in this derivation.

\begin{proof}[Completion of Theorem~\ref{ap:thm:critical}]
Choose an actual real counting-orthonormal eigenbasis of the real symmetric matrix
$L_a$. Retain the coordinates with eigenvalue $<c$, separately in each of the $h$
real field coordinates. Complexification used in the Fourier proof preserves real
eigenvalue multiplicities; it does not double them. Energy controls coefficients at
eigenvalues $\geq c$, while $c\|C(U)\|^2$ controls the rest. Equation~\eqref{ap:eq:count},
$p=D/2$, and $(1+\mu)/\mu\leq2$ give \eqref{ap:eq:rank}. For $n=1$ take the identity on all
$h$ coordinates. This constructs the actual global spectral observation.
\end{proof}

\section{Every lower positive inverse exponent is insufficient}
\begin{theorem}\label{ap:thm:subcritical}
For every fixed $0<r<D/2$, the analogue of Theorem~\ref{ap:thm:critical}
is false uniformly under the conditions
\[
 \operatorname{mean}((\mu/a)^r)\leq2,\qquad
 1\leq a\leq2\mu,\qquad \operatorname{mean}(a)=\mu.
\]
\end{theorem}
\begin{proof}
Choose a fixed integer $\ell>3D/c$, and put $P=2\ell+2$, $B=(2\ell+1)^D$.
The triangular window $t(s)=(1-|s|/\ell)_+$ has finite sums
\[
 T=\sum t^2=(2\ell^2+1)/(3\ell),\quad
 V=\sum(d^-t)^2=2/\ell,\quad W=\sum(d^+d^-t)^2=6/\ell^2.
\]
For a product window $\psi_y(x)=\prod_it(x_i-y_i)$ and nonzero actual adjoint $H$,
the actual curl $U_1=d^-_2\psi_y H$, $U_2=-d^-_1\psi_y H$, other components zero,
has backward divergence zero. Its field support and all nonzero forward-gradient
start sites lie in the box $y+[-\ell,\ell]^D$. Direct separable sums give
\[
 \|U\|^2=2VT^{D-1}\|H\|^2,\qquad
 \frac{\E_0(U)}{\|U\|^2}=\frac WV+(D-1)\frac VT
 =\frac3\ell+\frac{6(D-1)}{2\ell^2+1}\leq\frac{3D}\ell<c.
\]
Take sufficiently large multiples $n$ of $P$, set $\mu=n^2$ and
$J=\lfloor n^{D-2r}/B\rfloor$. These $J$ boxes fit on the separated $P$-grid
eventually, since $P^D\leq Bn^{2r}$. They satisfy
$J\geq n^{D-2r}/(2B)\to\infty$. Put $s=BJ/n^D\leq \mu^{-r}$ and set $a=1$ on
their union, $a=A=(\mu-s)/(1-s)$ elsewhere. Eventually $s\leq1/2$, so
\[
 \operatorname{mean}(a)=\mu,\quad1\leq a\leq2\mu,\quad A\geq \mu,
 \qquad\operatorname{mean}((\mu/a)^r)\leq s\mu^r+1\leq2.
\]
Each curl's gradients start only in its weak box. The actual $J$-dimensional span
has that same Rayleigh quotient, because both field and gradient supports are
disjoint. Section~\ref{ap:sec:objects} forces $\dim F\geq J$, whereas the target maximum at $\mu=n^2$
is $1$. This proves the refutation using original nonlinear data.
\end{proof}
Together the two theorems give the sharp inverse-moment exponent for this specified
family of global scalar conditions. They do not identify all sufficient coefficient
criteria, prove a physical coefficient hypothesis, or impose spatial locality on $C$.

\section{Spatial means and the critical dyadic obstruction}
For a partition $\mathcal P$ into integer rectangles, its actual normalized means are
\[
 C_{\mathcal P}(U)_{R,j}=\sqrt{|R|}\operatorname{mean}_R U_j,
 \qquad \dim F=h\#\mathcal P.
\]
A full dyadic cube tree splits each selected cube into all $2^D$ equal children and
retains unsplit siblings. The origin of the torus coordinate cube may be translated.
This class excludes arbitrary anisotropic rectangle cuts and compressed observations.

\begin{proposition}\label{ap:prop:dyadic}
There are fixed critical-moment profiles with target maximum $1$ for which every
coercive full dyadic cube-tree mean partition has $\Omega(\log n)$ leaves, even after
choosing a translated origin. The same profiles admit an explicit constant-size
rectangular mean partition satisfying the global inequality for every actual cochain.
\end{proposition}
\begin{proof}
Keep $\ell>3D/c$ fixed. Choose a dyadic $L=2^q\geq4\ell+2$, put
$\gamma=\max(1,c/4)$, $n=2^J$ with $J\geq q+1$, $\mu=\gamma n^2$, and
$S=\{0,\ldots,L-1\}^D$. Set $a=1$ on $S$ and $a=A=(\mu-s)/(1-s)$ elsewhere,
$s=L^D/n^D$. Then $1\leq a\leq2\mu$, its mean is exactly $\mu$, and
\[
 \operatorname{mean}((\mu/a)^{D/2})\leq1+L^D\gamma^{D/2}=:M_0.
\]
All these parameters except $n$ are fixed, and the target maximum is $1$.

In the hypothetical complete level $j_*=J-q$ of any translated dyadic grid, cells
have side $L$. Each coordinate interval of $S$ meets at most two such cyclic grid
intervals, so one intersection has integer length at least $L/2\geq2\ell+1$.
Taking these intersections independently places a whole gradient-start box for a
compact curl inside $S$ and inside one level-$j_*$ cell. If the proposed partition
had maximum depth at most $j_*$, that cell would lie in one actual leaf. The curl
then has zero means on every leaf, since each component has total sum zero, yet its
weighted Rayleigh quotient is $<c$. This contradicts coercivity. A single zero-extended
window needs no packing divisibility assumption on $n$; $n\geq2L$ suffices here.
The field may depend on the partition, since coercivity is asserted for every field.
Thus maximum depth is at least $J-q+1$, and the full-arity tree identity gives
\begin{equation}\label{ap:eq:dyadiclower}
 \#\mathcal P\geq1+(2^D-1)(J-q+1).
\end{equation}

For a successful unrestricted partition, isolate all $L^D$ vertices of $S$ as
singletons, and use the $D$ disjoint complement rectangles
\[
 R_i=[0,L-1]^{i-1}\times[L,n-1]\times[0,n-1]^{D-i}.
\]
Coverage follows from the unique first coordinate outside $[0,L-1]$.
Every internal edge starts at a strong vertex with coefficient $A$.
On a length-$l$ interval, counting variance and path Cauchy--Schwarz give
\[
 \sum\|u-\bar u\|^2\leq\frac{l(l-1)}{4A}\sum A\|du\|^2.
\]
Indeed the path resistance is at most $(l-1)/A$, and a cut at $t$ is crossed by
$t(l-t)\leq l^2/4$ pairs in the variance identity.
Product-variance tensorization on original fibres yields a rectangle constant
$n^2/(4A)$, with no extra factor $D$. Singletons have zero residual, internal edge
sets are disjoint, and omitted edges have positive energy. Since
$cn^2/(4A)\leq c/(4\gamma)\leq1$, exact orthogonal decomposition proves
\[
 c\|U\|^2\leq Q_w(U)+c\|C_{\mathcal P}(U)\|^2,
 \qquad \#\mathcal P=L^D+D.
\]
These rectangles have potentially unbounded aspect ratio; they leave the cube-tree class.
\end{proof}
The general critical moment question for unrestricted rectangular means remains open.
Changing only the stopping rule within full cube trees cannot resolve it. Classical
weighted Sobolev inequalities \cite{ap:ads} also yield block observations for inverse
moments of finite order strictly greater than $D/2$; the dossier contains the explicit
balanced-box proof. Proposition~\ref{ap:prop:dyadic} explains why its critical tree-counting
loss cannot be removed merely by choosing a better cube-tree stopping criterion.

\section{What a spatial observation proof must still establish}
\label{ap:sec:spatial-obstructions}
The preceding positive spectral result does not construct a rectangular mean map.
Three further checks prevent substitutions for that missing construction. They concern
local geometry, the penalty coefficient, and support of observations, respectively.
They are collected here as limitations of the same dossier, not as additional papers.

\subsection{Critical inverse mass does not remove rectangle geometry}
For a rectangle $R$ of side lengths $l_i$, set $V=\prod_i l_i$, $L=\max_i l_i$,
$\kappa=L^2/V^{2/D}$ and $\eta_R=\sum_R a^{-D/2}$. The endpoint Sobolev estimate
on balanced cubes, with $q=2D/(D+2)$, transfers to arbitrary rectangles by repeated
coordinate reflection and gives
\begin{equation}\label{ap:eq:aspect}
 \sum_R|u-\bar u_R|^2\leq A_D D^{2/D}\kappa\eta_R^{2/D}\E_{a,R}(u).
\end{equation}
Here the energy uses original internal edges and their actual starting-vertex weights.
The full reflection multiplicities and graph adapters to the classical estimate
\cite[Section 2.2]{ap:ads} are proved here in Section~\ref{sec:self-contained-box-sobolev}; this is an
analytic adapter, not a new Lean theorem.

The factor cannot be discarded even for uniform coefficients. With $u(x)=x_i-(L-1)/2$
along a longest side, direct sums give
\[
 \sum_Ru^2=V(L^2-1)/12,\qquad
 \E_{a,R}(u)=aV(L-1)/L.
\]
Their ratio is $L(L+1)/(12a)\geq\kappa\eta_R^{2/D}/12$.
On the thin rectangle of lengths $(L,1,\ldots,1)$, choosing $a=L^{2/D}$ makes
$\eta_R=1$ while this ratio diverges. Thus compressing a cube-tree chain into thin
rectangles cannot simply reuse a shape-independent inverse-mass certificate.
This prescribed-block test does not defeat every global partition: omitted crossing
edges and a different partition may help. The successful partition in
Proposition~\ref{ap:prop:dyadic} is not contradicted.

\subsection{Fixed and free penalties are different obligations}
Let $C$ be normalized means on a fixed partition of a finite real Euclidean space.
Then $CC^*=I$; $P=C^*C$ is the orthogonal block-mean projection and $Q=I-P$.
For the actual positive stiffness $L_a$, write the blocks on
$\operatorname{ran}P\oplus\operatorname{ran}Q$ as
\[
 A=PL_aP,\qquad B=PL_aQ,\qquad D_0=QL_aQ.
\]
Positivity implies $\ker A\subseteq\ker B^*$: otherwise the energy of $tp+v$
would have a nonzero linear coefficient with zero quadratic coefficient in $t$.
Completing the square with the finite spectral pseudoinverse $A^\dagger$ shows
\begin{equation}\label{ap:eq:shorted}
 L_a+cP\geq cI\quad\Longleftrightarrow\quad
 D_0-B^*A^\dagger B\geq cI\ \hbox{on }\operatorname{ran}Q.
\end{equation}
Indeed the first inequality is exactly $L_a-cQ\geq0$, and its quadratic form is
$\langle p+A^\dagger Bv,A(p+A^\dagger Bv)\rangle+
\langle v,(D_0-cI-B^*A^\dagger B)v\rangle$.
This is classical shorted-operator algebra \cite{ap:anderson}; it supplies no spatial
partition or target count. In particular $D_0\geq cI$ alone is insufficient.

\begin{proposition}\label{ap:prop:formal-cycle}
On the four-cycle with edge weights $(1,6,4,1)$ and blocks $\{0,1\},\{2,3\}$,
coercivity holds on zero block sums with $c=19/5$, but $L_a+cP\geq cI$ fails.
\end{proposition}
\begin{proof}
For $v=(x,-x,y,-y)$, with the counting Euclidean norm,
\[
 \E_a(v)-\frac{19}{5}\|v\|^2
 =\frac{17}{5}\left(x+\frac{35}{17}y\right)^2+\frac{84}{85}y^2\geq0.
\]
The actual normalized observation is
\[
 Cu=\bigl((u_0+u_1)/\sqrt2,\ (u_2+u_3)/\sqrt2\bigr).
\]
For $u=(19,-9,-5,-5)/14$, direct exact evaluation gives
\[
 \E_a(u)=52/7,\quad\|u\|^2=123/49,\quad\|Cu\|^2=25/49,
\]
and hence $\E_a(u)+(19/5)\|Cu\|^2-(19/5)\|u\|^2=-6/35$.
The weights have mean $3$ and lie between $1$ and $6$. Ten focal Lean theorems
preserve the constructed weights, normalized observations, projection identities,
kernel square identity and exact negative witness; Section~\ref{ap:sec:evidence}
identifies the successful immutable source. Those ten endpoints concern the scalar four-cycle. Section~\ref{sec:physical-cycle}
separately derives its actual physical embedding and the all-field kernel
premise; only its stated fixed-penalty endpoints have paired formal evidence.
\end{proof}

For the original class allowing an arbitrary finite $K$, put $J=D_0-cI$.
The necessary and sufficient condition is instead
\begin{equation}\label{ap:eq:free-penalty}
 \exists K\geq0:\ L_a+KP\geq cI
 \quad\Longleftrightarrow\quad J\geq0\ \hbox{and }B\ker J=0.
\end{equation}
Necessity follows by testing $v\in\ker C$, then $v+tp$ with $v\in\ker J$.
For sufficiency, $B\ker J=0$ implies $B^*p\in\operatorname{ran}J$.
Completing the square in $v+J^\dagger B^*p$ leaves
$A+(K-c)I-BJ^\dagger B^*$, so $K=c+\|BJ^\dagger B^*\|$ suffices.
In particular a proved all-cochain margin $D_0\geq(c+\epsilon)I$ gives
$K=c+\|L_a\|^2/\epsilon$. For actual $a\leq2\mu$, the elementary edge estimate
$|r-s|^2\leq2|r|^2+2|s|^2$ gives $\|L_a\|\leq8D\mu$, including the parallel
edge occurrences at $n=2$. A margin $2c$ therefore permits
$K=c+64D^2\mu^2/c$, without increasing the number of observations.
No margin or partition is proved by this conditional implication. Ordinary block
averaging need not preserve divergence, so an estimate only on the physical kernel
cannot silently replace the all-cochain hypothesis in this decomposition.

Even free $K$ needs care at equality. On the same cycle let
$s=\sqrt{85}$, $c_*=(17-s)/2$ and $v=(7,-7,6-s,s-6)$. Its kernel defect at $c_*$
is zero, since for general $(x,-x,y,-y)$ it equals
$(s-6)[x+7y/(s-6)]^2$. For
$p=(\delta/2,\delta/2,-\delta/2,-\delta/2)$ the full defect is
\[
 \E_a(v+p)+K\|C(v+p)\|^2-c_*\|v+p\|^2
 =(7+K-c_*)\delta^2-10(13-s)\delta.
\]
For each finite $K\geq0$ a sufficiently small positive $\delta$ makes it negative.
This additional equality example is exact analytic algebra and finite diagnostic
evidence, not one of the ten Lean theorems of Proposition~\ref{ap:prop:formal-cycle}.

\subsection{Sparse point samples cannot replace stable averages}
For uniform $a=\mu$, if an observation depends only on vertex values on a set $S$, $|S|=s$, and
$C(0)=0$, then the scalar field $u=1$ off $S$ and $0$ on $S$ has zero observation.
At most $2Ds$ forward-edge occurrences cross $S$, so
$\E_\mu(u)\leq2D\mu s$. Consequently all-cochain coercivity requires
\[
 s\geq\frac{cn^D}{2D\mu+c}.
\]
For an actual divergence-zero test, fix a component direction $j$ and let the field
be constant along every $j$-line, equal to zero on each line meeting $S$ and one on
the other lines, in one fixed internal coordinate. There are at most $s$ hit lines.
Its backward divergence is zero, its norm squared is at least
$n(n^{D-1}-s)$, and transverse crossing energy is at most $2(D-1)\mu ns$.
Thus necessarily
\[
 s\geq\frac{cn^{D-1}}{2(D-1)\mu+c}.
\]
If every transverse line is hit, the second bound is immediate without a nonzero test. These bounds give the required obstruction:
at $a=\mu=n^2$ and critical moment $M=1$, bounded samples fail for all cochains
when $D\geq3$ and for these physical tests when $D\geq4$.
For the same uniform coefficients one global mean succeeds for $0<c\leq16$,
by the nonzero Fourier gap $n^2\lambda(k)\geq16$.
This comparison rules out a point-sampling shortcut, not rectangular averages or
singletons supplemented by complementary block means. A nonlinear local observation
is tested on sufficiently small multiples of these fields.

\section{Constructive rectangular bounds and their limits}\label{ap:sec:spatial-new}
Retain $\mu\geq1$, $1\leq a\leq2\mu$, exact mean $\mu$, \eqref{ap:eq:moment},
and $p=D/2$. All assertions in this section remain analytic.

\subsection{A general threshold construction}
For an interval of length $l$, the identity
$\operatorname{Var}(u)=l^{-1}\sum_{x<y}|u(x)-u(y)|^2$ and path
Cauchy--Schwarz give coefficient at most $l(l-1)/4$ per edge: each
crossing pair has distance at most $l-1$, and at most $l^2/4$ pairs cross
any edge. Product variance tensorization therefore yields, for a rectangle
of maximal side $L$,
\begin{equation}\label{ap:eq:boxvar}
 \sum_R|u-\operatorname{mean}_R u|^2
 \leq \frac{L(L-1)}4\sum_{\text{internal forward edges of }R}|du|^2.
\end{equation}
There is no additional factor $D$ in this tensorization.

Fix $\tau\geq1$. If $\tau<c/4$ use singletons. Otherwise set
$\ell=\min\{n,\lfloor2\sqrt{\tau/c}\rfloor\}$ and start with a grid
whose sides have length at most $\ell$. Isolate each vertex with $a<\tau$
as a singleton. Its current rectangle minus that vertex splits into at most
$2D$ rectangles by the first coordinate differing from that vertex. Delete
empty pieces and preserve earlier singletons. Remaining rectangles have
$a\geq\tau$ and side at most $\ell$. Equation~\eqref{ap:eq:boxvar},
nonnegative omitted edges and normalized-mean orthogonality give
\begin{equation}\label{ap:eq:globalmeans}
 c\|U\|^2\leq\E_a(U)+c\|C_{\mathcal P}U\|^2,
 \quad (C_{\mathcal P}U)_{R,j}=\sqrt{|R|}\operatorname{mean}_R U_j
\end{equation}
for ALL cochains. Singleton residuals vanish, including $n=1$; parallel
edge occurrences for $n=2$ are counted separately, and wrap edges can be
omitted positively. With $B_c=(1+\sqrt c)^D$,
\[
 |\mathcal P|\leq B_c\max\{1,n^D/\tau^p\}+2D\#\{a<\tau\},\quad
 \#\{a<\tau\}\leq M n^D\tau^p/\mu^p.
\]
The singleton branch satisfies this count since $B_c\tau^{-p}>1$ there.
Taking $\tau=\sqrt \mu$ yields the weaker general bound
\[
 \dim F\leq h(B_c+2DM)\max\{1,n^D/\mu^{D/4}\}.
\]

\subsection{Every global threshold can lose the target order}
Let $\gamma=\max\{1,c/4\}$, $n=t^2$, $t\geq2$, $\mu=\gamma n^2$,
$b=\gamma n$, $S=\{0,\ldots,t-1\}^D$, and $f=n^{-p}$. Set $a=b$ on
$S$, $a=A=(\mu-fb)/(1-f)$ outside. Then $1\leq a\leq2\mu$, the mean is
$\mu$, and the critical moment is less than $2-f$. The target maximum is one.
For any global $\tau\geq1$, consider the preceding algorithm with any
translated initial grid or excision order, retaining grid boundaries and
every excised singleton. If $\tau\leq b$, no vertex is weak and the grid
volume bound gives
\[
 |\mathcal P|\geq n^D/(2\sqrt{b/c})^D
     =(c/(4\gamma))^p n^{D/2},
\]
or the singleton branch is larger. For large $n$ the grid side is less
than $n$. If $\tau>b$, all $t^D=n^{D/2}$ weak vertices remain distinct
singletons. Every such threshold therefore loses the target order.

For the SAME profile retain the whole weak cube and split its complement
by the first coordinate outside it, giving $D+1$ rectangles. On $S$,
$ct^2/(4b)=c/(4\gamma)\leq1$; outside, $A\geq \mu$ and
$cn^2/(4A)\leq1$. Equation~\eqref{ap:eq:boxvar} proves
\eqref{ap:eq:globalmeans} for this constant-size partition. This is a
mechanism obstruction, not a counterexample to every rectangular partition.

\subsection{Supplied scaled cubes attain the target order}
Suppose additionally that disjoint integer cubes $S_k$ of side $L_k$,
nonwrapping in one common root, satisfy for fixed $\beta\geq1$
\begin{equation}\label{ap:eq:scaled-assumptions}
 \begin{split}
 cL_k(L_k-1)/4\leq a(x)\leq\beta L_k^2 &\quad(x\in S_k),\\
 a(x)\geq cn(n-1)/4 &\quad(x\notin\cup_k S_k).
 \end{split}
\end{equation}
These are extra coefficient and geometric assumptions. Each cluster charges
at least $\beta^{-p}$ to $\sum_xa(x)^{-p}$; thus
$J\leq\beta^p M n^D/\mu^p$.

Section~\ref{sec:explicit-cube-partition} supplies a self-contained construction
for precisely these disjoint integer cubes. Use quarter-offset solid cubes,
then a compressed corner-potential subdivision. Its at most
$2^{D+1}(2^D-1)J$ terminal rectangles have no cube corner in their interiors.
This does not imply that a terminal rectangle meets only one cube: the
proper-intersection and covering bounds give at most
$T_D=2^{2D-1}$ cubes per terminal rectangle. Refining by all their facet
coordinates gives at most $(2T_D+1)^D$ subrectangles per terminal rectangle.
Consistent boundary assignment followed by taking integer points yields an
actual partition; every surviving rectangle is inside one supplied cube
or outside all supplied cubes. Fragmentation is permitted. Consequently,
with the explicit dimension-only constant
\[
 A_D=2^{D+1}(2^D-1)(2^{2D}+1)^D,
 \qquad |\mathcal P|\leq A_D\max\{1,J\},
\]
\[
 \dim F\leq h A_D(1+2^p\beta^pM)\max\{1,n^D/(1+\mu)^p\}.
\]
The earlier one-cluster-per-leaf inference is not used. The multiple-object
stress family and the boundary adapter are included in the cited section.
Inside a cluster fragment maximal side is at most $L_k$, outside at most
$n$. Equations~\eqref{ap:eq:scaled-assumptions} and~\eqref{ap:eq:boxvar} prove
\eqref{ap:eq:globalmeans}. Projection separation and a high-budget restriction
are unnecessary. This uses classical geometry; it does not derive the
supplied cubes from the original moment.

\subsection{The exact cube decomposition is not universal}
For fixed $c,\beta$ take arbitrarily large integer $n\geq3$, $\mu=n$,
$W=\{e_1,-e_1\}$, $f=2/n^D$, with $a=1$ on $W$ and
$a=A=(n-f)/(1-f)$ elsewhere. Then $n<A\leq2n$, mean $a=n$, and the
critical moment is less than $2n^{-p}+1-f\leq2$. Choose also
$n>\max\{\beta,8\beta/c,1+8/c\}$ and $n\geq c$.
Every vertex has $a<cn(n-1)/4$, so any family satisfying
\eqref{ap:eq:scaled-assumptions} must cover zero. A side-one cube containing
zero would require $A\leq\beta<n$. A cube of side $L\geq2$ contains at
least one of $e_1,-e_1$, in every translated root. Its lower bound forces
$L^2\leq8/c$, while its upper bound at zero requires
$A\leq\beta L^2\leq8\beta/c<n<A$. Even allowing overlap cannot repair
this exact decomposition.

The SAME profile nevertheless admits a strong grid with
$\ell=\min\{n,\lfloor2\sqrt{A/c}\rfloor\}$ and two singleton excisions.
Equation~\eqref{ap:eq:globalmeans} holds and
\[
 \dim F\leq h\bigl[2^p(1+\sqrt c)^D+4D\bigr]
                  \max\{1,n^D/(1+\mu)^p\}.
\]
There are at most $\lceil n/\ell\rceil^D+4D$ blocks. If $\ell<n$,
$\ell\geq\sqrt{A/c}$; if $\ell=n$, there is one grid box. Remaining
rectangles have coefficient $A$ and weak residuals are singletons. The
universal supplied-cube route is refuted. The original unrestricted
rectangular question is neither proved nor refuted.


% One article: self-contained proofs supporting the preceding claims.
\paragraph{Self-contained proof continuations.}
The next three sections complete the local and supplied-geometry arguments
and exhibit one actual physical observation map. The constants denoted $A_D$
are local to their sections: the Sobolev and partition constants differ.
Counting measure and unit lattice spacing are used throughout. These
continuations do not derive supplied geometry from the inverse moment.
% Proof fragment for the existing flat-coarse-rank dossier; no standalone paper.
% All norms below use counting measure. Analytic proof, not Lean verified.
\section{A finite proof of the critical rectangle estimate}
\label{sec:self-contained-box-sobolev}
This supplement derives the balanced-cube endpoint used by the rectangle
adapter, rather than importing a local Sobolev estimate as an external input.
The method is classical: finite entropy proves the projection inequality,
which yields relative edge isoperimetry, layer sets, and a power argument.
We claim neither a new general Sobolev principle nor a solution of the open
critical rectangular-partition problem. This is part of the existing dossier.

Fix $D\geq3$, and put
\[
 \begin{gathered}
 p=D/2,\quad q=\frac{2D}{D+2},\quad q'=\frac{2D}{D-2},\\
 \alpha=\frac{2(D-1)}D,\quad c_D=D(1-2^{-1/D}),\quad
 K_D=\frac{2\alpha D^{1/q'}}{c_D}.
 \end{gathered}
\]
An internal axis edge is an unordered adjacent pair, counted once. For its
positive-axis orientation $e=(x,x+e_i)$, write $du(e)=u(x+e_i)-u(x)$.

\subsection{Finite projection inequality, with proof}
For a nonempty finite $A\subset\prod_{j=1}^D X_j$, let $\pi_i$ delete
coordinate $i$. Then
\begin{equation}\label{eq:finite-projection}
 |A|^{D-1}\leq\prod_{i=1}^D|\pi_i A|.
\end{equation}
Here is a finite proof, including the entropy facts needed for this use.
For a finite probability vector $r$, define $H(r)=-\sum r\log r$ with
$0\log0=0$. Concavity follows from $(-t\log t)''=-1/t<0$ and continuity at
zero. Consequently mixing conditional distributions cannot decrease their
average entropy: $H(Y\mid Z)\geq H(Y\mid Z,W)$. Factoring finite joint
probabilities proves the chain rule
$H(X)=\sum_jH(X_j\mid X_1,\ldots,X_{j-1})$.
Also $H(Y)\leq\log|\operatorname{supp}Y|$: Jensen's inequality for $\log$
gives $\sum_y r_y\log(1/r_y)\leq\log\sum_y r_y/r_y$ on the positive support.
For $X$ uniform on $A$, $H(X)=\log|A|$, while
\[
 \begin{split}
 \sum_i H(X_{-i})
 &=\sum_i\sum_{j\ne i}H(X_j\mid X_{<j}\setminus X_i)\\
 &\geq\sum_i\sum_{j\ne i}H(X_j\mid X_{<j})
 =(D-1)H(X).
 \end{split}
\]
The notation in the conditioning set means omit $X_i$ only when $i<j$.
Bounding each $H(X_{-i})$ by $\log|\pi_i A|$ and exponentiating proves
\eqref{eq:finite-projection}. No projection inequality is an unproved
hypothesis in the sequel.

\subsection{Relative isoperimetry on a finite cube}
Let $Q_s=\{0,\ldots,s-1\}^D$, $s\geq2$, and let $\partial_E A$ consist
of internal cube edges crossing $A$ and its complement. If
$0<|A|\leq |Q_s|/2$, then
\begin{equation}\label{eq:relative-edge-iso}
 |\partial_E A|\geq c_D|A|^{(D-1)/D}.
\end{equation}
Indeed, put $a=|A|$, $P_i=|\pi_i A|$. At most $a/s$ occupied fibers in
direction $i$ are full. Every occupied nonfull fiber contains an internal
crossing edge, because the fiber is a connected interval. Edges in distinct
fibers/directions are different. Thus projection inequality and AM--GM give
\[
 |\partial_E A|\geq\sum_i(P_i-a/s)
 \geq D a^{(D-1)/D}-Da/s
 \geq D(1-2^{-1/D})a^{(D-1)/D}.
\]
Empty $A$ satisfies the same inequality trivially. The boundary is the EDGE
boundary, not the vertex boundary in a separate graph theorem.

\subsection{Endpoint counting-norm Sobolev--Poincare inequality}
For every real $u$ on $Q_s$, including $s=1$,
\begin{equation}\label{eq:self-contained-cube}
 \sum_{Q_s}|u-\bar u_{Q_s}|^2
 \leq K_D^2\left(\sum_{e\subset Q_s}|du(e)|^q\right)^{2/q}.
\end{equation}
For $s=1$ both sides vanish. For a nonnegative $v$ supported on at most
half the cube, layer-set decomposition and the triangle inequality in
$\ell^{D/(D-1)}$, followed by \eqref{eq:relative-edge-iso}, imply
\[
 \|v\|_{D/(D-1)}
 \leq\int_0^\infty|\{v>t\}|^{(D-1)/D}\,dt
 \leq c_D^{-1}\int_0^\infty|\partial_E\{v>t\}|\,dt
 =c_D^{-1}\sum_e|dv(e)|.
\]
The integrals are finite sums of step lengths between values of $v$;
the last equality follows edgewise from the length of the interval
between its two endpoint values. Apply this inequality to $v^\alpha$.
For $a,b\geq0$, the fundamental theorem of calculus gives
\[
 |a^\alpha-b^\alpha|
 \leq\alpha|a-b|(a^{\alpha-1}+b^{\alpha-1}).
\]
Since $\alpha D/(D-1)=2$ and $(\alpha-1)q'=2$, Holder and the degree
bound $\deg_{Q_s}(x)\leq2D$ give
\[
 \begin{split}
 \|v\|_2^\alpha
 &\leq\frac\alpha{c_D}\Big(\sum_e|dv(e)|^q\Big)^{1/q}
 \Big(\sum_{e=\{x,y\}}(v(x)^{\alpha-1}+v(y)^{\alpha-1})^{q'}\Big)^{1/q'}\\
 &\leq\frac{2\alpha D^{1/q'}}{c_D}
 \Big(\sum_e|dv(e)|^q\Big)^{1/q}\|v\|_2^{\alpha-1}.
 \end{split}
\]
The second inequality uses $(a+b)^{q'}\leq2^{q'-1}(a^{q'}+b^{q'})$ and
$\sum_e(v(x)^2+v(y)^2)\leq2D\|v\|_2^2$.
If $v=0$ the conclusion is immediate; otherwise division proves
$\|v\|_2\leq K_D\|dv\|_q$.

Choose a counting median $m$ of $u$ and set $v_+=(u-m)_+$,
$v_-=(m-u)_+$. Both supports have size at most half the cube. For each edge,
\[
 |dv_+(e)|^q+|dv_-(e)|^q\leq|du(e)|^q.
\]
For endpoint values on the same side of $m$ this is equality in the
nonzero part; on opposite sides it is $a^q+b^q\leq(a+b)^q$.
Since $2/q\geq1$, addition of the two squared estimates gives
\[
 \|u-\bar u_{Q_s}\|_2^2\leq\|u-m\|_2^2
 \leq K_D^2\big(\|dv_+\|_q^q+\|dv_-\|_q^q\big)^{2/q}
 \leq K_D^2\|du\|_q^2.
\]
This proves \eqref{eq:self-contained-cube} without an external graph,
ball-adapter, coefficient, or boundary-condition hypothesis.

\subsection{Reflection with all aspect-ratio multiplicities}
Let $R=\prod_i\{0,\ldots,l_i-1\}$, $V=\prod_i l_i$,
$L=\max_i l_i$, and $\kappa=L^2/V^{2/D}$. If $V=1$ the desired variance
bound is zero. Otherwise $L\geq2$. Put $s=2L+1$, $r=s^D/V$ and fold
$Q_s$ coordinatewise onto $R$ by the endpoint-repeating pattern of period
$2l_i$,
\[
 0,1,\ldots,l_i-1,l_i-1,\ldots,1,0,0,1,\ldots.
\]
For $l_i=1$ the fold is constant. Since $s/(2l_i)>1$, each original
coordinate value has at least $2\lfloor s/(2l_i)\rfloor\geq s/(2l_i)$
preimages and at most $s/l_i+2\leq3s/l_i$ preimages. An original internal
edge has at most the same upper count in its own direction; each full
period traverses it twice. Adjacent folded values either coincide or are
an original internal edge, possibly with reversed orientation. Thus for
the unweighted extension $\widetilde u$,
\[
 \operatorname{Var}_{Q_s}(\widetilde u)\geq2^{-D}r\operatorname{Var}_R(u),
 \qquad
 \sum_{e\subset Q_s}|d\widetilde u(e)|^q
 \leq3^D r\sum_{e\subset R}|du(e)|^q.
\]
The variance inequality holds by applying the preimage lower bound at the
mean of $\widetilde u$, and then minimizing over constants on $R$.
Apply \eqref{eq:self-contained-cube}. Since $2/q-1=2/D$ and $s\leq3L$,
\begin{equation}\label{eq:rectangle-unweighted-derived}
 \operatorname{Var}_R(u)\leq A_D\kappa
 \left(\sum_{e\subset R}|du(e)|^q\right)^{2/q},\qquad
 A_D=K_D^2 2^D 3^{2D/q+2}.
\end{equation}
The factor $r$ is not bounded independently of aspect ratio; it was kept
until it became $s^2/V^{2/D}$. Singleton coordinate sides cause no problem.

\subsection{Positive original edge weights and field coordinates}
For arbitrary positive ORIGINAL edge coefficients $b_e$, Holder on $R$
with exponents $2/q$ and $2/(2-q)$ gives
\[
 \left(\sum_e|du|^q\right)^{2/q}
 \leq\Big(\sum_e b_e|du|^2\Big)
 \Big(\sum_e b_e^{-q/(2-q)}\Big)^{(2-q)/q}.
\]
Here $q/(2-q)=D/2=p$. Consequently
\begin{equation}\label{eq:rectangle-weighted-derived}
 \operatorname{Var}_R(u)\leq A_D\kappa
 \Big(\sum_{e\subset R}b_e^{-D/2}\Big)^{2/D}
 \sum_{e\subset R}b_e|du(e)|^2.
\end{equation}
No reflected weights were introduced: reflection was completed before
this Holder step. For $b_{(x,x+e_i)}=a(x)>0$, each vertex supplies at most
$D$ internal starting edges. Thus the inverse edge mass is at most
$D\sum_Ra^{-D/2}$, proving exactly the scalar-start estimate
\[
 \operatorname{Var}_R(u)\leq A_D D^{2/D}\kappa
 \Big(\sum_Ra(x)^{-D/2}\Big)^{2/D}\mathcal E_{a,R}(u).
\]
For any fixed finite real internal/field dimension, apply this scalar
inequality to each coordinate and sum; its coefficient factor is common.
There is no additional internal-dimension factor and no replacement of
original edges by enlarged-region energy.

\subsection{Necessity of geometry and status}
For uniform $a>0$ and a longest coordinate of length $L\geq2$, put
$u(x)=x_i-(L-1)/2$. Direct sums give
\[
 \operatorname{Var}_R(u)=V(L^2-1)/12,\qquad
 \mathcal E_{a,R}(u)=aV(L-1)/L,
\]
so their ratio is $L(L+1)/(12a)\geq L^2/(12a)$. Since
$\kappa(\sum_Ra^{-D/2})^{2/D}=L^2/a$, the geometric dependence cannot
be uniformly dropped even for uniform weights. This is not a global
counterexample to a partition algorithm or to the physical divergence-zero
observation problem.

For attribution and comparison see S.~Andres, J.-D.~Deuschel and M.~Slowik,
\emph{Harnack inequalities on weighted graphs and some applications to the
random conductance model}, arXiv:1312.5473v6, Section~2.2, (2.5)--(2.8).
Their local normalized inequalities are classical prior context; this
proof does not import their theorem, graph assumptions or Harnack conclusions.
The finite projection/entropy, coarea and power steps above are classical,
not a new general method. The gain is a complete finite derivation of the
specific input previously left external. No Lean execution, independent
review, original-priority claim, new physical hypothesis removal or new
AIRR paper follows from this supplement. The frozen PDF candidate and its
actual reports/permissions remain unchanged; integration would create a
new exact artifact and must not inherit their assessments.

% Supplement to the existing scaled-cluster result, not a new standalone paper.
\section{A finite rectangular partition of disjoint lattice cubes}
\label{sec:explicit-cube-partition}
The linear partition for supplied disjoint cubes can be derived directly,
without importing a general fat-object or low-density theorem. We use the
classical compressed cube-splitting strategy of de Berg, UU-CS-1998-12,
Section~2.1 and Lemma~2.1, and prove the required special-case geometry and
all discrete/energy adapters below. This is not an original general BSP
method. No time-complexity, measured implementation or Lean claim is made.

Fix $D\geq3$ and an integer $n\geq1$. In the common root
$X=\{0,\ldots,n-1\}^D$, take pairwise disjoint nonempty integer cubes
\[
 S_k=\prod_i\{b_{ki},\ldots,b_{ki}+L_k-1\},\qquad 1\leq k\leq J.
\]
Set
\[
 Q_k=\prod_i[b_{ki}-1/4,b_{ki}+L_k-3/4].
\]
These are positive-side closed solid cubes, including $L_k=1$. They are
strictly disjoint: disjoint integer intervals in a separating coordinate
leave a real gap at least $1/2$. Their integer points are exactly $S_k$,
and each such point is interior to $Q_k$. All corners of the $Q_k$ lie
strictly inside $[-1,n]^D$.

\subsection{Proper-intersection corner bound}
Let $C$ be any axis-aligned cube of side $s>0$ with no corner of any $Q_k$
in its interior. Then at most $2^D$ of the $Q_k$ meet $\operatorname{int}C$.
Indeed, if $Q_k$ of side $a<s$ properly meets $C$, each of its coordinate
intervals has an endpoint in the interior of the corresponding interval
of $C$. Choosing these endpoints gives a corner of $Q_k$ in
$\operatorname{int}C$, a contradiction. Thus every properly meeting $Q_k$
has $a\geq s$. In each coordinate its interval then contains at least
one endpoint of the interval of $C$; otherwise it would lie strictly
inside that shorter interval. Hence $Q_k$ contains a corner of $C$.
Strict disjointness makes these corner choices injective, proving the bound.

The qualifier ``properly'' is essential. Closed boundary-only contact
must not be counted here. For $C=[0,1]^D$ and any integer $K\geq1$, put
$\epsilon=1/(4K+4)$ and
\[
 Q_j=[-\epsilon,0]\times[(2j+1)\epsilon,(2j+2)\epsilon]
       \times[\epsilon,2\epsilon]^{D-2},\qquad0\leq j<K.
\]
These $K$ strictly disjoint small cubes touch $C$ but none meets its
interior or has a corner in its interior. Thus the corresponding
closed-contact bound would be false. This boundary guard does not assert
that de Berg's theorem is false; our construction below explicitly uses
open-cell intersections and interior lattice points.

\subsection{A multiple-object stress family}
The corner-free property does not imply that a query cube meets only one
input cube. This can fail even for the quarter-offset solids of equal-side
integer cubes, not merely for arbitrary real objects. For $D\geq2$, put
$n=4$, $C=[1,2]^D$, and, for $\sigma\in\{0,1\}^{D-1}$, set
\[
 b_\sigma=(2\sigma_1,\ldots,2\sigma_{D-1},1),\quad L_\sigma=2,
 \qquad S_\sigma=\prod_{i=1}^D\{b_{\sigma i},b_{\sigma i}+1\}.
\]
The $2^{D-1}$ integer cubes are pairwise disjoint: distinct binary words
have a coordinate with the disjoint integer intervals $\{0,1\}$ and
$\{2,3\}$. Their solids have, in that coordinate, intervals
$[-1/4,5/4]$ and $[7/4,13/4]$, with a gap $1/2$.
In the last coordinate every solid has interval $[3/4,9/4]$.
Consequently every solid properly meets $C$, while no solid corner is
interior to $C$: every such corner has last coordinate $3/4$ or $9/4$.
Thus a universal proper-occupancy bound for corner-free query cubes in
this exact lattice class must be at least $2^{D-1}$. This neither proves
sharpness of $2^D$ nor asserts that this particular $C$ occurs as a leaf
of our root-based recursion. It rules out replacing the occupancy
argument by a one-object assumption. The finite unit139 baseline leaves
all had occupancy at most one, so those runs did not test this phenomenon.
\subsection{Constructing linearly many vertex-free rectangles}
Let $V$ be the set of all corners of the $Q_k$, so $N=|V|=2^DJ$.
For $J=0$ use the single root rectangle and skip this construction.
For $J>0$ begin with the cube $[-1,n]^D$. At a current cube $C$ keep
exactly $V_C=V\cap\operatorname{int}C$.
If $V_C$ is empty, stop at that cube. Otherwise consider its $2^D$ open
half-side children.

If the points are not all in one open child, perform all midpoint cuts
and recurse on the nonempty child point sets. Points on the cuts are
removed from subsequent interior sets. If all points lie in one open
child, reflect and translate coordinates so the common corner is zero
and $C=[0,s]^D$. Put
\[
 a=\max\{x_i:x\in V_C,\ 1\leq i\leq D\},\qquad C'=[0,a]^D.
\]
Then $0<a<s/2$. Successively cut at the $D$ upper facet planes of $C'$,
producing $C'$ and the $D$ boxes
\[
 E_i=[0,a]^{i-1}\times[a,s]\times[0,s]^{D-i},\qquad1\leq i\leq D.
\]
Use disjoint boundary ownership for these boxes. Only $C'$ can contain
points of $V$ in its interior, and at least one point of $V_C$ lies on
its boundary. Stop at the other boxes and recurse inside $C'$ if needed.
Every intermediate or terminal cell is full-dimensional; a boundary
ownership convention never creates a separate hyperplane cell.

Here is a direct finite termination and size proof. Among active cells,
let $n_C=|V_C|>0$ and use the nonnegative integer potential
\[
 \Psi=\sum_{C\ {\rm active}}(2n_C-1).
\]
Initially $\Psi=2N-1$. If a processed cell loses $r$ points to new cut
boundaries and has $t$ nonempty child point sets, the decrease is
$2r+t-1$. In the shrinking step $r\geq1$. In a midpoint step either
$r\geq1$ or $t\geq2$. Thus every event decreases $\Psi$ by at least one,
so there are at most $2N-1$ events. Each event creates at most $2^D-1$
additional leaves, since a shrinking step creates only $D$ extra leaves.
The number $M$ of terminal rectangles satisfies
\begin{equation}\label{eq:explicit-first-size}
 M\leq1+(2N-1)(2^D-1)
 \leq2^{D+1}(2^D-1)J.
\end{equation}
Tracking exactly $V\cap\operatorname{int}C$ ensures that every terminal
rectangle has no corner of an input cube in its interior. This argument
is a construction and a decreasing integer potential, not a supplied
partition certificate. It does not prove a fast implementation bound.

\subsection{A uniform occupancy bound for the terminal rectangles}
Terminal cubes meet at most $2^D$ input cubes properly, by the corner bound.
A terminal noncube has the form $E_i$ above after a coordinate reflection.
It is contained in the slab $[0,s]^{i-1}\times[a,s]\times[0,s]^{D-i}$.
Cover this slab by $2^{D-1}$ closed cubes of side $s-a$: in coordinate $i$
use $[a,s]$, and in every other coordinate choose either $[0,s-a]$ or
$[a,s]$. These intervals cover $[0,s]$ since $a<s/2$.
Every covering cube lies inside the parent cube and outside the interior
of $C'$, and its interior contains no point of $V$. The corner bound
therefore applies separately to each covering cube.

If an input cube properly meets $E_i$, the intersection contains an open
set. A point of that set avoiding the finitely many covering-facet planes
belongs to the interior of a covering cube. Thus all proper intersections
are counted by the cover, even if some particular point lies on a cover
boundary. Consequently each terminal rectangle properly meets at most
\begin{equation}\label{eq:explicit-occupancy}
 T_D=2^{D-1}2^D=2^{2D-1}
\end{equation}
input cubes. The covering argument, including its boundary qualification,
is proved here; no arbitrary-box unclutteredness is assumed.

\subsection{From open geometry to an actual integer partition}
For each terminal rectangle list the at most $T_D$ cubes meeting its
interior. Refine by all their facet coordinates. There are at most $2T_D$
cut coordinates per axis, so at most $(2T_D+1)^D$ refined rectangles per
terminal cell. All cuts are axis-parallel. Now assign every splitting
hyperplane to one child consistently, intersect with $X$, and delete empty
integer sets. The result is a disjoint cover of $X$ by integer rectangles.
Taking integer points and deleting empty sets does not increase the count.
More explicitly, use the uniform convention $[l_i,h_i)$ in every axis
of every cell. A coordinate split at $t$ then assigns the plane to the
upper child. The two intervals $[l,t)$ and $[t,h)$ partition $[l,h)$
exactly, even when $t$ is an integer. By induction through every cut,
these half-open cells partition $[-1,n)^D$; intersecting with $X$ gives
\[
 R=\prod_i\bigl(\{\lceil l_i\rceil,\ldots,\lceil h_i\rceil-1\}
                 \cap\{0,\ldots,n-1\}\bigr).
\]
This formula is also valid for negative or integer endpoints. An empty
axis makes the entire integer cell empty; no lattice point is assigned
twice or lost when those cells are deleted. The outer upper face at $n$
contains no point of $X$.

To check membership rigorously, suppose an assigned integer point $x$
lies in $S_k$. It is interior to $Q_k$, while its assigned terminal
rectangle has $x$ in its closure. A small ball about $x$ lies in $Q_k$
and meets the interior of that full-dimensional terminal rectangle.
Therefore $Q_k$ was in its list. None of the new facet coordinates is an
integer. Hence each refined integer rectangle is entirely inside a
listed $S_k$ or disjoint from every listed $S_k$. An unlisted cluster
cannot contain one of its integer points, by the same closure argument.
Strict disjointness then implies that each integer rectangle is inside
one cluster or outside every cluster.

In particular, with the explicit dimension-only constant
\begin{equation}\label{eq:explicit-final-size}
 A_D=2^{D+1}(2^D-1)(2^{2D}+1)^D,
 \qquad |\mathcal P|\leq A_D\max\{1,J\},
\end{equation}
we have derived the required partition. The constant is deliberately
loose and no sharpness or algorithm running-time claim is made.
For $n=1$ and for singleton input cubes all constructions remain valid.

\subsection{The actual coefficient and fixed-penalty consequence}
Suppose, in addition, that for fixed $c>0$ and $\beta\geq1$ the actual
positive scalar starting coefficients obey
\[
 \begin{gathered}
 a(x)\geq cL_k(L_k-1)/4\quad(x\in S_k),\\
 a(x)\geq cn(n-1)/4\quad(x\notin\cup_k S_k),\\
 a(x)\leq\beta L_k^2\quad(x\in S_k).
 \end{gathered}
\]
These coefficient/cube assumptions are retained, not inferred from the
critical moment. On an interval of $l$ integer vertices, pairwise variance,
path Cauchy--Schwarz and at most $l^2/4$ pairs through one cut give
\[
 \sum|u-\bar u|^2\leq\frac{l(l-1)}4\sum_{\rm edges}|du|^2.
\]
For completeness, product variance tensorization follows from the commuting
coordinate averages $P_i u=l_i^{-1}\sum_{t=0}^{l_i-1}u(x_1,\ldots,t,\ldots,x_D)$: the telescoping orthogonal decomposition of
$I-\prod_iP_i$ gives a sum bounded by
$\sum_i\|u-P_i u\|_2^2$. Applying the interval bound on its fibers proves
\[
 \operatorname{Var}_R(u)\leq\frac{L(R)(L(R)-1)}4
                      \sum_{e\subset R}|du(e)|^2.
\]
Inside a cluster $L(R)\leq L_k$, and outside all clusters $L(R)\leq n$.
The displayed coefficient lower bounds therefore give
$c\operatorname{Var}_R(u)\leq\mathcal E_{a,R}(u)$.
Singleton residuals are zero, so no zero lower bound is divided by.
Sum over all field/internal coordinates and all partition cells. Original
internal positive-axis edge occurrences are disjoint; omitted inter-cell
and periodic-wrap occurrences have nonnegative energy. Thus the ACTUAL
normalized means
\[
 (C_{\mathcal P}U)_{R,j}=\sqrt{|R|}\,\operatorname{mean}_R U_j
\]
satisfy
\[
 c\|U\|^2\leq\mathcal E_a(U)+c\|C_{\mathcal P}U\|^2
\]
for every cochain, with output dimension $h|\mathcal P|$,
$h=D(N_c^2-1)$. This uses an all-cochain variance estimate, not the invalid
inference from a kernel-only margin to fixed penalty $c$. Means need not
preserve the divergence constraint. The inequality consequently also
holds on that constraint subspace without changing $C_{\mathcal P}$.

If $p=D/2$, $\mu\geq1$, and the original critical mass bound is
$\sum_Xa^{-p}\leq Mn^D/\mu^p$, then each supplied cube contributes
at least $L_k^D/(\beta L_k^2)^p=\beta^{-p}$ to this mass. Hence
$J\leq\beta^pMn^D/\mu^p$, and \eqref{eq:explicit-final-size} gives
\[
 \dim F\leq hA_D(1+2^p\beta^pM)
                   \max\{1,n^D/(1+\mu)^p\}.
\]
This proves the restricted supplied-cube consequence with an explicit
constant and without a fat-object BSP theorem as a black box. It does
not construct the supplied cubes from the critical moment, remove either
scale bound or the strong complement bound, resolve the unrestricted
critical rectangular question, or establish a covariant/continuum gap.
No new originality or independent-discovery claim is made: the compressed
splitting strategy is attributed classical geometry, here derived and
adapted completely for this special class. The frozen reviewed PDF and
its full reports and pending permission remain unchanged.

\section{Physical cycle obstructions for an actual observation map}
\label{sec:physical-cycle}
This part closes a limitation of the scalar four-cycle example: the
obstruction is now realized inside the actual divergence-zero flat physical
cochain domain. The kernel estimate below is proved for all actual cochains,
not just the special fields used as negative witnesses. The conclusion
concerns this observation and this fixed penalty. It does not obstruct all
partitions or the optimal-rank question from the preceding parts.

\subsection{The actual form and rectangular observations}
Fix $D\geq3$, $N_c\geq2$ and $X=(\mathbb Z/4\mathbb Z)^D$.
Write $V=\mathbb R^{N_c^2-1}$ with its standard Euclidean norm and
$A=(A_\ell(x))_{x\in X,1\leq\ell\leq D}$ for an actual one-cochain.
All norms use counting measure. Choose the actual total coefficients
\[
 \alpha=(1,6,4,1),\qquad a(x)=\alpha_{x_1},\qquad
 w(i,x,\ell)=\alpha_{x_1}-1
\]
for every derivative direction $i$ and physical component $\ell$.
Define actual forward and backward differences by
\[
 D_i^+f(x)=f(x+e_i)-f(x),\qquad D_i^-f(x)=f(x)-f(x-e_i).
\]
The flat curl and backward divergence, with one curl for each unordered
pair, are
\[
 (\operatorname{curl}A)_{ij}=D_i^+A_j-D_j^+A_i\quad(i<j),\qquad
 \operatorname{div}^-A=\sum_iD_i^-A_i.
\]
Their counting quadratic form and the original componentwise repair are
\begin{align*}
 H_0(A)&=\sum_{i<j,x}\|D_i^+A_j(x)-D_j^+A_i(x)\|^2
       +\sum_x\Big\|\sum_iD_i^-A_i(x)\Big\|^2,\\
 R_w(A)&=\sum_{i,x,\ell}w(i,x,\ell)\|D_i^+A_\ell(x)\|^2.
\end{align*}
In operator notation the repair is $\langle A,G^*WG A\rangle$ for the
actual stacked gradient $G$ and diagonal multiplication $W$.
These formulas fix the energy on paper; no comparison certificate is supplied.

\paragraph{Deriving the flat identity.}
Periodic reindexing gives
\[
 \sum_x\langle D_i^+f,h\rangle=-\sum_x\langle f,D_i^-h\rangle.
\]
Coordinate shifts commute. Applying this formula twice and real symmetry
therefore yields
\[
 \sum_x\langle D_i^+f,D_j^+g\rangle
 =\sum_x\langle D_j^-f,D_i^-g\rangle.
\]
Also $\sum_x\|D_i^-f\|^2=\sum_x\|D_i^+f\|^2$ by translation.
Expand each curl square with $f=A_j$, $g=A_i$: it supplies the two
off-diagonal gradient norms and the negative mixed terms. Expansion of
the divergence square supplies the diagonal norms and the opposite
mixed terms. Cancellation for each $i<j$ proves, for every cochain,
\[
 H_0(A)=\sum_{i,x,\ell}\|D_i^+A_\ell(x)\|^2.
\]
No divergence condition or lift ansatz was used. Define $H=H_0+R_w$.
Adding the displayed repair, with $1+w(i,x,\ell)=\alpha_{x_1}$, gives
\begin{equation}\label{eq:actual-cycle-energy}
 H(A)=\sum_{x\in X}\sum_{i=1}^D\sum_{\ell=1}^D
 \alpha_{x_1}\|A_\ell(x+e_i)-A_\ell(x)\|^2.
\end{equation}
The cancellation was unweighted; the varying repair is added afterwards.
We do not move $\alpha$ through summation by parts or assert a weighted
curl--divergence identity. This is the paper counterpart of
\texttt{finite\_flat\_curl\_div\_energy} and
\texttt{componentwiseWeightedRepair\_energy} in the pinned formal source.
Thus the coefficient is the original shared scalar total coefficient,
not a substituted energy or an interacting Hessian.

Put $I_0=\{0,1\}$ and $I_1=\{2,3\}$. For every transverse
$y\in(\mathbb Z/4\mathbb Z)^{D-1}$ the blocks $I_b\times\{y\}$
are nonempty axis-aligned rectangles. They are disjoint and cover $X$.
For each physical component $\ell$ and each standard internal coordinate $k$,
define the actual normalized means
\begin{equation}\label{eq:actual-cycle-C}
 (CA)_{b,y,\ell,k}=\frac{1}{\sqrt2}\sum_{t\in I_b} A_\ell(t,y)_k.
\end{equation}
The output includes \emph{all} these coordinates. If
$h=D(N_c^2-1)$, its dimension is $2h4^{D-1}$; no optimality is asserted.

The spatial mean of $a$ is $\mu=3$, and $1\leq a\leq2\mu$.
For $p=D/2$ its original normalized inverse moment is
\begin{equation}\label{eq:actual-cycle-moment}
 \frac1{|X|}\sum_{x\in X}\left(\frac{\mu}{a(x)}\right)^p
   =\frac{2\cdot3^p+(1/2)^p+(3/4)^p}{4}=:M_D<\infty.
\end{equation}
This is an actual member of the stated bounded-moment coefficient class
when the allowed bound is at least $M_D$. Its fixed finite size and the
large output dimension do not disprove an asymptotic optimal-rank order.

\subsection{A fixed-penalty counterexample with a proved kernel premise}
For a real cyclic vector $u=(u_0,u_1,u_2,u_3)$ set
\begin{align*}
 E(u)&=(u_1-u_0)^2+6(u_2-u_1)^2+4(u_3-u_2)^2+(u_0-u_3)^2,\\
 |u|^2&=\sum_{t=0}^3u_t^2,\qquad
 O(u)=\frac{(u_0+u_1)^2+(u_2+u_3)^2}{2}.
\end{align*}
Choose a unit internal coordinate vector $T$. The actual lift is
\[
 (Ju)_2(x)=u_{x_1}T,\qquad (Ju)_\ell(x)=0\quad(\ell\ne2).
\]
Its backward divergence is zero: its only nonzero component is constant
in coordinate $2$. It is injective and direct summation gives
\begin{equation}\label{eq:cycle-transports}
 \|Ju\|^2=4^{D-1}|u|^2,\quad H(Ju)=4^{D-1}E(u),\quad
 \|CJu\|^2=4^{D-1}O(u).
\end{equation}
The last equality uses the entire output in \eqref{eq:actual-cycle-C},
not only a selected nonzero coordinate.

\paragraph{Kernel premise on all actual fields.}
For $v=(x,-x,y,-y)$, direct algebra yields
\begin{align}
 E(v)&=11x^2+14xy+23y^2,\qquad |v|^2=2x^2+2y^2,\notag\\
 E(v)-\frac{19}{5}|v|^2
   &=\frac{17}{5}\left(x+\frac{35}{17}y\right)^2
       +\frac{84}{85}y^2\geq0.\label{eq:cycle-kernel-square}
\end{align}
If $CA=0$, both pair sums vanish for every actual scalar profile
$t\mapsto A_\ell(t,y)_k$. Apply \eqref{eq:cycle-kernel-square}
to each such profile and sum over all $y,\ell,k$. These profiles account
for the full norm of $A$. Their energies are exactly the $i=1$ term
of \eqref{eq:actual-cycle-energy}; every other derivative term is
nonnegative because $\alpha_t\geq1$. Therefore
\begin{equation}\label{eq:full-physical-kernel}
 CA=0\quad\Longrightarrow\quad H(A)\geq\frac{19}{5}\|A\|^2
 \qquad\text{for every actual cochain }A.
\end{equation}
No divergence condition, scalar-lift ansatz or supplied coercivity
certificate was used in this proof.

\paragraph{Failure on an actual divergence-zero field.}
For $u=(19,-9,-5,-5)/14$ one has
\[
 E(u)=\frac{52}{7},\qquad |u|^2=\frac{123}{49},\qquad
 O(u)=\frac{25}{49}.
\]
Consequently the actual divergence-zero field $Ju$ satisfies
\begin{equation}\label{eq:physical-negative-defect}
 H(Ju)+\frac{19}{5}\|CJu\|^2-\frac{19}{5}\|Ju\|^2
       =-\frac{6}{35}4^{D-1}<0.
\end{equation}
Equations \eqref{eq:full-physical-kernel} and
\eqref{eq:physical-negative-defect} prove the complete fixed-map statement:
kernel coercivity does not imply the same coefficient as a fixed coarse
penalty, even on the actual flat physical divergence-zero domain.
They do not rule out larger penalties or different observations.

\subsection{The exact finite penalty at the subsharp target}\label{sec:cycle-subsharp}
The failure just proved concerns $K=19/5$. Here we determine exactly
which larger penalties work for this same observation and target.
This subsection is a paper proof with an exact rational coefficient
check; it is not part of the paired Lean theorem below.

Every $u\in\mathbb R^4$ has the unique decomposition
\[
 u=m(1,1,1,1)+(x,-x,y,-y)
       +(\delta/2,\delta/2,-\delta/2,-\delta/2),
\]
where $m=(\sum_tu_t)/4$, $x=(u_0-u_1)/2$, $y=(u_2-u_3)/2$
and $\delta=(u_0+u_1-u_2-u_3)/2$. Direct expansion gives
\begin{align*}
 |u|^2&=4m^2+2x^2+2y^2+\delta^2,\qquad O(u)=4m^2+\delta^2,\\
 E(u)&=11x^2+14xy+23y^2+7\delta^2-10\delta(x+y).
\end{align*}
Set $X=x-25\delta/2$ and $Y=y+75\delta/14$. For every real $K$,
\begin{align}
 E(u)+KO(u)-\frac{19}{5}|u|^2
 &=4\left(K-\frac{19}{5}\right)m^2
   +\frac{17}{5}\left(X+\frac{35}{17}Y\right)^2\notag\\
 &\quad+\frac{84}{85}Y^2
   +\left(K-\frac{1138}{35}\right)\delta^2.
 \label{eq:exact-subsharp-square}
\end{align}
Since $1138/35>19/5$, this proves sufficiency of $K\geq1138/35$
on every scalar profile. Necessity follows from
\[
 u_*=(13,-12,-41/7,34/7),\qquad
 O(u_*)=1,\quad |u_*|^2=\frac{18174}{49},\quad
 E(u_*)=\frac{67468}{49}.
\]
It has $m=0$, $\delta=1$ and $X=Y=0$, hence defect $K-1138/35$.

Apply the scalar inequality to every profile
$t\mapsto A_\ell(t,y)_k$ and sum over all $y,\ell,k$.
Their squared norms and observations exhaust the actual full norms,
and their energies exhaust the first-direction contribution in
\eqref{eq:actual-cycle-energy}. The remaining directions are
nonnegative. Conversely $Ju_*$ is actual divergence zero and multiplies
the scalar defect by $4^{D-1}$. Therefore the following are equivalent:
\begin{enumerate}
 \item $K\geq1138/35$;
 \item $H(A)+K\|CA\|^2\geq(19/5)\|A\|^2$ for every actual cochain;
 \item the same inequality for every actual divergence-zero cochain.
\end{enumerate}
This is a necessity-and-sufficiency statement for the fixed map, not
an obstruction to different partitions. In particular, a larger finite
penalty repairs the subsharp target even though $K=19/5$ fails.
The sharp kernel target in the next subsection behaves differently.

A second calculation records the classical block elimination behind the
constant. The kernel residual matrix and mixed vector are
\[
 M=\begin{pmatrix}17/5&7\\7&77/5\end{pmatrix},\quad
 \det M=84/25,\quad b=(5,5),\quad
 M^{-1}b=(25/2,-75/14).
\]
Thus $b^TM^{-1}b=250/7$ and the residual coefficient is
$K+16/5-250/7=K-1138/35$, agreeing with
\eqref{eq:exact-subsharp-square}. No new Schur-complement principle is claimed.

The unit151 exact verifier constructs $E$, $O$ and $|u|^2$ from the actual
four-coordinate linear forms and compares every coefficient of the
polynomial in $(m,x,y,\delta,K)$, rather than sampling values. Normal and
optimized Python in one own Colab Linux runtime each accepted five checks
and rejected three real coefficient mutations; both exits were zero and
their JSON outputs were identical. The 8794-byte archive has ten members
and nine checked inner hashes. It and the full hand derivation are
included in the reproduction package. The driver took about 0.56 seconds;
the runtime was released after recovery and exact billed consumption is
unknown. This check is neither a new Lean theorem nor independently
authored verification. Its source is
\begin{quote}\small\ttfamily
564b84230478cb8d40505cb67985470e9bae699d.
\end{quote}
\subsection{Sharp free-penalty extension and its distinct status}\label{sec:cycle-sharp}
The following is an analytic continuation, outside the paired Lean result.
Put $s=\sqrt{85}$ and $c_*=(17-s)/2$. Since $9<s<10$, the exact identity
\[
 E(v)-c_*|v|^2=(s-6)\left(x+\frac{7y}{s-6}\right)^2
\]
holds on the same block kernel. The nonzero vector
$v_*=(7,-7,6-s,s-6)$ attains equality. The all-profile argument above
also establishes the sharp kernel constant for the actual form in paper
proof, and the lift establishes sharpness.

For $p_\delta=(\delta/2,\delta/2,-\delta/2,-\delta/2)$, the two vectors
are orthogonal, $Cv_*=0$, and $O(p_\delta)=|p_\delta|^2=\delta^2$.
For any finite $K\geq0$, direct expansion gives
\begin{align*}
 &E(v_*+p_\delta)+K O(v_*+p_\delta)-c_*|v_*+p_\delta|^2\\
 &\hspace{18mm}=(7+K-c_*)\delta^2-10(13-s)\delta.
\end{align*}
The positive number $q_K=7+K-c_*$ permits the choice
$\delta=5(13-s)/q_K$, giving the strictly negative value
$-25(13-s)^2/q_K$. Lifting multiplies this complete defect by
$4^{D-1}$ and preserves actual divergence zero. Thus no finite quadratic
penalty repairs this same observation at the sharp kernel endpoint.
A penalty $\Phi(z)\leq K\|z\|^2$ only near zero is also defeated:
scale the negative witness by a sufficiently small nonzero real number.
Both radius conditions can be met, and its defect remains a negative
multiple of the square of that scale. This is not a strict-margin claim.

These exact analytic identities and arguments are supplied here, but their
complete physical sharp-endpoint/free-penalty statement has \emph{not}
been verified in Lean in the paired execution described next. Finite
checks in dimensions three and four, including $K=0,1,100$, are diagnostics
only, not the universal theorem or independent review.

\subsection{Formal scope and the paired reproduction record}
The fixed-penalty result, including the actual full-kernel premise and
all transport/comparison identities used to connect it, was compiled in
two separately allocated Colab Linux CPU sessions with high RAM. Each
started with a fresh clone at the same immutable source
\begin{quote}\small\ttfamily
0c706827dcbb11859e56bb6c7dc981eb37a02e69.
\end{quote}
The final theorem is named
\texttt{physicalCycleFour\_kernel\_fixed\_penalty\_counterexample}
in \texttt{PhysicalCycleFullForm.lean}. It asserts the proved
all-cochain kernel bound together with failure of the fixed-penalty
inequality on the actual divergence-zero domain. Four preceding printed
endpoints identify longitudinal energy, the full actual energy sum,
its comparison, and the kernel estimate. They do not assume transport
certificates or posit a structure containing the conclusion.

All five actual dependency readouts in each clone contain only
\texttt{propext}, \texttt{Classical.choice} and \texttt{Quot.sound}.
The normal and optimized validator in each clone accepted all five
readouts and rejected three actual oracle mutations. All actual decimal
exit sentinels were zero. Lean was version 4.29.0-rc6 with Mathlib pinned
in the included manifest. The recovered theorem source, oracle source,
compiled binary, oracle output, validator source, normal and optimized
outputs and driver have identical hashes across the two fresh executions.
The raw per-run reports still say first-clone pass and paired false:
the identical driver executes one clone per session. They remain
unmodified; the separate aggregate records the completed pair.

The first successful driver took about 374.31 seconds and the second
400.28 seconds. These are not billed connection times. Both environments
were disconnected after complete evidence recovery. No Lean compilation,
oracle or numerical campaign ran on Windows. Separate fresh executions
are not independently authored proofs or independent scientific review.
The previous failed attempt, its exact source and full compiler log are
preserved: mixed finite-index inference and nested sum type inference
were corrected without changing the theorem statements.

The package includes both public evidence archives. Their comparison is
recorded in the aggregate file:
\begin{quote}\small\ttfamily
paired-full-form148.json
\end{quote}
Full hashes and source bindings are in \texttt{inputs.json}.
Each archive has 45 checked inner manifest
members and excludes third-party installers and downloaded binary assets.
Their complete originals, including those assets, are retained in the
private evidence record. The compiled binary hash begins
\texttt{b0f8ff8654d79d02}, and the full value is in the paired comparison file.

\paragraph{Attribution and scope of novelty.}
The scalar square identities and negative vectors predate this physical
embedding in the dossier. The present consolidation adds the explicit
actual-domain embedding and the formally derived all-field comparison;
it claims no new Schur-complement principle, priority or exhaustive novelty
finding. Carlson, Haynsworth and Markham's classical paper,
\emph{A Generalization of the Schur Complement by Means of the Moore--Penrose
Inverse}, SIAM Journal on Applied Mathematics 26 (1974), 169--175,
\href{https://doi.org/10.1137/0126013}{doi:10.1137/0126013}, supplies algebraic
context. The publisher abstract and bibliography were inspected in the
source investigation; the full paywalled article was not read. The direct
finite proof here carries the mathematical claim.

\section{Evidence provenance and verification limits}\label{ap:sec:evidence}
The analytic sources, their exact revisions and the executable finite diagnostics are
collected in a hash-allowlisted companion package. Its map gives full SHAs and SHA256
values; short forms below are identifiers, not substitute hashes.
\begin{center}\small
\begin{tabular}{lll}
Item & Source revision & Verification status\\\hline
Actual flat Hodge identity & \texttt{dfd49d0b0730} & earlier formal source\\
Critical spectral theorem & \texttt{527fcdcd119a} & analytic adapters\\
Dyadic spatial separation & \texttt{fb65a67af360} & analytic adapters\\
Higher moment block proof & \texttt{6dcb893155e5} & analytic adapters\\
Exact finite fixture code & \texttt{4bb460e1537e} & own Colab109\\
Full diagnostic report & \texttt{602ae3f0b59a} & finite rational batch
\end{tabular}
\end{center}
Colab109 checked four actual profiles on $(\mathbb Z/2\mathbb Z)^4$ using rational
arithmetic. Each normal and optimized Python mode accepted four profiles and rejected
16 real matrix, inverse or flow mutations, with matching outputs and exit code zero.
It checked the constant mode, forward symbols, actual pseudoinverse identities,
feasible/minimizing flow energy, comparison inertia, strict low-spectrum multiplicities
and noncommuting counts. There are 12 distinct density examples, evaluated again for
each profile: 48 evaluations per mode, not 48 independent examples.

All eight files in the diagnostic archive's inner manifest were rehashed. Runtime
execution was recorded at 17:13:06.300923--17:13:16.201862 UTC on 2 November2026,
with automatic runtime release and zero sessions verified afterwards. Rounded resource
units do not prove zero cost. No local Windows mathematical execution occurred.
This single embedded-SHA finite batch is not two fresh matching clones, a Lean proof,
a theorem about arbitrary dimensions or odd lattices, or independent mathematical review.
This manuscript assembly adds no new Lean execution. Operational repository checks
on unchanged Lean files do not verify these new analytic proofs.

There is also subsequent, strictly focal formal evidence. Source
\texttt{41f8492c6190} supplies the ten four-cycle theorems and their oracle.
Own Colab118 compiled that same source in two fresh clones; both actual oracle logs
and compiled module hashes agree. Each of four normal/optimized validator runs
accepted ten readouts and rejected three real corruptions. All readouts use only
\texttt{propext}, \texttt{Classical.choice}, and \texttt{Quot.sound}.
The full archive has 69 manifest files and is included with the failed Colab116/117
compiler attempts. Source hashes, complete logs and exact reports are allowlisted.
This verifies Proposition~\ref{ap:prop:formal-cycle}'s finite identities, not the
geometric, shorted-operator, free-penalty or general rectangular adapters.
Own Colab115 separately checked 18 finite diagnostics and eight real mutations in
each mode of two fresh clones. Finite examples do not establish their general claims.
Section~\ref{ap:sec:spatial-new} consolidates analytic sources at
\texttt{0fa0d91000f6}, \texttt{312f476ea6e5}, \texttt{f52ae344f61f},
\texttt{1850be2aa6d5} and \texttt{b348f4488b3a}. The original spatial proofs and the self-contained Sobolev and cube
arguments remain analytic. The physical fixed-penalty endpoints have the
separately pinned paired formal record in Section~\ref{sec:physical-cycle}.
The exact subsharp threshold has an exact rational coefficient check, not a
Lean certificate; the sharp extension is analytic. Rendering does not upgrade
any verification status. Earlier exact PDFs and assessments remain preserved.

\section{Model assessments and review context}\label{sec:model-context}
This is a new exact-artifact revision candidate. The published parent v6
has PDF SHA256
\begin{center}\scriptsize\texttt{ccf53e0afc9819c3fc9b2008de9076a60f60d23ef6ac2c4ccd82adca30cbf389}.\end{center}
Its historical exact-PDF AI report154 records a 7.6/10 minor-revision result
in the displayed Gemini 4 Argon Alto interface, with no backend attestation
or established reviewer independence. That report remains bound only to
that 46-page PDF and is not an assessment of this revised manuscript.
No score, independent review, human peer review or editorial acceptance
is inherited. The new matrix continuation was prepared in the continuing
Codex session; exact backend identity is not exposed. Its complete
source/reproduction and review context are retained separately, and the
full revised PDF requires its own exact-hash review.

The prior exact v5 received a preparing-agent self-assessment of 6/10 on a
separate rubric, with eight retained objections. Its later Google Gemini report
used the AIRR assessment scale and scored 3.4/10 with a major-revision
recommendation. These scales are not pooled or converted. Neither number rates
this changed candidate. The full prior report, request, findings, exact-artifact
hash and operator response are preserved in the
\href{https://airr.science/papers/ARR-2026-7EPM3V3DP796BBRR/versions/v5/\#model-assessments}{v5 model-assessment record}.

The interface displayed Gemini 4 Argon with High reasoning. The operator
observed that display; no backend API attestation was available. The request
supplied the exact v5 PDF and withheld earlier reports and scores. Prior
participation, technical history isolation and memory isolation remain unknown.
The model reported a PDF-only review without external browsing, code execution
or inspection of the repository and Lean/Colab evidence. Its asserted proof
checks are model statements, not an operator correctness certificate. This is
neither established independent review nor human peer review.

The sole listed material objection alleged a missing factor $c$ in v5
Equation~\eqref{ap:eq:scaled-assumptions}. Exact TeX, PDF text and the hash-matched
rendered page show that $c$ was already present. A separate preparing-agent
response records this contradiction without editing the report or its grade.
The claim that Proposition~\ref{prop:volume-average} omitted the endpoint $R=1$
also contradicts the exact artifact. Other reported defects were real: this
candidate corrects stale weight symbols $m$ to $\mu$, adds a stable reference to
Section~\ref{ap:sec:objects}, and corrects an additional stale weight symbol in
the physical line-test energy. The original v5 PDF and all its reports remain
unchanged. No repeat review of unchanged v5 was requested to improve its grade.

\paragraph{Reports on distinct frozen successors.}
The 33-page revision136, SHA256 prefix \texttt{eb4df3}, received its
first Google Gemini report of 3.5/10 and minor revision. The operator
observed Gemini 4 Argon High; the raw backend identifier is unknown.
The review was non-blinded, with historical assessment context in the PDF.
Its five retained objections concern uninspected formal/Colab evidence,
non-Lean analytic continuations, primary sources not inspected by that
reviewer, the open unrestricted rectangular question, and the absent
interacting or continuum connection. The full report, request and context
are included separately as \texttt{review136-*} in the reproducibility package.

The distinct 14-page companion, SHA256 prefix \texttt{5c2774}, received
one first external Gemini report of 6.5/10 and minor revision. The selected
interface displayed Gemini 4 Argon High; the served backend is unexposed
and the report records unknown. No target grade or regeneration was used.
Historical context was visible in that PDF, so the review was non-blinded;
independence and context isolation are not established. All five objection
groups remain in the complete \texttt{review150-*} files: incomplete energy
definitions and uninspected execution archives; the local-to-global and
supplied-cube gap; distinct fixed, subsharp and sharp penalties; alleged
technical defects; and uninspected primary sources and review context.
Four technical allegations were contradicted by that exact artifact:
$a\geq s$ and $a,b\geq0$ were already printed, $2\,3^p$ was a product,
and $O(p_\delta)=|p_\delta|^2=\delta^2$ was complete across a line break.
The original report and grade, including inconsistent score subtotals,
remain unchanged. A separate operator disposition retains the evidence.

This successor supplies the missing energy definitions, cancellation and
exact finite subsharp threshold. Those additions do not make either model
inspect the archives, primary sources or this changed PDF. The 16-page
companion revision152 was an unreviewed consolidation, with no inherited
score. None of these historical scores rates the present combined artifact.
Model evaluations, public working-paper deposits, editorial acceptance and
human peer review remain distinct; no new acceptance or peer review is claimed.
This candidate is prepared by the same continuing Codex programme agent.
The served backend identity and effort for preparation are unexposed; no
technical isolation or independence is claimed. No numerical score or completed external exact-artifact assessment is assigned
to this changed candidate in the manuscript. The exact resulting PDF
must receive a separately retained first assessment after every page is inspected;
no earlier score or resolved-objection status transfers automatically.

The principal limitations and objections are these. The new finite Fourier, inverse,
field and partition adapters are analytic rather than kernel-verified. The classical
Rumin--Frank method and elementary refinement geometry limit any novelty claim; a
complete priority comparison is unfinished. The sufficient theorem uses one shared
scalar coefficient profile, and its distribution condition is not proved from arithmetic
means. Spectral coordinates are global; unrestricted critical rectangular means remain
unresolved. Exact finite diagnostics cover a small selected family, and some count bounds
are loose. The constants are safe rather than sharp. This manuscript makes no statement
about an interacting or continuum physical mass gap. The remaining programme obligations are not reduced by this manuscript; no continuum closure is claimed.

\section{The single remaining bounded question}
For fixed $D\geq3,c>0,N_c\geq2,M$, do all actual scalar profiles
$1\leq a\leq2\mu$, exact mean $\mu$, satisfying \eqref{ap:eq:moment}, admit
\eqref{ap:eq:rank} using normalized means on an unrestricted axis-aligned integer
rectangular partition? The mechanism must leave full cube trees, for example through
anisotropic cuts with a proved variance bound. A falsifier must defeat every such
rectangular partition, not merely every dyadic cube tree. After three substantive units
without a proved reduction or useful falsifier, the route must be changed or its concrete
blocker identified. The common-refinement route remains stopped at its recorded three-unit limit.
Formalization and consolidation do not reset that counter or create a second
active question. PR29 remains 20/41 with TermSource zero.


\clearpage
\section{Matrix gradients and critical rectangular penalties}\label{sec:matrix-continuation}
This grouped continuation retains its actual finite conventions. Its mathematical dossier is pinned at \texttt{66580c59}; the earlier sections and their verification referents are preserved. Its scalar rectangular optimum is not asserted to be the full restricted coexact optimum.
\subsection{Actual finite objects and observation scope}
Let $D\geq3$, $n\geq2$, $h\geq1$, $V=n^D$, $r=Dh$, and $p=D/2$.
On $X=(\mathbb Z/n\mathbb Z)^D$ use the counting inner product. For
$u:X\to\R^h$, define
\[
 du(x)=(u(x+e_i)-u(x))_{i=1}^D,\qquad
 Q_A(u)=\sum_x du(x)^T A(x)du(x),\quad L_A=d^*A d,
\]
where $A(x)$ is a real symmetric $r\times r$ matrix. A channel is one local
coefficient eigenvector, counted with multiplicity. Distinct matrices need
not commute. All derivatives have unit spacing. A physical derivative $nd$
requires an explicit rescaling of every energy and threshold.

The number $\Nlow(L_A)$ counts all eigenvalues strictly below $c>0$,
including negative and zero ones, with ordinary real multiplicity. Spectral
coordinates on this low space give a real observation $C$ and penalty
\begin{equation}\label{mx:eq:penalty}
 \Phi(y)=\sum_{\lambda<c}(c-\lambda)y_\lambda^2,
 \qquad Q_A(u)+\Phi(Cu)\geq c\|u\|^2.
\end{equation}
This is an existence result for an actual, generally nonlocal observation.
It does not construct the unresolved local rectangular observation.
For signed forms, $c\|y\|^2$ alone need not be sufficient.

\subsection{The positive matrix estimate}
Put
\[
 \alpha_D=\frac D2\left(\frac D{D-2}\right)^{D-1}
                 \left(\frac34\right)^D.
\]
\begin{matrixTheorem}\label{mx:thm:positive}
If every $A(x)$ is positive definite, then
\begin{equation}\label{mx:eq:positive}
 \Nlow(L_A)\leq h+\alpha_D c^p\sum_x\tr A(x)^{-p}.
\end{equation}
\end{matrixTheorem}
\begin{proof}
On the mean-zero space $H_0$, write $L_0=d^*d$ and $J=dL_0^{-1}$.
For any $b\in H_0$, the flow $j_*=A dL_A^{-1}b$ has $d^*j_*=b$.
Every other feasible flow $j=j_*+k$ has $d^*k=0$, and expansion gives
\[
 \langle j,A^{-1}j\rangle
 =\langle b,L_A^{-1}b\rangle+\langle k,A^{-1}k\rangle.
\]
Taking the actual feasible flow $Jb$ proves
$L_A^{-1}\leq J^*A^{-1}J$. No coefficients commute in this argument.

The unitary counting Fourier transform on $r$ components has
$L_1(\ell_2^r)\to L_\infty(\ell_2^r)$ norm $V^{-1/2}$. The symbol of $J$
is $z(k)/\lambda(k)$ with
$z_i(k)=e^{2\pi i k_i/n}-1$ and
$\lambda(k)=4\sum_i\sin^2(\pi k_i/n)$; set $J(0)=0$.
It is rectangular $r\times h$. Embed $\R^h$ isometrically in $\R^r$ by $I$
and use the square symbol $G(k)=J(k)I^*$. Thus $\|G(k)\|=\lambda(k)^{-1/2}$
for $k\ne0$. Writing $\mathcal F_r$ for the counting transform, the actual
factorization is $J=\mathcal F_r^*M_G\mathcal F_r I$.
Centered frequency representatives $b$ satisfy
$\lambda(k)\geq16|b|^2/n^2$. Counting the cube containing
$0<|b|<n/(4t)$ gives
\[
 \sup_{t>0}t^D\#\{k:\|G(k)\|>t\}\leq(3/4)^D V.
\]
When the radius is at most one the set is empty; otherwise
$(2\lfloor R\rfloor+1)^D\leq(3R)^D$ suffices.

Apply Frank's operator-valued Cwikel theorem~\cite[Theorem 3.1]{mx:Frank}
with $q=D$, $f=A^{-1/2}$, this square multiplier and counting measures.
Its constant is $\frac D2(\frac D{D-2})^{D-1}V^{-1}$, so the displayed
frequency $V$ cancels exactly once. It bounds the number of singular values
of $A^{-1/2}J$ strictly above $s$ by
$\alpha_Ds^{-D}\sum_x\tr A(x)^{-p}$. The right factor $\mathcal F_r I$ is an isometry
and cannot increase singular-value counts. Apply min-max to the inverse
comparison at threshold $1/c$ on $H_0$, then add the $h$ constant modes.
Complexification of a real symmetric matrix preserves multiplicities.
\end{proof}

\paragraph{Input and interpretation.}
The operator-valued estimate is classical; this proof fixes its finite
normalization, padding and actual form dictionary. It is not a new general
Cwikel theorem. A matrix arithmetic mean alone does not imply its inverse
moment hypothesis. For the formal flat cochain application one may set
$h=D(N_c^2-1)$, but an interacting Hessian or weighted curl/divergence form
must first be proved equal or comparable to this $Q_A$.

\subsection{Spectral repair and the capped budget}
\begin{matrixTheorem}\label{mx:thm:repair}
Let $A(x)$ now be arbitrary real symmetric matrices. In local orthonormal
eigenbases choose a set $S$ of channels containing all nonpositive ones.
Let $P_S$ be their orthogonal projection and $R(S)=\rank(P_Sd)$. Then
\begin{equation}\label{mx:eq:repair}
 \Nlow(L_A)\leq\min\left\{hV,\ h+R(S)+
          \alpha_Dc^p\sum_{j\notin S}\lambda_j^{-p}\right\}.
\end{equation}
Writing $\nu_0$ for the total multiplicity of nonpositive local eigenvalues,
one consequently has
\begin{equation}\label{mx:eq:capped}
 \Nlow(L_A)\leq\min\{hV,\ h+F_c(A)\},\quad
 F_c(A)=\nu_0+\sum_{\lambda_j>0}
                  \min\{1,\alpha_Dc^p\lambda_j^{-p}\}.
\end{equation}
\end{matrixTheorem}
\begin{proof}
Replace selected eigenvalues by $T$ larger than zero and all of them.
The resulting field $A_T$ is positive definite and
\[
 L_{A_T}-L_A=d^*P_S(T\id-A)P_Sd\geq0,
 \quad\rank(L_{A_T}-L_A)=R(S).
\]
The rank equality follows because $T\id-A$ is strictly positive on
$\im P_S$. Finite-rank min-max gives
$\Nlow(L_A)\leq\Nlow(L_{A_T})+R(S)$: intersect the low space of $L_A$
with the kernel of the update and the orthogonal complement of the low
space of $L_{A_T}$. A dimension larger than $R(S)$ would give a nonzero
vector whose identical quadratic forms are both strictly below and at least
$c\|u\|^2$. Theorem~\ref{mx:thm:positive} and $T\to\infty$ prove
\eqref{mx:eq:repair}. For \eqref{mx:eq:capped}, use $R(S)\leq|S|$ and select
all positive channels with $\alpha_Dc^p\lambda_j^{-p}\geq1$ as well.
\end{proof}

\paragraph{A strict extension of the full moment criterion.}
Let $m\geq1$ and $T_m=\max\{1,V/(1+m)^p\}$. If $F_c(A)\leq KT_m$,
then $\Nlow(L_A)\leq(h+K)T_m$. Every former SPD budget
\[
 \frac1{rV}\sum_x\tr((mA(x)^{-1})^p)\leq M
\]
implies this capped budget, with $K=2^p\alpha_Dc^prM$.
Conversely choose $m=n^2$, one actual diagonal edge channel
$\epsilon=n^{-4}$, and all others $m$. Now $T_m=1$ and
$F_c(A)\leq1+\alpha_Dc^pr$, whereas the normalized full inverse moment is
at least $n^{2D}/r$ and diverges. These bounds are derived for this explicit
family; no new moment bound is assumed.

\paragraph{Sparse contamination and its limitation.}
Selecting every channel with $\lambda<m$ gives
\[
 \Nlow(L_A)\leq h+R_m+\alpha_Dc^prVm^{-p},
 \qquad R_m=\rank(\boldsymbol1_{(-\infty,m)}(A)d).
\]
An actual selected-channel count $k_m\leq KT_m$ suffices since $R_m\leq k_m$.
If $A(x)\geq m\id$ outside a specified set $B$, then $k_m\leq r|B|$.
Neither condition follows from an arithmetic mean. This particular
prescribed-threshold count criterion does not subsume every full moment
family: $A=(m/2)\id$ has normalized inverse moment $2^p$ but every channel
lies below $m$. The two particular sufficient criteria are incomparable;
the general capped criterion \eqref{mx:eq:capped} does include the old one.

\paragraph{Actual rank and unavoidable costs.}
For diagonal direction/component channels the selected rows are graph
incidences. If $\kappa_a$ counts all components, including isolated vertices,
of the selected-edge graph for component $a$, then
$R(S)=\sum_a(V-\kappa_a)$: its kernel consists precisely of fields constant
on those components. Parallel rows at $n=2$ do not change that kernel.
This formula is not applied to rotated local eigenchannels.
Isolating $q$ nonadjacent vertices by zeroing their incident edges gives
at least $q+1$ kernel modes. Replacing these zeros by
$0<\epsilon<c/(2D)$ yields a $q$-dimensional strict low subspace of delta
functions. The defect-count cost cannot simply be deleted.

In \eqref{mx:eq:penalty}, a generic bound is $K\leq c+4Db$ with
$b=\max_x\|A_-(x)\|$, since $\|du\|^2\leq4D\|u\|^2$.
One negative channel $-B$ with all other channels $m$, at $n\geq3,h=1$,
has the trial vector $v=\delta_y-\delta_x$ on that edge and
\[
 \frac{Q_A(v)}{\|v\|^2}=-2B+(2D-1)m.
\]
Any normalized linear observation with penalty bounded by $K\|y\|^2$
therefore needs $K\geq c+2B-(2D-1)m$. Sparsity alone does not bound penalty.

\subsection{Stable signed families from actual resistance}
This section is independent of the Cwikel input. Its mechanism is classical
effective resistance and finite-rank congruence~\cite{mx:Zelazo,mx:Spielman}.
Let $e$ be one actual direction/component row and $v=d^*e$. With all
coefficients $m$ except this one set to $-B$, $B\geq0$,
\[
 H=mL_0-(m+B)vv^*,\qquad L_0=d^*d.
\]
On $H_0$ let $G=L_0^{-1}$. Translations and coordinate permutations make
every scalar row leverage equal. Their sum is
$\tr(dGd^*)=V-1$, so
\begin{equation}\label{mx:eq:rho}
 \rho=v^*Gv=\frac{V-1}{DV}.
\end{equation}
Components decouple; the same value holds for every $h$. The $DV$ rows
are counted individually even at $n=2$.

\begin{matrixProposition}\label{mx:prop:phase}
The exact threshold is $B_*=m(DV/(V-1)-1)$. On $H_0$, $H$ is positive
definite for $B<B_*$, has one zero direction at equality, and has exactly
one negative direction above it. If $m=n^2$, $B=\theta m$ and
$0\leq\theta<D-1$, its mean-zero gap is at least
\[
 g=16(D-1-\theta)/D>0.
\]
For every $0<c\leq g$, the normalized mean observation
$C(u)=V^{-1/2}\sum_xu(x)\in\R^h$ and $\Phi(y)=c\|y\|^2$ give coercivity
with optimal rank $h$, uniformly in $n$.
\end{matrixProposition}
\begin{proof}
Congruence by $G^{1/2}$ gives $m\id-(m+B)ww^*$ with $\|w\|^2=\rho$,
proving the phase claim. The equality witness is $Gv$.
Cauchy--Schwarz in the $L_0$ energy gives
$|v^*u|^2\leq\rho\langle u,L_0u\rangle$, hence
$H\geq[m-(m+B)\rho]L_0$ on $H_0$; equality holds on $Gv$.
For the stated scaling the coefficient is at least
$m(D-1-\theta)/D$. The actual first nonconstant Fourier eigenvalue of $L_0$
is $4\sin^2(\pi/n)\geq16/n^2$. This proves the gap. Decompose into
constants and $H_0$ to prove the displayed observation inequality.
The constant kernel forces rank at least $h$.
\end{proof}
For $\theta>0$, negative coefficient depth grows as $\theta n^2$, yet the
penalty is uniformly $c$. A bounded-depth hypothesis is sufficient for the
generic estimate, not necessary for every stable signed family. The actual
form is positive despite an indefinite coefficient field; this does not
contradict the fixed-$m$, $B\to\infty$ penalty obstruction above.

\subsection{Coupled channels and a nonlinear observation obstruction}
Choose $q$ actual rows of $d$ as $Z$, and $W=\diag(w_j)>0$. The modified
form is exactly $H=mL_0-Z^*WZ$. Define
\[
 \Gamma=ZGZ^*,\qquad \mathcal C=mW^{-1}-\Gamma,\quad N_0=h(V-1).
\]
Schur elimination of either positive diagonal block in
$\left[\begin{smallmatrix}mL_0&Z^*\\Z&W^{-1}\end{smallmatrix}\right]$
on $H_0\oplus\R^q$ proves
\begin{equation}\label{mx:eq:inertia}
 n_-(H)=n_-(\mathcal C),\quad n_0(H)=h+n_0(\mathcal C),\quad
 n_+(H)=N_0-q+n_+(\mathcal C).
\end{equation}
Indeed the two congruences give the identity
\[
 \operatorname{inertia}(H|_{H_0})+(q,0,0)
 = (N_0,0,0)+\operatorname{inertia}(\mathcal C/m).
\]
No disjoint-cycle or independent-row assumption occurs. In particular,
$H$ is positive semidefinite with exactly $h$ constant zeros if and only if
$\mathcal C$ is positive definite. General signed matrix criteria preexist
in~\cite{mx:Chen}; our $\Gamma$ uses the positive reference inverse.

\begin{matrixExample}\label{mx:ex:coupled}
Take $D=3,n=2,h=1,m=4$, and change the first two direction rows at one
vertex separately or together to coefficient $-84/11$. Then
\[
 \Gamma=\begin{pmatrix}7/24&5/48\\5/48&7/24\end{pmatrix},
 \quad w=128/11,\quad w/m=32/11.
\]
The diagonal follows from \eqref{mx:eq:rho}; the off-diagonal Walsh sum is
$(1/8)(1/2+1/3)=5/48$, with both forward orientations counted.
Each single channel is below its critical magnitude $68/7$ and has full
inertia $(7,0,1)$. Together the reduced matrix has diagonal $5/96$ and
off-diagonal $-5/48$, so its eigenvalues are $-5/96,5/32$.
Equation~\eqref{mx:eq:inertia} gives joint inertia $(6,1,1)$.

Each single form has mean-zero gap at least
$16(1-(32/11)(7/24))=80/33>1$. The same scalar mean observation with
$\Phi(y)=|y|^2$ repairs either single form at $c=1$. The joint form cannot
be repaired by any scalar observation $F:\R^8\to\R$ differentiable at zero
with $F(0)=0$ and $0\leq\Phi(z)\leq K|z|^2$ near zero. Its negative
eigenvector and constant vector span a two-dimensional nonpositive space.
The derivative of $F$ has a nonzero kernel vector $v$ there. Then
$F(tv)=o(t)$ and $\Phi(F(tv))=o(t^2)$, whereas $Q_A(tv)\leq0$, contradicting
coercivity at $c=1$ for small nonzero $t$.
\end{matrixExample}
This refutes the shortcut from individual to joint control, not the full
matrix resistance criterion or any theorem retaining a disjoint-cycle
hypothesis. The $q$-row perturbation also gives
$\Nlow(H)\leq h+q$ for $0<c\leq16m/n^2$. Extra zero or negative modes
provide the necessary rank lower bound in \eqref{mx:eq:inertia}; they do not
alone bound the penalty coefficient.

\subsection{Critical constants for rectangular mean observations}
For a disjoint partition $\mathcal P$ of a finite counting space, its
normalized means are $(Cf)_R=|R|^{-1/2}\sum_{x\in R}f(x)$. They satisfy
$CC^*=\id$ and $C^*C=P$, the orthogonal projection onto functions constant
on each cell. This concrete normalization matters for the penalty.

\begin{matrixTheorem}\label{mx:thm:critical}
Let $L=L^*$ on a finite real Hilbert space, $c>0$, and $C$ a coisometry
onto $\R^q$, $q\geq1$. On $\mathcal K=\ker C$ and
$\mathcal O=\im C^*$, write
\[
 L-c\id=\begin{pmatrix}A_0&X\\X^*&B_0\end{pmatrix}.
\]
A finite $k\geq0$ satisfying $L+kC^*C\geq c\id$ exists if and only if
\begin{equation}\label{mx:eq:compatibility}
 A_0\geq0,\qquad X^*(\ker A_0)=0.
\end{equation}
When these hold, its smallest value is
\[
 k_{\min}=\max\{0,\lambda_{\max}(X^*A_0^\dagger X-B_0)\}.
\]
Here $A_0^\dagger$ is the Moore--Penrose inverse on $\mathcal K$.
\end{matrixTheorem}
\begin{proof}
The restriction to $\mathcal K$ forces $A_0\geq0$. For $z\in\ker A_0$
the augmented form at $tz+y$ has linear term $2t\langle z,Xy\rangle$
and no quadratic term in $t$. Nonnegativity for every real $t$ forces
$X^*z=0$. This is equivalent to $\im X\subseteq\im A_0$. Under that
condition the full form at $z+y$ is exactly
\[
 \langle z+A_0^\dagger Xy,A_0(z+A_0^\dagger Xy)\rangle
 +\langle y,(B_0+k\id-X^*A_0^\dagger X)y\rangle.
\]
Finite spectral decomposition proves the criterion and the minimum.
\end{proof}
This is the classical generalized Schur criterion~\cite[Theorem 1]{mx:Higham},
applied to actual means; it is not a new general matrix theorem. Strict
positivity of $A_0$ suffices, whereas a weak bound at the target requires
the extra compatibility in~\eqref{mx:eq:compatibility}.

\begin{matrixExample}\label{mx:ex:rectangle}
On $(\mathbb Z/2\mathbb Z)^3$ take the same scalar coefficient in every
gradient direction,
\[
 a(x)=6+(-1)^{x_1}+4(-1)^{x_2},\qquad L=d^*ad.
\]
Its values are $1,3,9,11$, each twice, and its mean is $m=6$; hence
$1\leq a\leq2m$. Choose the four actual rectangles
$R_{yz}=\{0,1\}\times\{y\}\times\{z\}$ and
$(Cf)_{yz}=(f(0,y,z)+f(1,y,z))/\sqrt2$.
The best coercivity constant of $L$ on $\ker C$ is $16$, but no finite
quadratic penalty for these means attains the same constant on the full
space. Every $0<c<16$ is attainable. Thus the supremum $16$ is not attained.
\end{matrixExample}
\begin{proof}
For $S\subseteq\{1,2,3\}$ let
$\chi_S(x)=(-1)^{\sum_{i\in S}x_i}/\sqrt8$. These are an orthonormal real
Walsh basis. Retain all $24$ oriented forward rows, including opposite rows
at period two. The actual adjoint calculation gives
\begin{align*}
 L\chi_S={}&24|S|\chi_S
 +4\bigl(|S|-\boldsymbol1_{1\in S}\bigr)\chi_{S\triangle\{1\}}\\
 &+16\bigl(|S|-\boldsymbol1_{2\in S}\bigr)\chi_{S\triangle\{2\}}.
\end{align*}
Indeed $d_i\chi_S=-2\chi_S$ when $i\in S$, and the corresponding
adjoint term is $2(a+T_i a)\chi_S$; absent directions contribute zero.
In the kernel basis $(\chi_1,\chi_{12},\chi_{13},\chi_{123})$ the compression
has the two blocks
\[
 \begin{pmatrix}24&16\\16&48\end{pmatrix},\qquad
 \begin{pmatrix}48&32\\32&72\end{pmatrix}.
\]
After subtracting $16\id$, the first is positive semidefinite with kernel
vector $(-2,1)$, and the second is positive definite, with first pivot $32$
and determinant $768$. This proves the complete kernel bound, not just a
bound on selected test vectors.

Set $v=-2\chi_1+\chi_{12}$ and $u=\chi_2$. Direct evaluation gives
\[
 Cv=0,\quad\|v\|^2=5,\quad\|Cu\|^2=\|u\|^2=1,
 \quad(L-16\id)v=4u,\quad\langle u,(L-16\id)u\rangle=8.
\]
The augmented form on $u+tv$ is therefore $8+8t+k$, negative for
$t<-(8+k)/8$. No finite $k$ works. For $c<16$ the compressed block is
strictly positive and Theorem~\ref{mx:thm:critical} supplies a finite penalty;
for $c>16$, $v$ violates coercivity already in the kernel.
\end{proof}

\paragraph{A quantitative penalty obstruction.}
Write $\delta=16-c>0$. On the orthonormal pair $(v/\sqrt5,u)$ the
augmented form has matrix
\[
 \begin{pmatrix}\delta&4/\sqrt5\\4/\sqrt5&8+\delta+k\end{pmatrix}.
\]
Its determinant gives the necessary bound
\begin{equation}\label{mx:eq:blowup}
 k\geq\max\{0,16/(5\delta)-8-\delta\}.
\end{equation}
This lower bound need not equal the full optimum. It proves divergence
as the target approaches $16$ from below without numerical extrapolation.
If $F(0)=0$, $DF(0)=C$ and $0\leq\Phi(y)\leq k\|y\|^2$ near zero,
local coercivity would imply $L+kC^*C\geq c\id$ by substituting $sf$,
dividing by $s^2$ and taking $s\to0$. Higher order terms with this same
derivative do not repair the failure at $c=16$.

\paragraph{The exact penalty rather than only a lower bound.}
For this same scalar example, the full optimum at $c=16-\delta$,
$0<\delta<16$, is
\begin{equation}\label{mx:eq:exactpenalty}
 k_{\min}(c)=\max\left\{c,\frac{16(8+\delta)}{\delta(40+\delta)}-8-\delta\right\}.
\end{equation}
Here is a complete eight-dimensional check. Order the kernel basis as
$(\chi_1,\chi_{12})\oplus(\chi_{13},\chi_{123})$ and the observed basis as
$(\chi_\varnothing,\chi_2)\oplus(\chi_3,\chi_{23})$. At target $16$ the
kernel, observed and cross blocks of $L-16\id$ split into
\[
\begin{array}{lll}
 A_1=\begin{pmatrix}8&16\\16&32\end{pmatrix},&
 B_1=\operatorname{diag}(-16,8),&
 X_1=\begin{pmatrix}0&0\\0&4\end{pmatrix},\\[6pt]
 A_2=\begin{pmatrix}32&32\\32&56\end{pmatrix},&
 B_2=\begin{pmatrix}8&16\\16&32\end{pmatrix},&
 X_2=\operatorname{diag}(4,8).
\end{array}
\]
At $c=16-\delta$, add $\delta\id$ to each $A_i,B_i$, leaving $X_i$
unchanged. The first observed Schur block is diagonal, with entries
\[
 16-\delta,\qquad
 \frac{16(8+\delta)}{\delta(40+\delta)}-8-\delta.
\]
The second observed Schur block cannot exceed $(16-\delta)\id$:
$A_2-4\id$ has first pivot $28$ and determinant $432$, while $B_2\geq0$;
hence $X_2^*(A_2+\delta\id)^{-1}X_2\leq16\id$. The exact Schur minimum
in Theorem~\ref{mx:thm:critical} therefore proves both necessity and sufficiency of
\eqref{mx:eq:exactpenalty}, including all eight dimensions.
The branch transition is
\[
 \delta_0=\frac{\sqrt{3529}-59}{3}:
 \quad k_{\min}(c)=c\ \hbox{for }\delta\geq\delta_0,
\]
and the other branch is larger for $0<\delta<\delta_0$. Comparing
the entries reduces to $3\delta^2+118\delta-16\geq0$.
Partial fractions give the sharper critical expansion
\begin{equation}\label{mx:eq:residue}
 k_{\min}(c)=\frac{16}{5(16-c)}-\frac{192}{25}+O(16-c)
 \qquad(c\uparrow16).
\end{equation}
The previous lower bound~\eqref{mx:eq:blowup} alone is not the full optimum.

\paragraph{The general finite mechanism.}
For the blocks of Theorem~\ref{mx:thm:critical} at the least kernel eigenvalue $c_*>0$,
let $P$ project onto $\ker A_0$. At $c=c_*-\delta$,
\[
 k_{\min}(c)=\max\{0,\lambda_{\max}(
 X^*(A_0+\delta\id)^{-1}X-B_0-\delta\id)\}.
\]
Finite spectral splitting writes the inverse as $\delta^{-1}P$ plus a
bounded remainder converging to $A_0^\dagger$. The variational formula
for the largest eigenvalue consequently gives
\[
 \lim_{\delta\downarrow0}\delta k_{\min}(c_*-\delta)=\|PX\|^2.
\]
If $PX\ne0$, the penalty diverges as $\|PX\|^2/\delta+O(1)$ and the
critical target fails. If $PX=0$, it converges to the finite critical
Schur minimum. This uses classical finite Schur and resolvent principles,
not a claim of a new general matrix theorem.

\subsection{Transfer to actual flat coexact cochains}
The preceding obstruction also holds on the actual finite coexact domain,
with a proved form dictionary. Fix $N_c\geq2$ and an actual coordinate unit
$e\in\R^{N_c^2-1}$. At the trivial background, a physical one-cochain is
$A_j(x)\in\R^{N_c^2-1}$, $j=1,2,3$, with counting inner product. Use one
unordered plaquette copy for each $i<j$:
\[
 (d_1 A)_{ij}=d_iA_j-d_jA_i,\qquad
 \operatorname{div}A=\sum_j(A_j-T_j^{-1}A_j).
\]
Expansion of curl squares and cyclic summation by parts cancel the cross
terms against divergence, proving the all-cochain identity
\[
 \|d_1A\|^2+\|\operatorname{div}A\|^2=\sum_{i,j,x}\|d_iA_j(x)\|^2.
\]
These are the actual repository conventions, not a supplied comparison.
Take the actual componentwise repair coefficient $w(i,(x,j))=a(x)-1\geq0$.
The total Hodge-plus-repair form is consequently
\begin{equation}\label{mx:eq:physical-dictionary}
 \begin{aligned} Q_w(A)={}&\|d_1A\|^2+\|\operatorname{div}A\|^2+
 \sum_{i,j,x}w(x)\|d_iA_j(x)\|^2
 \\ ={}&\sum_{i,j,x}a(x)\|d_iA_j(x)\|^2. \end{aligned}
\end{equation}
It is a direct sum of $3(N_c^2-1)$ scalar copies. Componentwise normalized
means on the same four bars have output dimension $12(N_c^2-1)$, and
the whole ambient kernel has coercivity constant $16$ by the two blocks
in Example~\ref{mx:ex:rectangle}. Restriction to coexact fields preserves the
inequality; coisometry on the restricted domain is not assumed.

\begin{matrixProposition}\label{mx:prop:coexact}
At target $16$, no finite quadratic penalty for these actual componentwise
rectangle means supplies coercivity even on $\im d_1^*$. The same is true
for local nonlinear observations with derivative equal to those restricted
means and a locally quadratic penalty. On that domain every $0<c<16$ is
attainable, and the necessary bound~\eqref{mx:eq:blowup} still applies.
\end{matrixProposition}
\begin{proof}
Construct actual fields $U_3=ue$, $V_3=ve$ from Example~\ref{mx:ex:rectangle},
with all other components zero. They are isometric copies of $u,v$ and
independent of $x_3$, so their actual backward divergence vanishes.
More explicitly, use unordered two-cochains with only
\[
 B^U_{23}=-\tfrac12\chi_2 e,\qquad
 B^V_{13}=(\chi_1-\tfrac12\chi_{12})e
\]
nonzero. At side two $d_i^*=T_i^{-1}-\id=T_i-\id$. The adjoint of the
displayed curl puts $d_i^*B_{ij}$ in component $j$ and $-d_j^*B_{ij}$ in
component $i$. Since $d_3^*B^U=d_3^*B^V=0$ and each relevant character
has multiplier $-2$, these actual potentials satisfy $d_1^*B^U=U$ and
$d_1^*B^V=V$. Thus the entire line $U+tV$ is coexact.

Equation~\eqref{mx:eq:physical-dictionary} and the actual means give the same
augmented value $8+8t+k$. Scaling a violating coexact field towards zero
also excludes locally bounded quadratic penalties and the stated nonlinear
observations. Strict positivity on the ambient kernel supplies a finite
penalty for every $c<16$, which can then be restricted. The coexact field
$V$ excludes $c>16$. The same two-dimensional determinant proves the bound.
\end{proof}

\paragraph{Scope and source of the transfer.}
The pinned repository gives the actual form identities in these two sources:
\begin{quote}\small
\texttt{componentwiseFlatGradient\_eq\_hodge}\\
\path{YangMills/RG/ComponentwiseFlatHodge.lean}\\[4pt]
\texttt{actualTotalStiffness\_energy}\\
\path{YangMills/RG/ActualPacketSupportEnvelope.lean}
\end{quote}
Their proofs and the actual curl/divergence definitions were read; no new Lean compilation is
claimed. Unit 166 checks the newly constructed physical stencils rather
than relabelling a scalar diagnostic. For derivatives $nd_i$ at $n=2$,
energies and targets scale by four, so $16$ becomes $64$ with rescaled
potentials. This is finite rescaling, not a physical mass-gap claim.

The example has finite inverse moments at this one size. It refutes weak
kernel sufficiency for these prescribed rectangles at the same target;
it does not rule out another rectangular partition with the desired rank,
or establish an unbounded-size violation of the critical rank-order
conjecture. Only the trivial-background form has been identified here.


\paragraph{The sharp coexact leading coefficient.}
Let $k_{\mathrm{coex}}(c)$ be the minimum penalty for the actual means
restricted to $\im d_1^*$. The witness pair just proved and the ambient
direct-sum identity give, for $\delta=16-c\in(0,16)$,
\[
 \max\{0,16/(5\delta)-8-\delta\}
 \leq k_{\mathrm{coex}}(c)\leq k_{\min}(c).
\]
Multiplying by $\delta$ and using~\eqref{mx:eq:exactpenalty} proves
\[
 \lim_{c\uparrow16}(16-c)k_{\mathrm{coex}}(c)=\frac{16}{5}.
\]
No coisometry is assumed for the restricted mean map. This proves the
same leading coefficient on the actual coexact domain, not equality of
the full ambient and restricted optima or of their next constant terms.

\subsection{Reproduction and limitations}\label{mx:sec:evidence}
The source repository is
\href{https://github.com/lluiseriksson/THE-ERIKSSON-PROGRAMME/pull/186}{THE ERIKSSON PROGRAMME PR186},
stacked on PR185. All paths below are under
\texttt{reproducibility/matrix-moment-sweep-20261106/}.
Full source for this consolidation is pinned at
\begin{center}\small\texttt{66580c59ea9ef5f3c0374b339a165c9f89173db8}.\end{center}
Units 162 and 163 each used two distinct own CPU allocations, four actual
fresh sparse Git clones at one pinned source, and eight normal/optimized
Python processes. Outputs matched within each unit. Both runtimes were
deleted and disconnected after evidence export. Full stdout, stderr, exit
sentinels, Git logs, exact source bytes and manifests are in the corresponding
\texttt{evidence162-*.zip} and \texttt{evidence163-*.zip} archives.

\begin{center}\small
\begin{tabular}{llll}
\toprule
Unit & Executed source prefix & Fixture cases & Mutations per process\\
\midrule
162 & \texttt{f644c093} & four & seven types\\
163 & \texttt{378e879b} & four & seven types across four cases\\
\bottomrule
\end{tabular}
\end{center}
The 163 fixtures are $(D,n,h)=(3,2,1),(3,3,1),(3,3,2),(4,2,2)$.
They check actual Green and leverage identities, below/equal/above phase
inertia, equality witnesses, gap and mean-observation coercivity, and coupled
pair inertia. Unit 162 covers zero and near-zero stars, one large negative
channel and rationally rotated local eigendirections. Every acceptance check
raises explicitly on failure in both ordinary and \texttt{-O} execution.
There are 28 fixture-mutation pairs per 163 process, not 28 distinct types.
The two unit 163 recipe times were 53.83 and 29.98 seconds; requested-to-release
times were 179.71 and 102.74 seconds. Billed compute units were not exposed.

Units 165 and 166 each used two distinct own CPU allocations, four actual
fresh pinned Git clones and eight ordinary/optimized processes. Unit 165
checks the entire scalar torus matrix, the rectangular projection and three
distinct mutations per process. Unit 166 constructs the actual $72\times72$
physical matrix at $D=3,n=2,N_c=2$, with 72 unordered curl rows, 24 backward
divergence rows and 216 gradient rows. It checks the full Hodge-plus-repair
identity, both actual adjoint potentials, divergence zero and the rectangular
witnesses. Four distinct physical mutations are rejected per process: omit
divergence, double the unordered curl, reverse a potential sign and put the
field in the wrong component. Both 166 archives have 55 verified manifest
subjects. Their recipe times are 5.28 and 8.19 seconds; requested-to-release
times are 46.23 and 49.10 seconds. Both runtimes were disconnected after
export. Billed compute units remain unknown. The analytic all-$N_c$ coordinate
embedding and nonlinear limit are not extra executed cases.

Finite replay checks these finite matrices, not the universal analytical
statements. Arbitrary $q$, unequal weights, the all-size capped-family example
and the nonlinear derivative consequence are analytical arguments, not extra
executed fixtures. Earlier 159 diagnostics used source downloads rather than
fresh Git clones; their original limitation remains in their own records.
No Lean compilation certifies this appendix. A full independently prepared
audit of the imported Fourier adapter and the application form dictionary
remains necessary before broader claims.

Unit 168 derives the exact penalty and critical residue analytically.
Unit 169 then builds the actual scalar spatial matrix from all 24 forward
rows and the actual pair-average projection using rational arithmetic.
Generic full $8\times8$ positive-semidefinite elimination checks feasibility;
separate rational row reduction finds observed null vectors, which prove
that every strictly smaller penalty fails. Thus excess feasible penalties
are rejected rather than labelled optimal. Nine distinct rational
$\delta$ values cover both branches and bracket the transition without
rounding. Six mutation types per process reject the old lower bound used
as the optimum, an omitted constant-mode floor, doubled residue, excess
penalty, erased actual coupling and collapsed opposite rows.
Two own CPU allocations, four fresh pinned Git clones and eight ordinary
or optimized processes give identical payloads. Each complete archive has
53 verified manifest subjects and 16 real zero-exit sentinels. Recipe times
were 3.99 and 3.36 seconds; requested-to-release times were 36.32 and
41.92 seconds, including provisioning and export. Both runtimes were
disconnected after export; billed units are unknown. Complete records are
in \texttt{evidence169-*.zip} and \texttt{reproduction169-summary.json}.
These finite checks do not certify every real $\delta$, a new physical
$72\times72$ replay, the universal resolvent lemma or the coexact limiting
argument. Those remain the displayed analytic proofs. No new Lean build
certifies these extensions.

\subsection{Model assessments and review context}
This appendix was prepared and checked analytically in the continuing Codex
session that prepared its sources and diagnostics. The application does not
expose the exact backend model identity for that preparation. There is no
claim of independence, blinding, human peer review or editorial acceptance.
This exact new PDF has no external model assessment and no numerical score.
An assessment of an earlier PDF does not transfer to this artifact.

The retained objections are substantive: the Cwikel mechanism and Schur
principles are classical, so independent originality is unestablished; the
observation from spectral coordinates is generally nonlocal; sparse-count
and full-moment criteria must not be conflated; a diagonal graph-rank formula
does not cover rotated eigenchannels; negative depth is neither universally
dispensable nor universally necessary; the full off-diagonal interaction must
be retained; critical kernel compatibility cannot be omitted at the target; finite diagnostics do not prove all-size statements; and no
interacting physical form equality has been supplied. These limits are part
of the result, not editorial acceptance conditions represented as satisfied.

The original positive critical rectangular construction remains open.
The common-refinement route remains stopped after three unsuccessful units.
The recorded PR29 baseline stays 20/41 with TermSource 0. This appendix adds
supporting finite mathematical results and useful counterexamples without
changing those records. It is a draft for consolidation, not a deposit,
acceptance or new verified public release.


\begin{thebibliography}{99}
\small
\raggedright
\bibitem{eriksson} L. Eriksson, \emph{Exact-Kernel Fourier Families Obstruct
Every Penalty Factored Through a Block Map in a Flat Lattice Gauge Form},
ai.viXra 2608.0015v1 (2026), AIRR ARR-2026-7EPM3V3DP796BBRR v1.
\url{https://airr.science/papers/ARR-2026-7EPM3V3DP796BBRR/versions/v1/}.
Historical import, not peer reviewed or assessed by AIRR; checked 30 October 2026.
\bibitem{dodziuk} J. Dodziuk, \emph{Finite-difference approach to the Hodge
theory of harmonic forms}, American Journal of Mathematics \textbf{98}
(1976), 79--104. \url{https://doi.org/10.2307/2373615}.
Institutional source: \url{https://academicworks.cuny.edu/gc_pubs/180/}.
\bibitem{albert} A. Albert, \emph{Conditions for Positive and Nonnegative
Definiteness in Terms of Pseudoinverses}, SIAM Journal on Applied Mathematics
\textbf{17} (1969), 434--440. \url{https://doi.org/10.1137/0117041}.
\bibitem{carlson} D. Carlson, E. Haynsworth and T. Markham,
\emph{A Generalization of the Schur Complement by Means of the Moore--Penrose
Inverse}, SIAM Journal on Applied Mathematics \textbf{26} (1974), 169--175.
\url{https://doi.org/10.1137/0126013}.
\bibitem{dimock} J. Dimock, \emph{Covariant Axial Gauge},
arXiv:1410.2798v2 (2014). \url{https://arxiv.org/abs/1410.2798v2}.
\bibitem{mathlib-implicit} Y. Kudryashov and Mathlib contributors,
\emph{Implicit function theorem}, Mathlib source at commit
\texttt{07642720480157414db592fa85b626dafb71355b}.
\href{https://github.com/leanprover-community/mathlib4/blob/07642720480157414db592fa85b626dafb71355b/Mathlib/Analysis/Calculus/Implicit.lean}{Pinned primary formal source},
particularly \texttt{HasStrictFDerivAt.to\_implicitFunction} and its exact
level-set identity; inspected 30 October 2026. This attributes the classical parametrization; the separately pinned sector formalizations are identified in the evidence section.
\bibitem{spielman} D. A. Spielman, \emph{Graphs, Vectors, and Matrices},
Bulletin of the American Mathematical Society \textbf{54} (2017), 45--61.
Author-hosted primary exposition:
\url{https://cs-www.cs.yale.edu/homes/spielman/PAPERS/bull1557.pdf}.
\bibitem{levin-peres} D. A. Levin and Y. Peres, with contributions by
E. L. Wilmer, \emph{Markov Chains and Mixing Times}, second edition,
American Mathematical Society (2017), Chapter 13.
Author-hosted text: \url{https://pages.uoregon.edu/dlevin/MARKOV/markovmixing.pdf}.
\bibitem{zelazo} D. Zelazo and M. B{\"u}rger,
\emph{On the Definiteness of the Weighted Laplacian and its Connection to
Effective Resistance}, arXiv:1408.2187v1 (2014), Theorem III.3.
\url{https://arxiv.org/pdf/1408.2187}.
\bibitem{lin-wan-zhang} Y. Lin, S. Wan and H. Zhang,
\emph{Connection Laplacian on discrete tori with converging property},
arXiv:2403.06105v1 (2024), Theorem 1.2.
\url{https://arxiv.org/html/2403.06105v1}.
The sin-squared discrete-torus spectrum is attributed prior work, not a new
general spectral principle here.
\bibitem{ap:deberg} M.~de Berg, \emph{Linear Size Binary Space Partitions for
Uncluttered Scenes}, UU-CS-1998-12 (1998), Section2.1, Lemmas2.1 and2.4,
Theorem2.5 and Corollary2.9.
\url{https://ics-archive.science.uu.nl/research/techreps/repo/CS-1998/1998-12.pdf}.
\bibitem{ap:anderson} W.~N. Anderson, Jr. and G.~E. Trapp,
\emph{Shorted Operators. II}, SIAM J. Appl. Math. 28(1) (1975), 60--71.
\url{https://doi.org/10.1137/0128007}.
\bibitem{ap:frank} R.~L. Frank, \emph{Cwikel's theorem and the CLR inequality},
arXiv:1206.3325v1 (2012), Section2, Theorem2.1, Corollary2.2, Lemmas2.4--2.5.
\url{https://arxiv.org/abs/1206.3325}.
\bibitem{ap:ads} S.~Andres, J.-D. Deuschel and M.~Slowik,
\emph{Harnack inequalities on weighted graphs and some applications to the random
conductance model}, arXiv:1312.5473v6 (2019), Section2.2, equation(2.5) and Proposition2.1.
\url{https://arxiv.org/abs/1312.5473}.
\bibitem{mx:Frank} R. L. Frank, \emph{Cwikel's theorem and the CLR inequality},
arXiv:1206.3325v1 (2012), Theorem 3.1.
\url{https://arxiv.org/abs/1206.3325v1}.
\bibitem{mx:Zelazo} D. Zelazo and M. B\"urger, \emph{On the definiteness of the
weighted Laplacian and its connection to effective resistance},
IEEE CDC (2014), Theorem III.3.
\url{https://arxiv.org/abs/1408.2187}.
\bibitem{mx:Spielman} D. A. Spielman, \emph{Spectral Graph Theory Lecture 14},
15 October 2018, Claim 14.6.2.
\url{https://www.cs.yale.edu/homes/spielman/561/lect14-18.pdf}.
\bibitem{mx:Chen} W. Chen, J. Liu, Y. Chen, S. Z. Khong, D. Wang,
T. Ba\c{s}ar, L. Qiu and K. H. Johansson, \emph{Characterizing the positive
semidefiniteness of signed Laplacians via effective resistances},
IEEE CDC (2016), pp. 985--990.
\url{https://people.kth.se/~kallej/papers/network_cdc16chen.pdf}.
\bibitem{mx:Higham} N. J. Higham, \emph{What Is the Schur Complement of a Matrix?}, 1 June 2023, Theorem 1, generalized Schur complement. \url{https://nhigham.com/2023/06/01/what-is-the-schur-complement-of-a-matrix/}.

\end{thebibliography}
\end{document}
